\documentclass[11pt]{article}
\usepackage[T1]{fontenc}
\usepackage[utf8]{inputenc}
\usepackage{lmodern}
\usepackage[margin=1in]{geometry}
\usepackage{amsmath,amssymb,amsthm,mathtools,mathrsfs}
\usepackage{microtype}
\usepackage{graphicx,tikz}
\usepackage{booktabs,array,longtable}
\usepackage{enumitem}
\usepackage{xcolor}
\definecolor{linkblue}{RGB}{32,72,104}
\usepackage[colorlinks=true,linkcolor=linkblue,citecolor=linkblue,urlcolor=linkblue]{hyperref}
\hypersetup{pdftitle={Stable self-similar blowup for a supercritical defocusing Schrodinger equation on the torus},pdfauthor={OpenAI}}
\numberwithin{equation}{section}
\newtheorem{theorem}{Theorem}[section]
\newtheorem{proposition}[theorem]{Proposition}
\newtheorem{lemma}[theorem]{Lemma}
\newtheorem{corollary}[theorem]{Corollary}
\theoremstyle{definition}
\newtheorem{definition}[theorem]{Definition}
\theoremstyle{remark}
\newtheorem{remark}[theorem]{Remark}
\setlist{topsep=4pt,itemsep=2pt,parsep=0pt}
\allowdisplaybreaks[2]

\title{Stable self-similar blowup for a supercritical defocusing Schr\"odinger equation on the torus}
\author{OpenAI}
\date{September 24, 2026}

\begin{document}
\maketitle

\begin{abstract}
We prove stable self-similar finite-time blowup for a supercritical
defocusing nonlinear Schr\"odinger equation on the twelve-dimensional
torus. For a sufficiently large odd power, the blowup initial data contain
a nonempty open set in a high Sobolev space. As a consequence, Gaussian
Fourier initial data of arbitrarily high Sobolev regularity can have
positive probability of finite-time blowup.
\end{abstract}

\section{Introduction}
\label{sec:introduction}

Consider the defocusing nonlinear Schr\"odinger equation
\begin{equation}\label{eq:nls}
 i\partial_tu+\Delta u=|u|^{p-1}u
 \qquad\text{on }\mathbb T^d=(\mathbb R/2\pi\mathbb Z)^d,
\end{equation}
where $p\geq3$ is odd. Its scaling-critical Sobolev exponent is
$s_{\mathrm{cr}}=d/2-2/(p-1)$. In the energy-supercritical range
$s_{\mathrm{cr}}>1$, the positive conserved energy does not control
the regularity required for classical continuation. This leaves open
the possibility that singularities persist even for very regular data.
For odd $p$, the Sobolev algebra property and a Duhamel fixed point give
local well-posedness in $H^k(\mathbb T^d)$ for $k>d/2$.
Foundational Euclidean Cauchy theories include
\cite{GinibreVelo1979,CazenaveWeissler1990}. Boundedness of the critical
Sobolev norm implies global scattering in several whole-space
supercritical regimes \cite{KillipVisan2010}. Neither fact
supplies an a priori critical-norm bound from the conserved energy.
Bourgain surveyed global existence, blowup, and long-time dynamics
\cite{Bourgain2000}. Our result gives a robust deterministic
singularity in one such equation.

\begin{theorem}[An open set of blowup data]\label{thm:blowup}
There exist an odd integer $p\geq3$, a finite integer $k>8$, and a
nonempty open set $\mathcal U\subset H^k(\mathbb T^{12};\mathbb C)$
with the following property. For every $u_0\in\mathcal U$, the unique
local solution of \eqref{eq:nls} extends to
\[
 u\in C([0,T);H^k)\cap C^1([0,T);H^{k-2})
\]
for some finite $T>0$, but has no continuous $H^k$ extension through $T$.
More precisely, there are $x_*\in\mathbb T^{12}$ and $c_0>0$ such that
\[
 \lim_{t\uparrow T}(T-t)^{1/(p-1)}|u(t,x_*)|=c_0.
\]
The power can be chosen sufficiently large, and $k$ is then fixed after
that choice.
\end{theorem}

The prescribed-Gaussian question is motivated by the energy-supercritical
branch of Deng--Nahmod--Yue's Open problem~3
\cite[Section~9.2.3, p.~98]{DengNahmodYue2020}.
The eventual-$\alpha$, $H^m$ quantifiers below are the formulation
considered here, not a verbatim restatement of that problem.
We next state the probabilistic consequence with its quantifiers.
Let $g_n$, $n\in\mathbb Z^d$, be independent standard complex Gaussian
variables, and let $\langle n\rangle=(1+|n|^2)^{1/2}$.

\begin{corollary}[Failure of universal Gaussian global existence]
\label{cor:gaussian}
For the fixed power $p$ and index $k$ in Theorem~\ref{thm:blowup}, and
for every $\alpha>k+6$, the initial datum
\begin{equation}\label{eq:gaussian}
 u_0^\omega(x)=\sum_{n\in\mathbb Z^{12}}
       g_n(\omega)\langle n\rangle^{-\alpha}e^{in\cdot x}
\end{equation}
belongs to $H^k$ almost surely and has positive probability of finite-time
blowup. In particular, the following assertion is false: for every
energy-supercritical pair $(d,p)$ there is a finite $\alpha_0(d,p)$
such that, for every $\alpha>\alpha_0(d,p)$ and every
\[
 d/2+2<m<\alpha-d/2,
\]
the local $H^m$ solution with the Gaussian datum extends almost surely
to $C(\mathbb R;H^m(\mathbb T^d))$.
\end{corollary}

The probability in this corollary may depend on $\alpha$. We do not
claim probability-one blowup or a lower probability bound uniform in
$\alpha$. A single fixed pair $(d,p)$ is enough to refute the universal
assertion, and an explicit numerical value of the permissible power is
not needed.

\paragraph{Context.}
An earlier counterexample for a specially designed vector-valued,
forced defocusing Schr\"odinger system showed that conservation and
scaling alone cannot settle scalar global regularity
\cite[Theorem~1.2]{Tao2018}. For the unforced scalar energy-supercritical equation,
Merle, Rapha\"el, Rodnianski, and Szeftel
\cite[Theorem~1.1]{MRRS2022} obtain radial whole-space type-II blowup
on a finite-codimension family for the pairs
$(5,9),(6,5),(8,3),(9,3)$, subject to a stated profile nondegeneracy
condition. Cao-Labora, G\'omez-Serrano, Shi, and Staffilani
\cite[Theorem~1.1]{CLGSS2024} construct periodic nonradial blowup for
the cubic equation in dimension eight, again on a finite-codimension
family. Those constructions use hydrodynamic profiles with spatial
scale $(T-t)^{1/r}$, where $r>2$.
More recently, Buck, Carballeira, G\'omez-Serrano, and Shi
\cite[Theorem~1.1 and Section~4.3]{BuckEtAl2026} construct asymptotically self-similar
scalar defocusing blowup at the Schr\"odinger scale from smooth radial,
compactly supported data on $\mathbb R^d$ for
$(d,p)=(3,27),(4,9),(5,7)$. Their computer-assisted profile matching
and nonlinear stability yield a finite-codimension family; their
theorem does not assert a physical Sobolev-open blowup set.

Here the spatial scale is $(T-t)^{1/2}$, the amplitude is
$(T-t)^{-1/(p-1)}$, and there is a logarithmic phase rotation. This is
the meaning of self-similar, or type-I, blowup in this paper.
The open neighborhood in Theorem~\ref{thm:blowup} is essential for
Corollary~\ref{cor:gaussian}: full support of a Gaussian measure charges
every nonempty open set, whereas a finite-codimension existence theorem
alone does not give this conclusion. We prove the profile, spectral
stability, periodic localization, and openness needed here directly;
the cited blowup constructions provide historical and methodological
context, not inputs to those proofs.

Physical-data openness has precedents for the focusing equation.
Merle, Rapha\"el, and Szeftel \cite[Theorem~1.2]{MRS2010} construct
an $H^1$-open set of self-similar blowup data for slightly
mass-supercritical, energy-subcritical powers in dimensions one through
five. Under a mode-stability hypothesis, Li
\cite[Theorems~1.4--1.5]{Li2023} obtains asymptotic stability and
$H^1$-open blowup sets by matching symmetry parameters to a stable
graph. Here the spectral classification and parameter matching are
proved for the defocusing profile and the periodic evolution.

The probabilistic conclusion differs from global-flow results for
particular invariant ensembles in periodic and supercritical NLS
\cite{Bourgain1994,Sy2021,GueyeEtAl2025}. The Gaussian law in
Corollary~\ref{cor:gaussian} has full support in the stated $H^k$ space;
its positive-probability blowup event follows from the deterministic
physical-data openness, not from stochastic forcing or a random-data
construction of singular solutions.

\paragraph{Structure and methods.}
Write $a=1/(p-1)$. As $a\to0$, a radial profile develops a flat core
with positive limiting pressure and a free exterior. This limit turns
the spectral matching problem into a pair of explicitly defined
confluent hypergeometric solutions. A Laguerre recurrence gives analytic
tail bounds, including at a removable singular parameter. Exact integer
and rational certificates then count all low-angular roots and produce
a nonzero profile-matching degree; a complex-path multiplier treats
all remaining angular degrees. Section~\ref{sec:free} establishes these
free facts, and Appendix~\ref{cert:arithmetic} gives the finite
arithmetic recipes and their reproducible certificate.
Matched asymptotics for focusing NLS profiles appear in
\cite{BMR2021}; rigorous computer-assisted focusing profile
constructions appear in \cite{DahneFigueras2026,
DonningerSchorkhuber2026}, with exact fraction arithmetic in the latter.
Buck and coauthors \cite{BuckEtAl2026} use validated interval
arithmetic for the defocusing profile in different dimensions.
The present large-power limit and finite certificates are proved here
for the profile and full angular spectrum required by our theorem.

Section~\ref{sec:profiles} constructs genuine finite-power profiles and
the uniform bounds needed for the spectrum. Section~\ref{sec:spectrum}
shows that the only modes in the relevant half-plane are the phase,
twelve translations, and blowup time. Two points are central: a uniform
outgoing estimate excludes modes escaping to large temporal or angular
frequency, and compact convergence of a constrained pressure problem
preserves algebraic multiplicity in the bounded spectral region.
For slightly mass-supercritical focusing NLS, Li identifies the
low-energy symmetry spectrum \cite[Theorem~1.1]{Li2025Low}, excludes
other unstable eigenvalues, and deduces asymptotic stability with an
$H^1$-open blowup set in dimensions one through ten
\cite[Theorems~1.2--1.3]{Li2025High}. Our outgoing and multiplicity
arguments treat the defocusing large-power profile in dimension twelve.

The remaining analysis separates a local compact error from global
dissipation. Section~\ref{sec:linear} uses this estimate to transfer
decay from the whole-space evolution to expanding tori, retaining the
norm that might escape in a merely local limit. Finally,
Section~\ref{sec:nonlinear} constructs a graph of decaying perturbations
and varies the physical blowup parameters. A degree argument turns that
graph into an open set of physical data. The Gaussian support argument
then proves Corollary~\ref{cor:gaussian}.

\paragraph{Notation and choices of parameters.}
Throughout the construction $d=12$, $a>0$ is small, and
$D=y\cdot\nabla$. The integer $p=1+1/a$ is restricted to odd values
when polynomial smoothness is used. After choosing an admissible profile
and fixing $a$, we choose a sufficiently large finite integer $k$.
Later choices of a time-step length, a torus scale, and a data
neighborhood are made in that order. Implicit constants can depend on
previously fixed parameters. When a limit as $a\to0$ is used, the
required uniformities are stated explicitly. Complex-valued evolution
spaces are regarded as real spaces for the linearization, and are
complexified only for spectral analysis.

\section{Free exterior matching}\label{sec:free}

The profile construction and the spectral count both use a free radial
equation in dimension twelve. We study it first as an independent matching
problem. For a real parameter $b$, a spectral parameter $\lambda\in\mathbb C$,
an angular degree $\ell\in\{0,1,2,\ldots\}$, and a channel index $h=\pm1$,
the equation is
\begin{equation}\label{free:radial-equation}
 v_h''+\left(\frac{11}{r}+\frac{hi r}{2}\right)v_h'
 +\left(b+hi\lambda-\frac{\ell(\ell+10)}{r^2}\right)v_h=0.
\end{equation}
Here $\ell(\ell+10)$ is the eigenvalue of the negative spherical Laplacian
on degree-$\ell$ harmonics on $\mathbb S^{11}$. The substitution
$v_h=r^\ell H(q,6+\ell,x)$, with
$q=\lambda+\ell/2-hib$ and $x=-hi r^2/4$, gives the confluent
hypergeometric equation. Its slow solution is
\[
 r^\ell H(\lambda+\ell/2-hi b,6+\ell,-hi r^2/4),
\]
where Lemma~\ref{free:slow} defines $H$ by its integral and its
nonoscillatory asymptotic expansion. Primes on $v_h$ denote radial
derivatives; primes on $H$ denote derivatives in its last argument.
We need two kinds of boundary data
from this solution. At $\lambda=\ell=0$ in channel $h=1$, the quotient
$j=-H'/H$ vanishes exactly when the radial derivative vanishes; this
matches the exterior to a constant core. For general $\lambda,\ell$,
the two value-and-derivative columns in Proposition~\ref{free:main}
define the matching determinant whose zeros we count.

Section~\ref{sec:profiles} derives the similarity equation and constructs
the finite-power profiles. In those variables $b$ is the logarithmic phase
frequency and a spectral mode has time factor $e^{\lambda s}$, where
$s$ is logarithmic similarity time.
Lemma~\ref{spec:interface} will identify the determinant studied here
with the limiting constrained spectral problem. The calculations in this
section concern its free exterior, where the power parameter is $a=0$
and the pressure is zero.

All terminating decimals in this section denote exact rational numbers.  Set
\begin{equation}\label{free:disk}
 \begin{gathered}
 b_*=.33477606871236,\qquad Z_*=2.70506819293654,\qquad
 \delta=10^{-8},\\
 \mathcal D=\{(b,Z)\in\mathbb R^2:(b-b_*)^2+(Z-Z_*)^2\le\delta^2\}.
 \end{gathered}
\end{equation}
The radius of the limiting core is $R=2\sqrt Z$.
We prove the following statement, including its finite
arithmetic ingredients.  The definition of $H$ follows immediately below.
The degree on $\mathcal D$ means Brouwer degree on its interior.

\begin{proposition}[Free matching]\label{free:main}
For every $(b,Z)\in\mathcal D$, the following conclusions hold.
\begin{enumerate}
\item The quotient
\[
 j(b,Z)=-\frac{H'(-ib,6,-iZ)}{H(-ib,6,-iZ)}
\]
is well defined and satisfies $|j(b,Z)|<3\cdot10^{-9}$.  It is nonzero on
$\partial\mathcal D$, and, identifying $\mathbb C$ with $\mathbb R^2$ in the
order real part, imaginary part,
\[
 \deg(j,\mathcal D,0)=-1.
\]
\item For every real $z\ge Z$, $H(-ib,6,-iz)\ne0$ and
\begin{equation}\label{free:exterior-bounds}
 \operatorname{Re}j(b,z)<.043.
\end{equation}
For $z\ge3$ one also has $\operatorname{Im}j(b,z)>0$.
\item For $\ell\in\{0,1,2,\ldots\}$ and $h\in\{1,-1\}$ define
\begin{equation}\label{free:jet}
 q_h=\lambda+\ell/2-hi b,\qquad
 J_h(\lambda)=
 \begin{pmatrix}H(q_h,6+\ell,-hiZ)\\-H'(q_h,6+\ell,-hiZ)\end{pmatrix}.
\end{equation}
The columns $J_h$ are holomorphic and nonzero on a neighborhood of
$\operatorname{Re}\lambda\ge-1/32$.  The determinant
$\Delta_\ell(\lambda;b,Z)=\det(J_+(\lambda),J_-(\lambda))$ has no zero on
$\operatorname{Re}\lambda=-1/32$.  In the closed half-plane
$\operatorname{Re}\lambda\ge-1/32$ its numbers of zeros, counted with
holomorphic multiplicity, are
\[
 \begin{array}{c|ccc}
 \ell&0&1&\ell\ge2\\ \hline
 \#\{\Delta_\ell=0\}&2&1&0.
 \end{array}
\]
\end{enumerate}
The locations of the symmetry zeros are identified later at parameters
where $j(b,Z)=0$; the count in this proposition holds on the entire disk.
\end{proposition}

\subsection{The slow solution and the Laguerre cone}

The function $H(q,m,x)$ below is Tricomi's confluent hypergeometric
function $U(q,m,x)$; see \cite[Equations~13.4.4 and~13.3.22]{DLMF}.
We prove the regularized representation and the uniform symbol bounds
needed on the imaginary rays.

\begin{lemma}[Regularized slow solution]\label{free:slow}
For $\operatorname{Re}q>-1$, $\operatorname{Re}x\ge0$, $x\ne0$, define, with
principal powers,
\begin{equation}\label{free:H-definition}
 H(q,m,x)=x^{-q}\left\{1+\frac1{\Gamma(q)}
 \int_0^\infty e^{-u}u^{q-1}
 \left[(1+u/x)^{m-1-q}-1\right]\,du\right\}.
\end{equation}
Here $m$ may in particular be any positive integer used below.  This
definition is holomorphic in $q$, including $q=0$, and in $x$ on the
interior of the indicated region, with holomorphic continuation across
each nonzero point of its imaginary boundary.  Primes on $H$ denote
derivatives in its last argument.  It satisfies
\begin{equation}\label{free:H-equation}
 xH''+(m-x)H'-qH=0,\qquad
 H'(q,m,x)=-qH(q+1,m+1,x),\qquad H(0,m,x)=1.
\end{equation}
For every $N\ge1$ it has the expansion
\begin{equation}\label{free:H-expansion}
 H(q,m,x)=x^{-q}\left\{
 \sum_{j=0}^{N-1}\frac{(-1)^j(q)_j(1+q-m)_j}{j!x^j}
 +O(|x|^{-N})\right\}.
\end{equation}
Here $(q)_0=1$ and $(q)_j=\prod_{n=0}^{j-1}(q+n)$ for $j\ge1$.
The expansion is uniform for $\operatorname{Re}x\ge0$, $|x|\ge1$, and compact parameter
sets.  The same remainder estimate holds after any fixed number of
$x\partial_x$ derivatives of the bracket.  The symbol expansion specifies
the slow solution, without an oscillatory component on the imaginary
rays.
\end{lemma}

\begin{proof}
The bracket in the integrand of \eqref{free:H-definition} is $O(u)$ at
zero, locally uniformly in the parameters; at infinity it and its
parameter derivatives grow at most polynomially, with additional powers
of $\log u$.  Thus the regularized integral is locally holomorphic for
$\operatorname{Re}q>-1$.  The reciprocal gamma function vanishes at zero,
which proves the last assertion in \eqref{free:H-equation}.

For $\operatorname{Re}q>0$ the expression in braces can be combined into
the usual integral.  On $x>0$, substitution $u=xv$ gives
\[
 H(q,m,x)=\frac1{\Gamma(q)}\int_0^\infty
 e^{-xv}v^{q-1}(1+v)^{m-1-q}\,dv.
\]
Differentiation and integration by parts prove both the equation and the
derivative identity; holomorphic continuation proves them on the stated
domain.  The branches in the regularized integral are regular on a
neighborhood of each nonzero imaginary $x$, so the boundary values used
in the matching problem are genuine holomorphic values.

To obtain \eqref{free:H-expansion}, split the integral at $u=|x|/2$.
Taylor's formula applies uniformly to $(1+u/x)^{m-1-q}$ on the first
piece.  Its $j$th coefficient integrates to
$(-1)^j(q)_j(1+q-m)_j/j!$.  On the second piece the exponential weight
bounds the integral by any prescribed inverse power of $|x|$; the same
argument applies after symbol differentiation.  This proves the full
assertion about remainders.  For completeness, the other local solution
at infinity obtained by reduction of order has the form
$e^xx^{q-m}$ times an inverse-power expansion.  On either imaginary ray
its repeated symbol derivatives acquire arbitrarily high powers of
$|x|$.  It therefore cannot be added to \eqref{free:H-expansion} while
preserving all its symbol estimates.  This also explains the slow
normalization on the rays.
\end{proof}

Fix $h,\ell,b,Z$ and write
\[
 M=\ell+5,\quad s=hiZ,\quad q=\lambda+\ell/2-hi b,
 \quad G(t)=H(q,M+1,t-s)\quad(t\ge0).
\]
All parameter analyticity statements below concern $\lambda$; estimates
are locally uniform in the other displayed parameters.

\begin{lemma}[Laguerre recurrence]\label{free:laguerre}
Let $L_n(t)=e^t\partial_t^n(e^{-t}t^n)/n!$ be the Laguerre polynomials,
orthonormal for $e^{-t}\,dt$ on $[0,\infty)$.  Write
\[
 G=\sum_{n\ge0}g_nL_n,\qquad
 B_n=\sum_{k\ge n}g_k,\qquad C_n=\sum_{k\ge n}(k-n)g_k.
\]
These series and their parameter derivatives converge rapidly.  With
$t_n=q+n$ and $d_n=t_n-M$, one has
\begin{align}
 &(B_0,C_0)=(G(0),-G'(0)),\label{free:initial-tail}\\
 &t_ng_n=-MB_{n+1}-sC_n,\qquad
 d_nB_{n+1}=t_nB_n+sC_n,\qquad
 C_{n+1}=C_n-B_{n+1}.\label{free:recurrence}
\end{align}
The real quadratic form
\begin{equation}\label{free:cone-form}
 p_n=\frac M2|B_n|^2+\operatorname{Re}(s\overline{B_n}C_n)
\end{equation}
satisfies the exact identity
\begin{equation}\label{free:cone-increment}
 p_{n+1}-p_n=
 (\operatorname{Re}\lambda+n-5/2)|g_n|^2.
\end{equation}
\end{lemma}

\begin{proof}
Integration by parts gives orthonormality and
$[t\partial_t^2+(1-t)\partial_t]L_n=-nL_n$.
Polynomials are complete: a function orthogonal to all powers has a
Laplace transform whose derivatives vanish at an interior point of its
domain of holomorphy, hence its Laplace transform, and then the function,
vanishes.  The function $G$ and all its derivatives are smooth at zero
and of polynomial growth at infinity by Lemma~\ref{free:slow}.
Repeated integration by parts with this Laguerre operator shows that
$g_n$, and the Laguerre coefficients of each derivative of $G$, decrease
faster than every inverse power of $n$.  The boundary terms vanish
because the coefficient of the derivative at zero is $t$, and the
weight at infinity is exponential.  These arguments are uniform on
compact parameter sets, also after holomorphic differentiation.

The generating function gives, initially for $0\le x<1$,
\[
 \sum_{n\ge0}g_nx^n=
 \int_0^\infty e^{-u}G((1-x)u)\,du.
\]
Taking its value and first derivative as $x\uparrow1$ proves
\eqref{free:initial-tail}.  Since
$L_k'=-\sum_{n<k}L_n$, the coefficients of $G'$ are $-B_{n+1}$,
whereas those of $G''-G'$ are $C_n$.  Rewriting the differential equation
as
\[
 [t\partial_t^2+(1-t)\partial_t]G+MG'-s(G''-G')-qG=0
\]
proves the first recurrence.  The other two follow from
$g_n=B_n-B_{n+1}$ and the definition of $C_n$.

Finally insert $B_{n+1}=B_n-g_n$ and
$C_{n+1}=C_n-B_{n+1}$ into \eqref{free:cone-form}.  Since $s$ is purely
imaginary, the cross term $\operatorname{Re}(s|B_{n+1}|^2)$ vanishes.
Use $t_ng_n=-MB_{n+1}-sC_n$ to simplify the remaining terms.  The
coefficient of $|g_n|^2$ becomes
$\operatorname{Re}t_n-M/2=\operatorname{Re}\lambda+n-5/2$, which is
\eqref{free:cone-increment}.
\end{proof}

\begin{lemma}[Strict normalized cone]\label{free:cone}
Suppose $\operatorname{Re}\lambda\ge-1/32$ and $K\ge3$.  For $n\ge K$
the normalized coefficients $\widehat g_n=g_n/(q)_K$, and their sums
$\widehat B_K=B_K/(q)_K$, $\widehat C_K=C_K/(q)_K$, extend
holomorphically across $q=0$.  They satisfy
\begin{equation}\label{free:strict-cone}
 \frac M2|\widehat B_K|^2+
 \operatorname{Re}(s\overline{\widehat B_K}\widehat C_K)
 =-\sum_{n\ge K}(\operatorname{Re}\lambda+n-5/2)
 |\widehat g_n|^2<0.
\end{equation}
Consequently $\widehat C_K\ne0$, and
\begin{equation}\label{free:tail-disk}
 R_h=\widehat B_K/\widehat C_K\quad\hbox{is holomorphic},\qquad
 |R_h+s/M|<Z/M.
\end{equation}
\end{lemma}

\begin{proof}
The only zero of $(q)_K$ in the specified half-plane is $q=0$, and it is
simple.  It can occur only for $\ell=0$.  At this point $G=1$, so every
$g_n$ with $n\ge1$ vanishes.  Holomorphic division therefore removes
the apparent singularity for $n\ge K$.  Local Cauchy estimates preserve
the rapid decay established above, justifying division of the sums and
continuation of their identities.

The normalized tail cannot vanish identically in its sequence index.
Indeed repeated use of \eqref{free:H-equation} gives
\begin{equation}\label{free:normalized-derivative}
 \frac{\partial_t^K G}{(q)_K}
 =(-1)^KH(q+K,M+1+K,t-s).
\end{equation}
The right side is nonzero by its leading asymptotic coefficient, also at
$q=0$.  If all normalized coefficients with $n\ge K$ vanished, the
function $[G-\sum_{n<K}g_nL_n]/(q)_K$, with its removable value at zero,
would vanish.  Its $K$th derivative would be zero, contradicting
\eqref{free:normalized-derivative}.

Sum \eqref{free:cone-increment} from $K$ to infinity and divide by
$|(q)_K|^2$, taking the removable value when necessary.  Rapid decay
makes the limiting form zero.  Each weight in the resulting sum is
positive because $K\ge3$ and $\operatorname{Re}\lambda\ge-1/32$.
The preceding nontriviality gives strictness in
\eqref{free:strict-cone}.  If $\widehat C_K$ were zero the left side
would be nonnegative.  Dividing by its nonzero squared modulus and
completing the square proves \eqref{free:tail-disk}.
\end{proof}

\begin{lemma}[Analytic matching columns]\label{free:columns}
Define
\begin{equation}\label{free:backward}
 F_n=\begin{pmatrix}d_n-s&-s\\t_n&t_n\end{pmatrix},\qquad
 P_K=F_0F_1\cdots F_{K-1}.
\end{equation}
Then, including parameters where an individual $t_n$ or $d_n$ vanishes,
\begin{equation}\label{free:column-identity}
 P_K\begin{pmatrix}R_h\\1\end{pmatrix}
 =\frac1{\widehat C_K}
 \begin{pmatrix}H(q,M+1,-s)\\-H'(q,M+1,-s)\end{pmatrix}.
\end{equation}
The scalar on the right is holomorphic and nonzero.  The column is
nonzero.
\end{lemma}

\begin{proof}
The recurrences imply the division-free identity
$F_n(B_{n+1},C_{n+1})^t=t_n(B_n,C_n)^t$.  Multiplying gives
$P_K(B_K,C_K)^t=(q)_K(B_0,C_0)^t$.  Divide by $(q)_K$ and use
Lemma~\ref{free:cone}; analytic continuation gives the identity at the
removable point.  A simultaneous zero of $H$ and $H'$ at $x=-s\ne0$
would, by uniqueness for its second order equation, make $H$ identically
zero, contrary to \eqref{free:H-expansion}.
\end{proof}

\subsection{Exclusion of high angular degrees}

\begin{lemma}\label{free:high-angular}
For $(b,Z)\in\mathcal D$, $\ell\ge4$, and
$\operatorname{Re}\lambda\ge-1/32$, one has $\Delta_\ell(\lambda;b,Z)\ne0$.
\end{lemma}

\begin{proof}
Suppose the determinant vanished.  Since the two columns are nonzero,
multiply their functions by nonzero constants to arrange
$H_+(-iZ)=H_-(iZ)$ and $H_+'(-iZ)=H_-'(iZ)$.
Here $H_h(x)$ denotes the resulting multiple of
$H(\lambda+\ell/2-hi b,M+1,x)$ and $M=\ell+5\ge9$.
The substitution $U_h=e^{-x/2}x^{M/2}H_h$ gives
\begin{equation}\label{free:gauged}
 -(x\partial_xU_h)_x+\left(\frac x4+\frac{M^2}{4x}\right)U_h
 =(3-\lambda+hi b)U_h.
\end{equation}

Integrate this equation along the conjugate paths
\[
 \begin{gathered}
 x_h(\eta)=u-hi v,\quad x_h(0)=-hiZ,\quad
 g_h=x_h'=e^{-hi\vartheta},\\
 \vartheta=\min(2\eta,\vartheta_0),\quad
 \cos\vartheta_0=\frac{99}{101},\quad
 \sin\vartheta_0=\frac{20}{101}.
 \end{gathered}
\]
Write $\eta_0=\vartheta_0/2$ and
\[
 x_h/g_h^2=L+hiN,\qquad
 L=u\cos(2\vartheta)+v\sin(2\vartheta),\quad
 N=u\sin(2\vartheta)-v\cos(2\vartheta).
\]
In particular $L>0$ for $\eta>0$ and
\[
 L'=\cos\vartheta-2\vartheta'N,\qquad
 N'=\sin\vartheta+2\vartheta'L.
\]
The real part of $x$ tends to $+\infty$ linearly, so the exponential
factor in $U_h$ justifies every boundary limit and integration below.

For a bounded real function $\alpha$, smooth on each side of $\eta_0$,
put $B_c=\alpha L+\vartheta'N/2$.  Pair
\eqref{free:gauged} with $\overline U_h\,d\eta$ and take real parts.
The kinetic expression after integration by parts is
\[
 L|U_h'|^2+\operatorname{Re}
 \big(hi\vartheta'(L+hiN)\overline U_hU_h'\big).
\]
Completing the square writes this as
$L|U_h'+(\alpha-hi\vartheta'/2)U_h|^2
-B_c(|U_h|^2)'-L(\alpha^2+\vartheta'^2/4)|U_h|^2$.
The real lower endpoint terms before this completion sum to
$\sum_h(M/2)|U_h(0)|^2$.  To see the cancellation, expand
$x\overline U_h\partial_xU_h$ using $U_h=e^{-x/2}x^{M/2}H_h$:
the terms containing $x\overline H_hH_h'$ cancel between the two
matched jets, and the real contribution from the gauge is $M/2$.
Since $B_c(0)=-Z$, integration of the displayed derivative, separately
on the two path pieces, now yields
\begin{align}\label{free:multiplier-identity}
 0=\sum_{h=\pm1}\Bigg\{&
 (M/2-Z)|U_h(0)|^2+
 [B_c(\eta_0+)-B_c(\eta_0-)]|U_h(\eta_0)|^2\notag\\
 &+\int_0^\infty
 \left[L|U_h'+(\alpha-hi\vartheta'/2)U_h|^2
       +V_c|U_h|^2\right]\,d\eta\Bigg\},
\end{align}
where, with $\sigma=\operatorname{Re}\lambda$,
\begin{equation}\label{free:multiplier-potential}
 V_c=\sigma-3+\frac u4+\frac{M^2u}{4(u^2+v^2)}
 +\frac{\vartheta'\sin\vartheta}{2}
 +\frac34L\vartheta'^2+(\alpha L)'-L\alpha^2.
\end{equation}
There is no missing corner term: it is precisely the jump displayed in
\eqref{free:multiplier-identity}.

On the arc take $\alpha=.3$.  Here
$u=\sin\vartheta/2$, $v=Z+(1-\cos\vartheta)/2$, and
\[
 u\le.1,\qquad 2.70<v<2.72,\qquad
 N\le.04-2.70\cdot.92=-2.444,
 \qquad L'>10.7.
\]
It follows directly from \eqref{free:multiplier-potential}, dropping
its nonnegative $u$ terms and sine term, that
$V_c>-3-1/32+.3(10.7)+(3-.09)L>0$.
At the join $1.13<L(\eta_0)<1.19$.  On the straight piece, with
$t=\eta-\eta_0$, choose
\[
 \alpha L=-1.5+\frac{5t}{1+1.5t}.
\]
The jump is $-1.5-.3L(\eta_0)-N(\eta_0)>0$.
The geometry of the ray gives
\[
 u_0=.099+.98t\le u\le.100+.981t,\qquad
 v\le2.72+.20t,\qquad L\ge L_0=1.13+.98t.
\]
Set $d_0=(.100+.981t)^2+(2.72+.20t)^2$ and $q_0=1+1.5t$.
Since $M^2\ge81$, $V_cL_0d_0q_0^2$ is bounded below by
\begin{align}\label{free:high-polynomial}
 P(t)={}&(-3-1/32+u_0/4)L_0d_0q_0^2
 +\frac{81}{4}u_0L_0q_0^2+5L_0d_0
 -(-1.5+2.75t)^2d_0\notag\\
 >&2+t(33-72.2t+44.8t^2)
       +t^4(17.9-4.61t+.54t^2)>0.
\end{align}
The first strict comparison is coefficientwise exact arithmetic; the
last two quadratic discriminants are $-17519/25$ and
$-174119/10000$.  Appendix~\ref{cert:high-angular} provides the full
polynomial recipe and verifier.  Thus all terms of
\eqref{free:multiplier-identity} are nonnegative, the bulk potential is
strictly positive, and $M/2>Z$.  It follows that $U_h=0$, a
contradiction.
\end{proof}

\subsection{Boundary homotopy and the low angular counts}

We use $K=8$ for $0\le\ell\le3$.  For $0\le\rho\le1$ replace the tail
ratio in \eqref{free:column-identity} by $\rho R_h$.  The closure of
the disk in \eqref{free:tail-disk} is convex and contains zero, so every terminal
pair $(\rho R_h,1)$ satisfies $p_K\le0$.  The resulting columns are
holomorphic functions of $\lambda$, continuous in $\rho$.  We must
exclude determinant zeros on a common counting contour.

\begin{lemma}[Counting boundary]\label{free:boundary}
There is a rectangle with left edge $\operatorname{Re}\lambda=-1/32$
and other edges sufficiently far to the right and above and below,
uniformly for $(b,Z)\in\mathcal D$ and $0\le\rho\le1$, on whose
boundary the homotoped determinant is nonzero.  The same exclusions
hold on every larger such rectangle.
\end{lemma}

\begin{proof}
Choose the edges so that $t_n,d_n\ne0$ for $n<8$ on them.  If the
columns became dependent at such a point, their nonzero common initial
jet $(B_0,C_0)$ would have two forward images satisfying $p_{8,h}\le0$.
Indeed the backward matrices are invertible there and their inverses
are exactly the forward recurrences, up to nonzero scalar factors.
For the same initial jet in both channels,
\begin{equation}\label{free:summed-increment}
 p_{8,+}+p_{8,-}=M|B_0|^2+
 \sum_{h=\pm1}\sum_{n=0}^7
 (\sigma+n-5/2)|g_{n,h}|^2.
\end{equation}
For $\sigma>5/2$ all increment weights are positive.  If $B_0\ne0$
the first term is positive; if $B_0=0$ and $C_0\ne0$, then
$g_{0,h}=-s_hC_0/d_{0,h}\ne0$.  This excludes the right edge,
which can be taken beyond every possible zero of its divisors.

For the horizontal edges let $v=\operatorname{Im}\lambda$,
$C_0=vc$, and normalize $|B_0|^2+|c|^2=1$.
Uniformly for $\sigma$ in a bounded interval and $n<8$,
\[
 g_{0,h}=-(MB_0+s_hC_0)/d_{0,h}=-hZc+o(1),\qquad
 g_{n+1,h}=\frac{t_{n,h}g_{n,h}+s_hB_{n+1,h}}{d_{n+1,h}}
             =g_{n,h}+o(1).
\]
Consequently \eqref{free:summed-increment} equals
\[
 M|B_0|^2+2(8\sigma+8)Z^2|c|^2+o(1)>0
\]
for large $|v|$, since $8\sigma+8\ge31/4$.
This proves the two horizontal exclusions uniformly.

For the left edge set $s=iZ$ and form the forward matrix without its
divisors:
\begin{equation}\label{free:forward-matrix}
 T=T_{K-1}\cdots T_0,\qquad
 T_n=\begin{pmatrix}t_n&s\\-t_n&d_n-s\end{pmatrix},\qquad K=8.
\end{equation}
Thus the actual forward jet is $T(B_0,C_0)^t/\prod_{n<K}d_n$.
Define
\begin{equation}\label{free:boundary-form}
 \mathsf H(v)=T^*
 \begin{pmatrix}M&s\\-s&0\end{pmatrix}T,
 \qquad \mathsf S(v)=\mathsf H(v)+\overline{\mathsf H(-v)}.
\end{equation}
The second summand is the form for $h=-1$.  Each summand evaluates to
$2|\prod d_{n,h}|^2p_{K,h}$ on the common initial pair.  We shall prove
$\mathsf S(v)>0$ for every real $v$, excluding the two simultaneous
nonpositive values.

First, its determinant is a real even polynomial of degree at most
$4K-2=30$ for every real $(b,Z)$ in the parameter box.  Hermitian
symmetry gives real diagonal entries, and the reflection in
\eqref{free:boundary-form} makes the diagonal entries even and the
off-diagonal entry satisfy $C(-v)=\overline{C(v)}$.  To check the degree
bound structurally, hold the initial jet fixed and let $|v|\to\infty$.
The forward recurrences give $g_{n,h}=O((|B_0|+|C_0|)/|v|)$ for
$n<K$.  Also $|\prod d_{n,h}|^2=|v|^{2K}(1+O(1/|v|))$.
Using \eqref{free:cone-increment}, each form divided by $|v|^{2K}$
has its initial cone form as leading term and an $O(v^{-2})$
increment.  The opposite initial cross terms cancel in their sum.
Thus the first diagonal, off-diagonal, and second diagonal entries of
$\mathsf S/|v|^{2K}$ have respective orders $O(1)$, $O(v^{-1})$,
$O(v^{-2})$.  This proves the degree bound, independently of any
cancellation at a central parameter value.  Its first diagonal is
positive for sufficiently large $|v|$.

The exact coefficient calculation in Appendix~\ref{cert:boundary}
proves $\det\mathsf S(v)>0$ for all real $v$ and all parameters in
$\mathcal D$.  A continuous Hermitian matrix with strictly positive
determinant cannot change its signature.  Its positive first diagonal
at large $|v|$ therefore implies $\mathsf S(v)>0$ on the whole line.
This completes the left-edge exclusion and the proof.
\end{proof}

\begin{lemma}[Winding count]\label{free:winding}
For $0\le\ell\le3$ the zero counts in Proposition~\ref{free:main} are
$2,1,0,0$, respectively.
\end{lemma}

\begin{proof}
Lemma~\ref{free:boundary} and the argument principle preserve total
zero multiplicity throughout $\rho:1\to0$.  At $\rho=1$,
Lemma~\ref{free:columns} identifies this count with the true matching
determinant by nonzero holomorphic scalar factors.  Individual
homotoped columns may vanish at interior singular factors of $P_K$;
these are zeros of a holomorphic determinant and do not affect the
boundary argument.  No assertion of interior nonvanishing is needed.

We may also move $(b,Z)$ to $(B/H_0,S/H_0)$, where
$H_0=10^8$, $B=33477607$, $S=270506819$.  This point lies in
$\mathcal D$, and the preceding exclusions hold on the connecting
segment.  At zero tail ratio, put $\lambda=-1/32+X$ and let
$(\widetilde x,\widetilde y)^t=P_K(0,1)^t$ in channel $h=1$.
The real polynomial
\[
 P_\ell(X)=\operatorname{Im}_{\mathrm{coeff}}
 (\widetilde x(X)\overline{\widetilde y}(X))
\]
is a nonzero constant multiple of the determinant; the bar here
conjugates coefficients.  Induction in the backward recurrence shows
that its degree is $15=2K-1$, with unscaled leading coefficient $-KZ$:
$\widetilde y$ has leading term $X^K$ and $\widetilde x$ has leading
term $-KiZ X^{K-1}$.  Write
\begin{equation}\label{free:axis-polynomials}
 P_\ell(iv)=\mathcal R_\ell(v^2)+iv\mathcal I_\ell(v^2).
\end{equation}
Both real polynomials have degree at most seven.

The rational signs in Appendix~\ref{cert:winding} give the following
root counts; all roots in the table are simple:
\[
\begin{array}{c|cc|cc}
 &\multicolumn{2}{c|}{\text{positive roots}}&
   \multicolumn{2}{c}{\text{negative roots}}\\
 \ell&\mathcal R_\ell&\mathcal I_\ell&\mathcal R_\ell&\mathcal I_\ell\\\hline
 0&5&6&2&1\\
 1&6&6&1&1\\
 2,3&7&7&0&0
\end{array}
\]
Here is why signs alone suffice.  Successive positive separators have
alternating one-coordinate sign changes, following the quadrant cycle
$++,-+,--,+-$.  They therefore give the listed distinct positive
roots, interlaced between the two polynomials.  For $\ell=0,1$ the
additional values at $-1/4,-4$ supply the listed negative roots.
These distinct roots exhaust each degree bound of seven.  Hence there
are no further roots, none repeated, and no unrecorded axis crossings.

The curve $P_\ell(iv)$ for $v>0$ crosses each successive axis in the
positive rotational direction and finishes in the first quadrant.
Let $n=5,6,7$ be the positive root count of $\mathcal R_\ell$.
Starting with the real value at zero and using the final odd-degree
asymptotic, the continuous argument change on $[0,\infty)$ is
$(n+1/2)\pi$.  This also follows by counting the consecutive quadrants
from the initial signs $-+$, $++$, $--$ in the three rows.  Conjugation
doubles that change on the full oriented imaginary axis.  A single
factor $iv-z$ contributes $+\pi$ if $\operatorname{Re}z<0$ and
$-\pi$ if $\operatorname{Re}z>0$.  Thus
\[
 N_{\mathrm{left}}-N_{\mathrm{right}}=2n+1,
 \qquad N_{\mathrm{left}}+N_{\mathrm{right}}=15,
\]
and $N_{\mathrm{right}}=(15-(2n+1))/2=2,1,0$.
The counting boundary has no zeros, so these are also the counts in
the closed half-plane.  The count transfers back by the homotopies.
\end{proof}

\subsection{Profile degree and nonvanishing of the exterior}

\begin{lemma}[Profile matching degree]\label{free:profile-degree}
The first conclusion of Proposition~\ref{free:main} holds.  More
precisely there is an invertible real-linear map
$\mathcal L:\mathbb R^2\to\mathbb C$ with negative determinant and
minimum stretch greater than $.048$ such that
\begin{equation}\label{free:profile-linearization}
 |j(b,Z)-\mathcal L(b-b_*,Z-Z_*)|<2.5\cdot10^{-11}
 \qquad((b,Z)\in\mathcal D).
\end{equation}
In particular $|j|>4.55\cdot10^{-10}$ on $\partial\mathcal D$.
\end{lemma}

\begin{proof}
Use $h=1$, $\ell=\lambda=0$, $K=34$ in
\eqref{free:backward}, dividing its $n$th factor by $n+1$, and call
the resulting product $L(b,Z)$.  Set
\[
 U_0=L(b_*,Z_*),\quad X_0=\partial_bL(b_*,Z_*),\quad
 Y_0=\partial_ZL(b_*,Z_*),\quad m_0=(U_0)_{01}.
\]
All matrix indices here are zero-based.  The exact rational
enclosures of Appendix~\ref{cert:profile} give
\begin{gather}
 9<|m_0|<10,\qquad
 (U_0)_{00}/m_0\in(2.835,2.837)+i(1.426,1.428),\notag\\
 |(U_0)_{11}/m_0|<10^{-12},\qquad
 (X_0)_{11}/m_0\in(-.0811,-.0810)+i(.1361,.1363),\notag\\
 (Y_0)_{11}/m_0\in(-.0001,.0001)+i(.1237,.1240),\notag\\
 \max_{i,j}|(X_0)_{ij}/m_0|,
 \max_{i,j}|(Y_0)_{ij}/m_0|<20,\notag\\
 \max_{i,j}|L_{ij}-(U_0)_{ij}-(b-b_*)(X_0)_{ij}
                   -(Z-Z_*)(Y_0)_{ij}|<50000|m_0|\delta^2,\notag\\
 |\det L|<6.2\cdot10^{-12}|m_0|^2.\label{free:profile-enclosures}
\end{gather}
The same certificate, propagating these estimates by rational
inequalities, gives
\begin{equation}\label{free:denominator-gap}
 |L_{01}-iZL_{00}/5|-Z|L_{00}|/5>.6|m_0|,
 \qquad |L_{01}/m_0-1|<5\cdot10^{-7}.
\end{equation}
Since the tail ratio lies in the disk with center $-iZ/5$ and radius
$Z/5$, the first component $L_{00}R_++L_{01}$ cannot vanish.
Lemma~\ref{free:columns} proves the desired nonvanishing of $H$ before
we take its logarithmic derivative.

The second-to-first component ratio is $j(b,Z)$.  Subtracting its value
at zero tail gives the exact identity
\[
 j-\frac{L_{11}}{L_{01}}
 =-\frac{R_+\det L}{L_{01}(L_{00}R_++L_{01})}.
\]
Using $|R_+|\le2Z/5$ and \eqref{free:denominator-gap} bounds its
modulus by
\[
 \frac{(2Z/5)|\det L|}{.6|m_0|\,|L_{01}|}.
\]
Define
\[
 \mathcal L(\eta,\zeta)=
 \eta\frac{(X_0)_{11}}{m_0}+\zeta\frac{(Y_0)_{11}}{m_0}.
\]
Appendix~\ref{cert:profile} gives the exact propagation of the displayed
bounds to \eqref{free:profile-linearization}, the minimum stretch
bound, and $|\mathcal L(\eta,\zeta)|\le.283\delta$ on the disk.
Its determinant is negative already from the two enclosing
rectangles.  The boundary size of the linear map exceeds
$.048\delta=4.8\cdot10^{-10}$; subtracting the error proves the
stated boundary margin.  Straight homotopy to $\mathcal L$ avoids
zero on the boundary, so its degree is $-1$.  Finally
$.283\delta+2.5\cdot10^{-11}<3\cdot10^{-9}$ proves the disk bound.
\end{proof}

\begin{lemma}[Exterior inequalities]\label{free:exterior}
For $(b,Z)\in\mathcal D$, $H(-ib,6,-iz)$ is nonzero for all $z\ge Z$,
and the inequalities in \eqref{free:exterior-bounds} and
$\operatorname{Im}j(b,z)>0$ for $z\ge3$ hold.  In radial variables,
for any nonzero constant multiple $Q(r)$ of $H(-ib,6,-ir^2/4)$,
\begin{equation}\label{free:radial-quotients}
 \frac{A'}A=-\frac r2\operatorname{Im}j(b,r^2/4),\qquad
 \frac wr=\frac12+\operatorname{Re}j(b,r^2/4),
 \quad A=|Q|,\quad w=r/2+2\partial_r\arg Q.
\end{equation}
\end{lemma}

\begin{proof}
Use the forward matrix \eqref{free:forward-matrix} at
$\lambda=\ell=0$, $s=iz$, $K=5$, and let
\[
 \mathsf H=T^*\begin{pmatrix}5&iz\\-iz&0\end{pmatrix}T
   =\begin{pmatrix}A_0&\overline c\\c&d_0\end{pmatrix}.
\]
The exact coefficient checks in Appendix~\ref{cert:exterior} prove
\begin{align}\label{free:exterior-forms}
 &d_0>0,\qquad .043d_0+\operatorname{Re}c-1130z>0
 &&(z\ge2.704),\notag\\
 &-\operatorname{Im}c-1130z>0&&(z\ge3),\notag\\
 &|\det T|=\prod_{n=0}^4|t_nd_n|<1130.
\end{align}
Every $b$ in the disk is positive, so all forward divisors are
nonzero.  The exact jet has nonpositive terminal cone form, hence
$J^*\mathsf HJ\le0$.  If its first component were zero, positivity
of $d_0$ would force its second component to be zero, contradicting
Lemma~\ref{free:columns}.  Thus $j$ is well defined.  Moreover
\[
 \det\mathsf H=-z^2|\det T|^2,
 \qquad \left|j+\frac c{d_0}\right|
 \le\frac{z|\det T|}{d_0}.
\]
Taking the largest possible real part and smallest possible imaginary
part in this disk, and using \eqref{free:exterior-forms}, proves the
assertions.  All $Z$ in \eqref{free:disk} exceed $2.704$.
Finally $Q'/Q=(ir/2)j$ gives \eqref{free:radial-quotients}.
\end{proof}

Proposition~\ref{free:main} now follows from
Lemmas~\ref{free:high-angular}, \ref{free:winding},
\ref{free:profile-degree}, and \ref{free:exterior}.

\section{Smooth profiles with a pressurized core}\label{sec:profiles}

We work in dimension $12$.  The small parameter is $a>0$; later we fix
$a=(p-1)^{-1}$ with $p$ an odd integer.  In this section $a$ may take every
sufficiently small positive real value.  The profiles never vanish, so the
nonlinearity is smooth along them even without the restriction on $p$.
We use the parameter disk from Proposition~\ref{free:main}; all
terminating decimals below denote exact rational numbers.  Thus
\[
 \mathcal D=\{(b,Z):|(b,Z)-(b_*,Z_*)|\le10^{-8}\},\qquad
 b_*=.33477606871236,\quad Z_*=2.70506819293654.
\]
Put
\begin{equation}\label{profile:radii}
 \varepsilon=10^{-4},\qquad r_b=2\sqrt{Z_*}+\varepsilon,
 \qquad R(Z)=2\sqrt Z,\qquad r_0=r_b-6\varepsilon.
\end{equation}
The radius $r_b$ is fixed throughout the construction.

\subsection{Similarity variables and the amplitude equations}

For $\tau=T-t$, $s=\log(T/\tau)$ and
$y=(x-x_*)/\sqrt\tau$, write
\[
 u(t,x)=e^{i\theta}\tau^{-a}(\tau/T)^{ib}U(s,y).
\]
Since $\partial_t s=\tau^{-1}$ and
$\partial_t y=y/(2\tau)$, substitution into the defocusing equation gives
\begin{equation}\label{profile:similarity}
 i\partial_sU+\Delta U+i\left(\frac{y\cdot\nabla}{2}+a\right)U+bU
       =|U|^{1/a}U.
\end{equation}
We seek a nonvanishing radial stationary solution $Q=Ae^{i\phi}$, with
$A>0$ and $\phi(0)=0$.  Throughout this section a prime denotes a radial
derivative, except when explicitly attached to the last argument of $H$.
Set
\begin{equation}\label{profile:notation}
 P=A^{1/a},\quad \mu=A^2,\quad c_*=6-2a,\quad
 w=\frac r2+2\phi',\quad D=r\partial_r.
\end{equation}
The same letter $w$ denotes the vector field $w(r)y/r$, extended smoothly
at zero, and $W=w\cdot\nabla$.  The imaginary part of the stationary
equation is
\[
 2\frac{A'}A\phi'+\phi''+\frac{11}{r}\phi'
       +\frac r2\frac{A'}A+a=0.
\]
It follows that
\begin{equation}\label{profile:amplitude}
 \begin{split}
 w(r)&=c_*r^{-11}A(r)^{-2}\int_0^r A(t)^2t^{11}\,dt,\\
 -\Delta A+A^{1+1/a}&=VA,\qquad
 V=b+\frac{r^2}{16}-\frac{w^2}{4},\\
 w'&=c_*-\frac{11w}{r}-2w\frac{A'}A.
 \end{split}
\end{equation}
Here the integration constant in the first line is zero by regularity at
the origin.  The real part gives the second line because
$\phi'^2+r\phi'/2=w^2/4-r^2/16$.  In particular,
\begin{equation}\label{profile:divergence}
 \operatorname{div}(\mu w)=c_*\mu.
\end{equation}
Conversely, positive amplitudes satisfying the first two equations in
\eqref{profile:amplitude}, with $A'(0)=0$, give a stationary solution by
integrating $\phi'=(w-r/2)/2$.

Here is the complete profile statement used subsequently.  A constant
with a subscript $j$ may depend on that fixed derivative order.  Uniform
statements as $a\downarrow0$ refer to all profiles selected by the
construction, without requiring a continuous selection of their matching
parameters.

\begin{proposition}[Profiles and uniform estimates]\label{profile:existence}
There is $a_0>0$ such that, for every $0<a<a_0$, there are
$(b_a,Z_a)\in\mathcal D$ and a smooth nonvanishing radial stationary
solution $Q_a=A_ae^{i\phi_a}$ of \eqref{profile:similarity}, normalized by
$Q_a(0)>0$.  They have the following properties.
\begin{enumerate}
\item The identities \eqref{profile:amplitude}--\eqref{profile:divergence}
hold, $0<P_a\le1$, and there are constants $c,C>0$, independent of small
$a$, such that
\begin{equation}\label{profile:coercive-bounds}
 C^{-1}\langle r\rangle^{-4a}\le\mu_a(r)
       \le C\langle r\rangle^{-4a},\qquad
 \nabla w_a\ge c I.
\end{equation}
\item If $a_n\downarrow0$ and $(b_{a_n},Z_{a_n})\to(b,Z)$, then
$j(b,Z)=0$, where $j$ is the free matching function of
Proposition~\ref{free:main}.  With $R=2\sqrt Z$, the profiles converge in
$C^1$ on every compact ball to
\begin{equation}\label{profile:limit}
 Q_0(r)=
 \begin{cases}
 1,&0\le r\le R,\\[2pt]
 \displaystyle\frac{H(-ib,6,-ir^2/4)}{H(-ib,6,-iZ)},&r\ge R.
 \end{cases}
\end{equation}
The exterior amplitude is strictly decreasing and has a positive limit
at infinity.  The convergence is $C^\infty$ on compact subsets of
$\{r>R\}$.  On every fixed ball,
\begin{equation}\label{profile:pressure-limit}
 P_{a_n}\ \stackrel{*}{\rightharpoonup}\ b\mathbf1_{\{r<R\}}
       \quad\hbox{in }L^\infty,
\end{equation}
and $\mu_{a_n},w_{a_n}$ converge in $C^1$.  In particular
$P_{a_n}/a_n\to0$ locally uniformly on $\{r>R\}$.
\item On the fixed exterior $r\ge r_b$,
\begin{equation}\label{profile:tail}
 Q_a(r)=r^{\nu_a}f_a(r),\qquad \nu_a=-2a+2ib_a,
 \qquad 0<\kappa\le|f_a(r)|\le\rho<1,
\end{equation}
with $\kappa,\rho$ independent of $a$.  There are coefficients
$m_{a,j}$, with $|m_{a,0}|\ge\kappa$, such that for every pair of fixed
nonnegative integers $N,J$,
\begin{equation}\label{profile:symbol-expansion}
 \left|D^j\left(f_a(r)-\sum_{q=0}^{N}m_{a,q}r^{-2q}\right)\right|
       \le C_{N,J}r^{-2N-2},\qquad 0\le j\le J,\quad r\ge r_b.
\end{equation}
The constants are uniform as $a\downarrow0$.  In particular,
\begin{equation}\label{profile:exterior-pressure}
 |D^j(P_a/a)|\le\eta_j(a)r^{-2},\qquad
 |D^je^{2i\phi_a}|\le C_j,\qquad \eta_j(a)\longrightarrow0
\end{equation}
for each fixed $j$.  On the inner ball,
\begin{equation}\label{profile:pressure-transport}
 W_a(P_a/a)\le Ca(P_a/a),\qquad r\le r_b,
\end{equation}
whereas $|W_a(P_a/a)|\le\eta(a)r^{-2}$ for $r\ge r_b$, with
$\eta(a)\to0$.  Consequently the positive part of $W_a(P_a/a)$ is
bounded uniformly on the whole space.
\item For each fixed $a$, every Cartesian derivative obeys
\begin{equation}\label{profile:fixed-power-symbols}
 |\partial^\beta Q_a(y)|\le C_{a,\beta}\langle y\rangle^{-2a-|\beta|},
 \qquad
 \left|\partial^\beta\big((P_a/a)e^{2i\phi_a}\big)\right|
       +|\partial^\beta(P_a/a)|
       \le C_{a,\beta}\langle y\rangle^{-2-|\beta|}.
\end{equation}
The time-symmetry combination has the improved exterior order
\begin{equation}\label{profile:time-symbol}
 \left|D^j\left((a-ib_a+D/2)Q_a\right)\right|
       \le C_j r^{-2a-2},\qquad r\ge r_b.
\end{equation}
\end{enumerate}
\end{proposition}

We prove the proposition in stages.  The inner problem is solved on the
whole parameter disk, as is the slow exterior problem.  Only their final
matching uses a choice of $(b_a,Z_a)$.

\subsection{The inner boundary value problem}

Let $F(r;b,R')$ solve
\begin{equation}\label{profile:free-initial}
 F''+\left(\frac{11}{r}+\frac{ir}{2}\right)F'+bF=0,
       \qquad F(R')=1,\quad F'(R')=0,
\end{equation}
and set
\[
 F_{\rm in}(r;b,Z)=F(r;b,R(Z)),\qquad
 h_b(b,Z)=|F_{\rm in}(r_b;b,Z)|.
\]
The subscript on $h_b$ indicates the fixed boundary radius, not
dependence on $b$ alone.

\begin{lemma}[The boundary modulus]\label{profile:boundary-modulus}
Uniformly on $\mathcal D$,
\begin{equation}\label{profile:boundary-gap}
 1-.2\varepsilon^2<h_b<1-.15\varepsilon^2.
\end{equation}
For fixed $b$ in the projection of $\mathcal D$, the function
$R'\mapsto|F(r_b;b,R')|$ is strictly increasing on $[r_0,r_b)$.
\end{lemma}

\begin{proof}
On these intervals the coefficient $11/r+ir/2$ has modulus less than
$6$, and $.334<b<.335$.  Put $u=r-R'$, so $0\le u\le6\varepsilon$.
Integrating the two first-order equations in \eqref{profile:free-initial}
gives, by a bootstrap starting at $u=0$,
\[
 |F'|\le .35u,\qquad |F-1|\le .175u^2,
       \qquad |F''+b|<.003.
\]
For completeness, assuming the first bound up to a first failure gives
$|F-1|\le.175u^2$ and
$|F''+b|\le6(.35)u+.335(.175)u^2<.0013$; integration then gives
$|F'|<.337u<.35u$, excluding that failure.  The displayed weaker
$.003$ bound follows as well.  Thus
\[
 F(r_b)=1-\frac b2(r_b-R')^2+E,\qquad
       |E|\le.0015(r_b-R')^2.
\]
Since $|R(Z)-2\sqrt{Z_*}|<10^{-8}$, taking $R'=R(Z)$ proves
\eqref{profile:boundary-gap}, including the negligible quadratic error
when passing from the complex value to its modulus.

Differentiate the initial-value problem with respect to $R'$.  The
derivative $G=\partial_{R'}F$ satisfies the same linear equation and has
data $G(R')=0$, $G'(R')=b$.  Integration, or the same bootstrap, gives
$|G(r_b)/(r_b-R')-b|<.003$.  As $F(r_b)=1+O((r_b-R')^2)$,
\[
 \partial_{R'}|F(r_b)|
   =\frac{\operatorname{Re}(\overline{F(r_b)}G(r_b))}{|F(r_b)|}>0.
\]
\end{proof}

\begin{lemma}[The inner amplitude]\label{profile:inner}
For all sufficiently small $a>0$ and all $(b,Z)\in\mathcal D$, the
system \eqref{profile:amplitude} has a positive radial solution on
$[0,r_b]$ with $A'(0)=0$ and $A(r_b)=h_b$.  Within the bounds
\begin{equation}\label{profile:inner-bounds}
 h_b\le A\le .5^a,
\end{equation}
it is unique.  It depends continuously on $(b,Z)$, and obeys
\begin{equation}\label{profile:bulk-barrier}
 A(r)\ge .2^a-.02(r-r_0)_+^2.
\end{equation}
The amplitudes are bounded uniformly in $C^{1,1}[0,r_b]$.
\end{lemma}

\begin{proof}
First take $a$ small enough that $h_b^{1/a}\le.3$ uniformly on the disk;
this also ensures $h_b<.5^a$.  For a continuous radial input $A$ satisfying
\eqref{profile:inner-bounds}, define $w,V$ by the first and second lines
of \eqref{profile:amplitude}.  The useful nonsingular formula is
\begin{equation}\label{profile:w-ratio}
 \frac{w(r)}r=c_* A(r)^{-2}\int_0^1A(rs)^2s^{11}\,ds.
\end{equation}
Since $h_b>.999$, $r_b^2<11$, and $a$ is small, it implies
\begin{equation}\label{profile:V-bounds}
 .3\le V\le .5,\qquad
 \|V[A_1]-V[A_2]\|_\infty\le10\|A_1-A_2\|_\infty.
\end{equation}
Indeed $c_*h_b^2/12\le w/r\le c_*/(12h_b^2)$; inserting these bounds
into $V=b+r^2(1/16-(w/r)^2/4)$ gives the first assertion for small $a$.
The Lipschitz constant of $A\mapsto w/r$ is at most
$c_*(2h_b^{-3}+2h_b^{-2})/12<2.01$; multiplying by
$r_b^2(\sup w/r)/2$ gives a constant less than $10$ for $V$.

For this prescribed $V$, solve the scalar Dirichlet problem
\begin{equation}\label{profile:scalar-dirichlet}
 -\Delta H+H^{1+1/a}=VH,\qquad H(r_b)=h_b,\quad H'(0)=0.
\end{equation}
The constants $h_b$ and $.5^a$ are respectively a lower and an upper
solution.  One can obtain $H$ by monotone iteration: choose $L>0$ so that
$t\mapsto(L+V(r))t-t^{1+1/a}$ is increasing on $[h_b,.5^a]$, and at
each step invert $-\Delta+L$ with the given boundary value.  This inverse
is positive by the minimum principle.  Its existence, if needed directly
in the radial class, follows by minimizing the strictly convex functional
$\int(|\nabla u|^2+L|u|^2-2fu)/2$ for zero boundary data; the radial
equation then reads $(r^{11}u')'=r^{11}(Lu-f)$ and gives the classical
solution.  Iteration from $h_b$ is increasing and stays below $.5^a$.
For fixed $a$ its right sides are uniformly bounded; radial integration
gives compactness in $C^1$, and its limit solves
\eqref{profile:scalar-dirichlet}.

We record a comparison argument that also applies below.  The radial
Dirichlet inverse of $-\Delta$ is
\[
 (Kf)(r)=\int_r^{r_b}t^{-11}\int_0^t s^{11}f(s)\,ds\,dt,
       \qquad \|K\|_{\infty\to\infty}=r_b^2/24<11/24.
\]
Thus the positive geometric series defines
$T=(-\Delta-.5)^{-1}=\sum_{n\ge0}(.5K)^nK$ with zero boundary values.
Also $\psi=1+K(1+.5\psi)$ has a positive solution and satisfies
$(-\Delta-.5)\psi=1$, $\psi(r_b)=1$.  For every $q\ge0$ and $V\le.5$,
\[
 (-\Delta-V+q)\psi\ge1.
\]
If a function with nonnegative boundary values violated the comparison
principle for $-\Delta-V+q$, its ratio to $\psi$ would have a negative
interior minimum, where the product rule gives a strictly negative value
for that operator.  This is a contradiction.  Taking $q$ to be the
divided difference of $t^{1+1/a}$ proves uniqueness in
\eqref{profile:scalar-dirichlet} and comparison with lower solutions.

The function
$B(r)=.2^a-.02(r-r_0)_+^2$ is such a lower solution.  For small $a$ it
is positive, $B(r_b)<h_b$ by \eqref{profile:boundary-gap},
$B^{1/a}\le.2$, and
\[
 -\Delta B\le .04+.44\frac{6\varepsilon}{r_0}<.05
       \le(.3-.2)B.
\]
At $r_0$ the first derivative is continuous, so no interface measure
occurs in this inequality.  Comparison proves
\eqref{profile:bulk-barrier} for every output $H$.

Let $H_1,H_2$ be outputs for two inputs.  Their difference solves
\[
 (-\Delta-V_1+q)(H_1-H_2)=(V_1-V_2)H_2,
 \quad
 q=\int_0^1(1+1/a)(tH_1+(1-t)H_2)^{1/a}\,dt.
\]
The barrier implies $q\ge.2/a$ on $r\le r_0$, whereas $q\ge0$
everywhere.  A positive supersolution for unit forcing is
\begin{equation}\label{profile:small-inverse}
 S=10a+2T\mathbf1_{\{r>r_0\}}.
\end{equation}
On the inner ball, the constant term contributes at least $2-5a$ and
the second term contributes a nonnegative quantity.  On the shell, the
second term contributes at least $2$, while the constant term is bounded
below by $-5a$.  Thus the comparison principle applies, including at
the shell interface.  To bound $S$, radial integration gives
\[
 \|K\mathbf1_{\{r>r_0\}}\|_\infty
 \le\int_{r_0}^{r_b}(t-r_0)\,dt=(6\varepsilon)^2/2,
 \qquad
 \|T\mathbf1_{\{r>r_0\}}\|_\infty\le(6\varepsilon)^2.
\]
Here the second inequality uses
$(1-.5\|K\|)^{-1}<2$.  Consequently the input-output map satisfies
\[
 \|H_1-H_2\|_\infty
 \le10\big(10a+2(6\varepsilon)^2\big)\|A_1-A_2\|_\infty.
\]
Our already fixed shell width makes the second contribution less than
$1/4$; taking $a$ small makes the whole factor less than $1/2$.
The closed interval \eqref{profile:inner-bounds} in the radial continuous
functions is complete, so successive iteration gives its unique fixed
point.  This proves the coupled inner problem.

Finally $0<P\le.5$ and \eqref{profile:V-bounds} give
\[
 A'(r)=r^{-11}\int_0^r t^{11}(P(t)-V(t))A(t)\,dt,
       \qquad |A'|\le Cr,\quad |A''|\le C.
\]
These are the asserted uniform bounds.  For fixed $a$, convergent
parameter sequences have $C^1$ convergent subsequences; the equation and
uniqueness identify their limits.  This proves continuous dependence,
including when the boundary value $h_b$ varies.
\end{proof}

\begin{lemma}[Identification of the core]\label{profile:core-limit}
For the inner solutions of Lemma~\ref{profile:inner}, suppose
$a_n\downarrow0$ and $(b_n,Z_n)\to(b,Z)\in\mathcal D$.  Then their
complex reconstructions converge in $C^1[0,r_b]$ to the function equal to
$1$ on $[0,R(Z)]$ and to $F_{\rm in}$ on $[R(Z),r_b]$.
Their pressures converge weak star to $b\mathbf1_{\{r<R(Z)\}}$.
These convergence assertions are uniform over the compact parameter disk
in the sequential sense just stated.
\end{lemma}

\begin{proof}
The preceding bounds give $C^1$ compactness of the amplitudes.  Every
limit $A_0$ is at most $1$ and, by \eqref{profile:bulk-barrier}, equals
$1$ on $[0,r_0]$.  Formula \eqref{profile:w-ratio} and the last equation
in \eqref{profile:amplitude} give $C^1$ convergence of $w$, including at
zero.  On every compact interval where $A_0<1$, the pressure tends
uniformly to zero.  There the limit satisfies $\Delta A_0=-V_0A_0$,
with $V_0\ge.3$.

Each component of $\{A_0<1\}$ has a left endpoint $R'$ where
$A_0(R')=1$, $A_0'(R')=0$.  On this component
$(r^{11}A_0')'=-r^{11}V_0A_0<0$, so $A_0'<0$ and the component cannot
have a right endpoint below $r_b$.  As $A_0(r_b)=h_b<1$, there is exactly
one such component, $(R',r_b]$, and $R'\in[r_0,r_b)$.  On the plateau
\eqref{profile:w-ratio} gives $w_0=r/2$ and hence $\phi_0=0$.
Reconstruction gives the free complex equation with data $1,0$ at $R'$.
Its boundary modulus is $h_b$, so
Lemma~\ref{profile:boundary-modulus} forces $R'=R(Z)$.

The pressures are bounded by $.5$.  Any weak-star limit $P_0$ satisfies
in distributions
\[
 P_0A_0=V_0A_0+\Delta A_0.
\]
There is no surface measure at $R$ because $A_0$ and $A_0'$ are
continuous.  On the plateau the right side is $b$, and on its exterior
it is zero.  Since $A_0$ is positive, this identifies $P_0$ uniquely.
The preceding identification applies to every subsequence; compactness
therefore proves all convergence assertions.  Reconstruction by
$\phi'=(w-r/2)/2$ gives the asserted $C^1$ convergence of the complex
profiles as well.
\end{proof}

\begin{figure}[tbp]
  \centering
  \begingroup
  \definecolor{coreblue}{RGB}{48,91,120}
  \definecolor{coreorange}{RGB}{163,92,43}
  \resizebox{0.95\linewidth}{!}{%
  \begin{tikzpicture}[
    x=1cm,y=1cm,font=\small,
    axis/.style={->,thin,black!75},
    guide/.style={densely dashed,thin,black!45},
    amplitude/.style={very thick,coreblue},
    pressure/.style={very thick,coreorange},
    every node/.style={inner sep=2pt}]
    \begin{scope}
      \fill[coreorange!7] (0,0) rectangle (2.2,2.6);
      \draw[axis] (0,0) -- (5.9,0) node[right] {$r$};
      \draw[axis] (0,0) -- (0,2.8) node[above] {$A_0(r)$};
      \draw[guide] (2.2,0) -- (2.2,2.55);
      \draw[amplitude] (0,2.1) -- (2.2,2.1)
        .. controls (2.85,2.1) and (3.15,1.3) .. (3.85,1.02)
        .. controls (4.4,.8) and (5.0,.7) .. (5.65,.66);
      \draw (0,2.1) -- (-.08,2.1) node[left] {$1$};
      \draw (2.2,0) -- (2.2,-.09) node[below] {$R$};
      \node[below left] at (0,0) {$0$};
      \node at (1.1,1.15) {constant core};
      \node[align=center] at (4.3,2.25) {free exterior\\strictly decreasing};
      \node at (2.85,-.65) {(a) Limiting amplitude};
    \end{scope}
    \begin{scope}[xshift=7.4cm]
      \fill[coreorange!7] (0,0) rectangle (2.2,2.6);
      \draw[axis] (0,0) -- (5.9,0) node[right] {$r$};
      \draw[axis] (0,0) -- (0,2.8) node[above] {$P_0(r)$};
      \draw[guide] (2.2,0) -- (2.2,2.55);
      \draw[pressure] (0,2.1) -- (2.2,2.1);
      \draw[pressure] (2.2,0) -- (5.65,0);
      \draw[pressure,fill=white] (2.2,2.1) circle (.045);
      \draw[pressure,fill=white] (2.2,0) circle (.045);
      \draw (0,2.1) -- (-.08,2.1) node[left] {$b$};
      \draw (2.2,0) -- (2.2,-.09) node[below] {$R$};
      \node[below left] at (0,0) {$0$};
      \node at (1.1,1.15) {pressure $b$};
      \node at (4.15,1.15) {zero pressure};
      \node at (2.85,-.65) {(b) Limiting pressure};
    \end{scope}
  \end{tikzpicture}%
  }
  \endgroup
  \caption{Schematic of the singular profile limit as $a\to0$.
  The amplitude $A_0$ equals $1$ for $r\le R$ and decreases strictly
  on the free exterior; its limiting tail is nonzero. The pressure
  converges weakly to $P_0=b\mathbf 1_{\{r<R\}}$; its value at the
  jump is immaterial. Compact support of this limiting pressure does
  not assert compact support of the profiles $Q$.
  The spectral interface conditions for amplitude and phase perturbations
  are derived in Lemma~\ref{spec:interface}.
  Heights and radial distances are schematic.}
  \label{fig:core-limit}
\end{figure}


The pressure limit in Figure~\ref{fig:core-limit} is compatible with
$A_a\to1$ in the core: the power $1/a$ retains the order-$a$ variation of
the amplitude.  Weak-star convergence, rather than differentiation of
this limiting step function, is the fact needed in the constrained mode
problem.

For each $(b,Z)\in\mathcal D$, the inner solution now has prescribed
boundary modulus $h_b$ and a determined complex logarithmic derivative
at $r_b$. We next construct a slow exterior solution with the same
boundary modulus. A constant phase rotation can then match the boundary
values; equality of the logarithmic derivatives supplies the two
remaining real matching conditions. Lemma~\ref{profile:matching} will
impose these conditions by varying $(b,Z)$.

\subsection{A uniform slow solution on the exterior}

Write
\[
 G_b(r)=H(-ib,6,-ir^2/4),\qquad
 j(b,z)=-\frac{\partial_xH(-ib,6,-iz)}{H(-ib,6,-iz)}.
\]
Proposition~\ref{free:main} says that $G_b$ never vanishes for
$r\ge R(Z)$.  For its amplitude and phase,
\begin{equation}\label{profile:free-j}
 \frac{A'}A=-\frac r2\operatorname{Im}j(b,r^2/4),\qquad
 \frac wr=\frac12+\operatorname{Re}j(b,r^2/4)<.543.
\end{equation}
At $R$, $w/r>0$ because $|j(b,Z)|\le3\cdot10^{-9}$.
The divergence equation with $c_*=6$ shows that $w$ stays positive.
For $R\le r\le2\sqrt3$, it follows that
\[
 V=b+r^2/16-w^2/4>.18.
\]
Normalize temporarily by $|G_b(R)|=1$.  Its initial amplitude derivative
has magnitude at most $(R/2)3\cdot10^{-9}$, and the equation
$(r^{11}A')'=-r^{11}VA$ makes $A'(r_b)<0$: on $[R,r_b]$ the amplitude
stays, for example, between $1/2$ and $2$, while integrating the negative
right side decreases $r^{11}A'$ by more than its initial magnitude.
It continues to decrease up to $2\sqrt3$.  Beyond that radius the strict
inequality $\operatorname{Im}j>0$ proves $A'<0$ directly.  Thus $|G_b|$
is strictly decreasing on $[r_b,\infty)$, uniformly over the disk in the
sense of strict negativity on each compact interval.  If $j(b,Z)=0$,
the same argument gives strict decrease throughout $(R,\infty)$.

For reference define the normalized free exterior
\begin{equation}\label{profile:normalized-free}
 Q^{\rm out}_0(r;b,Z)=h_b(b,Z)\frac{G_b(r)}{G_b(r_b)}.
\end{equation}
It has positive real value $h_b$ at $r_b$.  The free symbol expansion
shows that $r^{-2ib}Q^{\rm out}_0$ tends to a nonzero number $m_0(b,Z)$.
By the preceding monotonicity its modulus is at most $h_b<1$ everywhere.
Nonvanishing, the nonzero limit, and compactness of the parameter disk
therefore give uniform lower and upper bounds strictly between $0$ and
$1$ for its modulus.

\begin{lemma}[Slow exterior construction]\label{profile:exterior}
For all sufficiently small $a>0$ and $(b,Z)\in\mathcal D$, there is a
nonvanishing exterior stationary solution $Q^{\rm out}_a$ with
$Q^{\rm out}_a(r_b)=h_b$.  It depends continuously on $(a,b,Z)$, down to
the limit \eqref{profile:normalized-free}, with convergence in every
fixed derivative on compact exterior intervals.  Its factor
$f=r^{2a-2ib}Q^{\rm out}_a$ satisfies the bounds and expansions in
\eqref{profile:tail}--\eqref{profile:symbol-expansion}, uniformly over
the entire parameter disk.
\end{lemma}

\begin{proof}
Writing $\nu=-2a+2ib$ and $Q=r^\nu f$ transforms the equation into
\begin{equation}\label{profile:f-equation}
 (D+\nu)(D+\nu+10)f+\frac{ir^2}{2}Df=N_a(f),
       \qquad N_a(f)=|f|^{1/a}f.
\end{equation}
Set $N_0=0$.  Choose a small fixed neighborhood of the compact family
$(m_0(b,Z),b)$ such that the corresponding free functions $f$ stay in
an annulus
\begin{equation}\label{profile:annulus}
 0<2\kappa\le|f|\le\rho_0<1,\qquad r\ge r_b.
\end{equation}
This is possible because at $a=0$ the equation is complex linear in the
coefficient $m=\lim f$.  Enlarge the annulus slightly, still keeping its
upper radius below $1$.  On this enlarged annulus every fixed real
derivative of $N_a$ tends uniformly to zero as $a\downarrow0$: its norm
is bounded by $C_j a^{-j}\rho_1^{1/a}$ for some $\rho_1<1$ (adjusting
$C_j$ and the power of $a$ does not affect this conclusion).  Thus all
fixed differentiability orders needed below are uniformly controlled.

For prescribed $m,b$, construct a formal expansion
$f_N=\sum_{q=0}^Nm_qr^{-2q}$ with $m_0=m$.  Substituting in
\eqref{profile:f-equation} determines $m_{q+1}$ from the coefficient of
$r^{-2q}$: its coefficient in $(ir^2/2)Df_N$ is $-i(q+1)m_{q+1}$.
All the other terms at that order involve only $m_0,\ldots,m_q$ and
derivatives of $N_a$ at $m$.  This proves existence and uniqueness of
the coefficients, boundedness at every fixed order, continuous
dependence down to $a=0$, and the same bounds after a fixed number of
derivatives in the real and imaginary parts of $m$ and in $b$.
The residual of $f_N$ is $O(r^{-2N})$, with its fixed symbol and
parameter derivatives.

Here is a convergent construction behind this expansion.  Put $t=\log r$
and $Y=(e,De)$ for $e=f-f_N$.  The error equation has the form
\begin{equation}\label{profile:tail-system}
 \partial_tY=
 \begin{pmatrix}0&0\\0&-ie^{2t}/2\end{pmatrix}Y
       +\mathcal B Y+\mathcal N(t,Y)+\mathcal R_N(t),
\end{equation}
where the constant matrix $\mathcal B$ has bounded norm,
$\mathcal N(t,0)=0$, its Lipschitz constant on the annulus is bounded
uniformly in $a,m,b$, and
$|\mathcal R_N(t)|\le C_Ne^{-2Nt}$.  The propagator $\Phi(t,s)$ of the
displayed diagonal part is unitary.  The backward integral equation is
\begin{equation}\label{profile:tail-contraction}
 Y(t)=-\int_t^\infty\Phi(t,s)
       \big(\mathcal B Y(s)+\mathcal N(s,Y(s))+\mathcal R_N(s)\big)\,ds.
\end{equation}
On $t\ge t_*$ use the norm
$\|Y\|_N=\sup e^{2Nt}|Y(t)|$.  If $C$ bounds the combined Lipschitz
constant of the two terms involving $Y$, their integral has norm at
most $C/(2N)$.  Choose one integer $N_*$ with this ratio below $1/2$,
and then choose $t_*$ large enough that the corresponding ball around
$f_{N_*}$ remains in the enlarged annulus.  The residual term has norm
at most $C_{N_*}/(2N_*)$, so iteration in a ball of twice this radius
converges geometrically to a solution.  No decay estimate for an
oscillatory propagator has been used; unitarity and the weight make the
integral absolutely convergent.

The same construction works at any larger truncation order.  Two
solutions so obtained differ by $O(r^{-2N_*})$ in both entries after
subtracting their common expansion through order $N_*$.  Their
difference solves \eqref{profile:tail-contraction} without the residual,
on a common sufficiently long tail.  The contraction estimate forces
it to vanish there, and uniqueness for the finite-interval ODE makes
the solutions identical.  This proves an expansion to every order for
the single solution already constructed.  Symbol derivatives cause no
loss in this conclusion.  Indeed \eqref{profile:tail-system} bounds the
next derivative of an error by at most $C r^2$ times its two entries and
the residual; differentiating it $J$ times loses at most $2J$ powers of
$r$.  Constructing at an order greater than $N+J+2$ first therefore
gives the remainder estimate of \eqref{profile:symbol-expansion} through
the prescribed order $J$.  This also explains why the constants need
only be uniform at each fixed order.

Differentiating the integral equation with respect to $m$ or $b$ gives
the same contraction for its derivative, with a bounded inhomogeneous
term.  Difference quotients converge to that solution, proving real
$C^1$ dependence; repeated differentiation works at every fixed order.
Continuity at $a=0$ follows directly from the uniform contraction and
convergence of its coefficients.  Continue the solution from the fixed
large radius back to $r_b$.  On this finite interval the right side of
the first-order form of \eqref{profile:f-equation} and all its fixed
derivatives are bounded and converge on the enlarged annulus.  The
integral equation on finitely many short intervals, followed by
Gronwall's inequality for differences, keeps the solution close to its
free counterpart and proves existence and parameter convergence down
to $r_b$.  This also proves convergence of all fixed radial derivatives
and the first parameter derivatives there.  In view of
\eqref{profile:annulus}, the final functions satisfy
$\kappa\le|f|\le\rho<1$ for fixed $\kappa,\rho$, uniformly in small
$a$ and the chosen parameter neighborhood.

It remains to choose $m$ so that $Q(r_b)=h_b$.  At $a=0$ the map
$m\mapsto Q(r_b)$ is complex multiplication by a nonzero factor.
Its real two-dimensional inverse is uniformly bounded over the compact
parameter disk.  The just proved $C^1$ convergence implies that, on a
fixed small neighborhood of each $m_0$, the equation for $m$ is a
contraction after multiplying by that limiting inverse.  Its residual
at $m_0$ tends to zero uniformly, so it has a unique root near $m_0$.
These roots agree on overlapping neighborhoods and depend continuously
on $(a,b,Z)$.  This imposes the required boundary value.  All tail bounds
extend from the initially fixed large radius to $r_b$ by the compact
interval bounds, completing the proof.
\end{proof}

\subsection{Matching, pressure transport, and deformation}

\begin{lemma}[Matching by degree]\label{profile:matching}
For every sufficiently small $a>0$ there is $(b_a,Z_a)\in\mathcal D$
at which the inner and outer profiles have the same logarithmic
derivative at $r_b$.  After a constant phase rotation of the outer
profile they form a smooth nonvanishing stationary profile.  Every
parameter limit as $a\downarrow0$ is a zero of $j$ in $\mathcal D$.
\end{lemma}

\begin{proof}
Let $Q^{\rm in}_a$ be reconstructed with phase zero at the origin, and
define the continuous complex-valued matching map
\[
 M_a(b,Z)=\frac{(Q^{\rm in}_a)'(r_b)}{Q^{\rm in}_a(r_b)}
       -\frac{(Q^{\rm out}_a)'(r_b)}{Q^{\rm out}_a(r_b)}.
\]
Lemmas~\ref{profile:core-limit} and \ref{profile:exterior} give uniform
convergence on $\mathcal D$ to
$M_0=F_{\rm in}'(r_b)/F_{\rm in}(r_b)-G_b'(r_b)/G_b(r_b)$.
Both denominators are nonzero.

The Wronskian $\mathcal W=F_{\rm in}G_b'-F_{\rm in}'G_b$ satisfies
$\mathcal W'=-(11/r+ir/2)\mathcal W$.  At $R=2\sqrt Z$,
$G_b'(R)=(iR/2)j(b,Z)G_b(R)$, and hence
\begin{equation}\label{profile:wronskian-match}
 M_0(b,Z)=K(b,Z)j(b,Z),\qquad
 K=-\frac{iR}{2}\left(\frac R{r_b}\right)^{11}
 e^{-i(r_b^2-R^2)/4}
 \frac{G_b(R)}{F_{\rm in}(r_b)G_b(r_b)}.
\end{equation}
The factor $K$ is continuous and never zero on the entire disk.
Contracting its parameter argument to the disk center gives a homotopy
through nonzero factors to a constant.  Complex multiplication by that
constant preserves orientation.  Proposition~\ref{free:main} therefore
gives
\[
 \deg(M_0,\mathcal D^\circ,0)
       =\deg(j,\mathcal D^\circ,0)\ne0.
\]
In particular $M_0$ is separated from zero on the boundary.  Uniform
convergence preserves the boundary homotopy and the degree for every
sufficiently small $a$, so $M_a$ has a zero.

At such a zero the values have equal modulus $h_b$.  Rotate the outer
profile to match the value of the inner one; equality of logarithmic
derivatives then matches the first derivatives as well.  The ODE gives
smoothness across $r_b$, since the profile is nonzero.  To verify
smoothness through the origin without a regular-singular shortcut,
write $q(x)=Q(\sqrt x)$ for $x\ge0$.  Integrating the complex radial
equation gives
\[
 q'(x)=\frac14\int_0^1t^5e^{-ix(1-t)/4}
       \big(|q(xt)|^{1/a}-b-ia\big)q(xt)\,dt.
\]
The right side is continuous at $x=0$.  It first proves $q\in C^1$ and
then, by successive differentiation, $q\in C^\infty$ because $q$ is
bounded away from zero on this compact interval.  Thus
$Q(y)=q(|y|^2)$ is smooth at the origin.  The normalization at zero was
already imposed.  Finally, a convergent sequence of matching parameters
satisfies $M_0(b,Z)=0$, so \eqref{profile:wronskian-match} gives
$j(b,Z)=0$.  At such a root the free initial-value solution equals
$G_b/G_b(R)$, proving \eqref{profile:limit}.
\end{proof}

\begin{lemma}[The pressure derivative and positive deformation]
\label{profile:pressure}
The matched profiles satisfy \eqref{profile:coercive-bounds} and
\eqref{profile:exterior-pressure}--\eqref{profile:pressure-transport}.
\end{lemma}

\begin{proof}
First consider the inner amplitude.  On every component of $\{P<.2\}$,
\[
 (r^{11}A')'=r^{11}(P-V)A<0.
\]
At a positive left endpoint $A'\le0$, since the amplitude is entering
the sublevel set $A<.2^a$; at the origin $A'=0$.  Thus $A'<0$ at every
positive radius in this sublevel set.  If
$M=\max_{[0,r_b]}(A')_+$ is positive, it is attained at a point
$r_m\in(0,r_b)$: the derivative is zero at the origin and strictly
negative at $r_b$ for small $a$ by Lemma~\ref{profile:core-limit}.
At this point $P\ge.2$, $A''=0$, and $A'''\le0$.

Let $q=w(r_m)/r_m$.  The bounds $A\le.5^a$, $A(r_m)\ge.2^a$ and
\eqref{profile:w-ratio} give
\[
 .49<q\le q_+:=\frac{c_*}{12}(2.5)^{2a}=\frac12+O(a).
\]
The fixed lower bound follows from $h_b>.999$ and small $a$; it need
not approach $1/2$ at any prescribed rate.  Differentiating $V$ and
using \eqref{profile:amplitude} gives the exact identity
\[
 V'=r\left(\frac18-\frac{c_*}{2}\frac wr
                  +\frac{11}{2}\left(\frac wr\right)^2\right)
          +w^2\frac{A'}A.
\]
The polynomial $q\mapsto1/8-c_*q/2+11q^2/2$ is increasing for
$q\ge.49$ and small $a$, and its value at $q_+$ is $O(a)$, since its
value at $(a,q)=(0,1/2)$ is zero.  Thus at $r_m$,
\[
 V'\le Ca+w^2M/A.
\]
Differentiate the amplitude equation.  At this maximum,
\[
 (-\Delta A)'=-A'''-\frac{11}{r}A''+\frac{11}{r^2}A'\ge0,
\]
so
\[
 \frac{.2}{a}M\le VM+V'A\le CM+Ca.
\]
It follows that $M\le Ca^2$.  Since $w\ge0$ and $w/A$ is bounded on
the inner ball,
\[
 W(P/a)=\frac{wPA'}{a^2A}\le CP=Ca(P/a),
\]
which proves \eqref{profile:pressure-transport}.

In the exterior, \eqref{profile:tail} gives
$P/a=r^{-2}|f|^{1/a}/a$.  Repeated Euler differentiation and
$\kappa\le|f|\le\rho<1$ bound every fixed derivative by
$C_j a^{-N_j}\rho^{1/a}r^{-2}$ for some finite integer $N_j$.
This tends to zero uniformly after multiplication by $r^2$.
Also
$e^{2i\phi}=r^{4ib}f/\overline f$, so all its fixed Euler derivatives
are uniformly bounded.  As $w/r$ is uniformly bounded in the exterior,
these facts give \eqref{profile:exterior-pressure} and the exterior
bound on $W(P/a)$.  Together with $P\le.5$ in the interior they give
the uniform upper bound on its positive part.  By taking $a$ smaller
if necessary, the exterior formula also gives $P\le1$ everywhere.
The same bounds for $f$, and \eqref{profile:inner-bounds}, prove the
two-sided comparison for $\mu$ in \eqref{profile:coercive-bounds}.

It remains to prove positive deformation.  The two eigenvalues of the
derivative of a radial field are $w'$ and $w/r$.  For every limiting
matched profile they equal $1/2$ in the core.  On the exterior,
$A_0'<0$, $w_0>0$, and $w_0/r<.543$, so
\[
 w_0'=6-11w_0/r-2w_0A_0'/A_0\ge6-11(.543)=.027.
\]
Integrating from zero also gives $w_0/r\ge.027$.  The compact $C^1$
convergence of $A$ and formula \eqref{profile:w-ratio}, followed by
the last equation of \eqref{profile:amplitude}, give convergence of
both eigenvalues on every fixed compact interval, including zero.
For the remaining tail, the uniform expansion of Lemma~\ref{profile:exterior}
gives
\[
 \frac wr=\frac12+\frac{4b}{r^2}+O(r^{-4}),\qquad
 w'=\frac12-\frac{4b}{r^2}+O(r^{-4}),
\]
uniformly in small $a$.  Choose a large fixed radius to make these tail
bounds positive, and only then take $a$ small enough for the compact
convergence.  A subsequence contradiction, using compactness of the
parameter disk, makes this choice uniform for all matched profiles.
For example $c=.01$ is admissible after reducing $a_0$.  This proves
\eqref{profile:coercive-bounds}.
\end{proof}

\begin{proof}[Completion of Proposition~\ref{profile:existence}]
The inner construction, exterior construction and degree argument give
existence for every $0<a<a_0$ after finitely many reductions of $a_0$.
The order of choices is fixed: the disk and $\varepsilon$ first, then
the core comparison estimates, the exterior annulus and a single large
truncation order and starting radius, and finally sufficiently small
$a$.  Any later fixed derivative order is obtained by taking a higher
auxiliary truncation of the same exterior solution; it does not require
changing that solution or the already fixed geometric parameters.

The limiting statements on the inner ball are
Lemma~\ref{profile:core-limit}, and on $[r_b,\infty)$ follow from
Lemma~\ref{profile:exterior} and matching.  On a compact subset of
$(R,r_b)$, the limiting amplitude is strictly below $1$; uniform
$C^1$ convergence therefore places the profiles in a fixed annulus
bounded away from both $0$ and $1$.  All derivatives of
$z\mapsto|z|^{1/a}z$ tend to zero there.  The complex radial ODE,
successively differentiated, then proves $C^\infty$ convergence and
$P/a\to0$ on that compact set.  The weak-star pressure assertion on
any larger fixed ball follows by joining the inner assertion to this
exterior convergence.  The $C^1$ convergence of $\mu,w$ has already
been proved from their amplitude formulas.

Finally, radial Euler estimates imply Cartesian symbol estimates away
from zero: an ordinary derivative of order $j$ is a linear combination
of $r^{-j}$ times Euler derivatives with bounded smooth angular
coefficients.  The fixed-$a$ smoothness on a compact ball supplies the
remaining bounds in \eqref{profile:fixed-power-symbols}; their constants
are not asserted to be uniform as $a\downarrow0$.  As
$a-ib+\nu/2=0$,
\[
 (a-ib+D/2)Q=\tfrac12r^\nu Df,
\]
and \eqref{profile:symbol-expansion} gives
\eqref{profile:time-symbol}.  All parts of the proposition follow.
\end{proof}

\section{The outgoing spectrum}\label{sec:spectrum}

Throughout this section, $Q=Ae^{i\phi}$ is one of the profiles of
Proposition~\ref{profile:existence}, and
\[
 P=A^{1/a},\quad \mu=A^2,\quad c_*=6-2a,\quad
 w=(r/2+2\phi')\partial_r,\quad W=w\cdot\nabla,\quad D=r\partial_r.
\]
The same letter $w$ denotes the radial component of the vector field when
it appears in a one-dimensional formula. Complexification is always
complexification of a real linear equation. Thus its two circular
components $v_+,v_-$ are independent; the original real space is the
subspace $v_- =\overline{v_+}$.

\begin{theorem}[Classification of regular modes]\label{spectrum:classification}
For all sufficiently small admissible $a>0$, the only regular outgoing
modes with $\operatorname{Re}\lambda\ge-1/32$ are the phase, translation,
and time modes. Their first circular components and eigenvalues are
\[
 iQ\quad(\lambda=0),\qquad
 \partial_jQ\quad(\lambda=1/2, 1\le j\le12),\qquad
 (a-ib+D/2)Q\quad(\lambda=1).
\]
The second component of each listed real symmetry generator is the
conjugate of the displayed component. Their complex spans are the
corresponding complexified eigenspaces. In
each fixed spherical harmonic, each eigenvalue that occurs is
algebraically simple; degrees $\ell\ge2$ have no such eigenvalue.

Here ``outgoing'' means that, in each component, for every fixed integer
$j\ge0$,
\begin{equation}\label{spec:slow-condition}
 D^jv_h=O(r^{-2(a+\operatorname{Re}\lambda)})\qquad(r\longrightarrow\infty).
\end{equation}
The same classification holds for regular solutions for which
$\partial_r^kv_h\in L^2(r^{11}\,dr)$ for some integer $k\ge7$.
There is no regular generalized mode above any of these eigenvectors
with this top-derivative integrability. Algebraic simplicity for
outgoing modes means simplicity of the analytic regular/outgoing
boundary problem, including its differentiated outgoing condition.
\end{theorem}

We prove the theorem in several steps. All statements about $a\to0$
allow a convergent subsequence of the profile parameters $(b,Z)$.
Uniform conclusions for all small $a$ will follow by contradiction and
subsequence extraction, so a continuous choice of matched profiles is
not needed.

\subsection{Linearization and the boundary condition at infinity}

Let a real spherical harmonic of degree $\ell$ have unit $L^2$ norm on
$\mathbb S^{11}$. Its eigenvalue for $-\Delta_{\mathbb S^{11}}$ is
$\ell(\ell+10)$. The mode equation in this harmonic is
\begin{equation}\label{spec:circular}
 v_h''+\left(\frac{11}{r}+\frac{hi r}{2}\right)v_h'
 +\left[b+hi(a+\lambda)-\frac{\ell(\ell+10)}{r^2}\right]v_h
 =Pv_h+\frac{P}{2a}(v_h+e^{2hi\phi}v_{-h}),\qquad h=\pm1.
\end{equation}
Set
\[
 v_h=Ae^{hi\phi}(f+hi g),\qquad X=(f,g)^t,
 \qquad J=\begin{pmatrix}0&1\\-1&0\end{pmatrix}.
\]
We include the angular harmonic in $f,g$ when using spatial notation.
The coordinates $f,g$ are independent complex variables in the spectral
problem and are real on the original real space.
Cancellation of the stationary profile terms gives
\begin{equation}\label{spec:real-system}
 \mathcal H X+J(\lambda+W)X=0,\qquad
 L=-\mu^{-1}\nabla\cdot(\mu\nabla),\qquad
 \mathcal H=\begin{pmatrix}L+P/a&0\\0&L\end{pmatrix}.
\end{equation}
Indeed the terms remaining after division by $Ae^{hi\phi}$ in
\eqref{spec:circular} are $-L+hiW$ and the additional pressure
$(P/a)f$. The evolutionary generator in these variables is therefore
$\mathcal G=J\mathcal H-W$. We use repeatedly
\begin{equation}\label{spec:divergence}
 \operatorname{div}(\mu w)=c_*\mu,
 \qquad W^*=-W-c_*
\end{equation}
in the weighted volume inner product, with the appropriate boundary
term on a ball.

\begin{lemma}[Block reduction at infinity]\label{spec:block-reduction}
Suppose that
\[
 DY=(r^2\Lambda(r)+B(r))Y
\]
on a half-line, where $\Lambda$ is diagonal, its selected blocks have
uniformly separated spectra, and $\Lambda,B$ and their fixed Euler
derivatives are bounded. For every fixed $m$ there is an invertible
change $I+O_m(r^{-2})$, with the same symbol bounds, after which the
off-block remainder is $O_m(r^{-m})$. The bounded part in each diagonal
block changes by $O_m(r^{-2})$. The changes preserve holomorphic
dependence on any parameters on which the original system depends
holomorphically.
\end{lemma}

\begin{proof}
Suppose an off-block entry is $b_{jk}=O(r^{-q})$. A change with entry
$-r^{-2}b_{jk}/(\Lambda_j-\Lambda_k)$ cancels it by commutation with
$r^2\Lambda$. Differentiating this change or multiplying it by the
bounded part gives $O(r^{-q-2})$; its products with the other changes
have the same or better order. The denominator and all its fixed
Euler derivatives are uniformly bounded by the spectral gap.
Perform this operation simultaneously on every off-block entry, and
iterate a finite number of times. The product of the changes is
invertible for sufficiently large $r$ by its $I+O_m(r^{-2})$ bound.
Differentiating the construction proves the symbol estimates. All
operations are algebraic operations and differentiation with nonzero
denominators, proving the last assertion.
\end{proof}

\begin{lemma}[Slow bases and exclusion of fast modes]\label{spec:slow-basis}
For bounded $\lambda,\ell$ and small $a$, the exterior solution space
of \eqref{spec:circular} has a two-dimensional outgoing subspace,
locally holomorphic in $\lambda$. Its bases and parameter derivatives
converge locally on $r\ge r_b$ to free outgoing bases as $a\to0$.
Every homogeneous solution satisfying the top-derivative condition
in Theorem~\ref{spectrum:classification} belongs to this subspace.
\end{lemma}

\begin{proof}
Use $(v_h,v_h'/r)$ for each channel. The coefficient of $r^2$ in its
$D$ equation is
\[
 \begin{pmatrix}0&1\\0&-hi/2\end{pmatrix}.
\]
Diagonalize by the eigenvectors with first component one, group the
two zero eigenvalues, and apply Lemma~\ref{spec:block-reduction}.
Writing
\[
 \nu_h=-2(a+\lambda)+2hib,
\]
the slow block is $\operatorname{diag}(\nu_h)+O(r^{-2})$, and the fast
entries are $-hi r^2/2-12-\nu_h+O(r^{-2})$.
For example, before diagonalization the bounded contribution to the
second row is $-12(v_h'/r)-(b+hi(a+\lambda))v_h$; this verifies both
constant exponents.

Factor out $r^{\nu_h}$ in a slow coordinate and
$e^{-hi r^2/4}r^{-12-\nu_h}$ in a fast coordinate. The diagonal errors
are integrable in $d r/r$. The off-block errors are also integrable if
$m>|12+2\operatorname{Re}\nu_h|$ with a fixed positive margin on the
parameter neighborhood. The factored equation is thus an integrable
perturbation of a constant vector. On a sufficiently late half-line,
its integral from infinity is a contraction on bounded vectors;
prescribing its four limits gives four independent solutions.
Continuation gives a basis on the entire exterior. This construction
is holomorphic on a slightly smaller parameter neighborhood because
the integrable majorants are uniform there.

Setting both fast limits to zero gives exactly the outgoing subspace.
In fact the fast coordinates forced by slow coordinates are
$O(r^{\operatorname{Re}\nu_h-m})$ before factoring, by the backward
integral. Taking $m$ larger before each fixed number of differentiations
gives \eqref{spec:slow-condition}. Increasing the reduction order does
not change the limiting normalization: the additional transformations
have arbitrarily high order relative to the prescribed normalization,
and uniqueness of the backward integral applies on a common tail.
Uniform exterior symbol bounds for the profiles and their convergence
give convergence of these integral equations by dominated convergence;
ordinary ODE continuation then proves the asserted local convergence.
At $a=0$ the slow columns are, up to their nonzero normalizations,
\begin{equation}\label{spec:free-bases}
 r^\ell H(\lambda+\ell/2-hib,6+\ell,-hi r^2/4)
\end{equation}
in their separate channels, by Lemma~\ref{free:slow}.

If the fast limit $d_h$ in channel $h$ is nonzero, reconstruction and
$k$ radial differentiations give the leading fast term
\begin{equation}\label{spec:fast-derivative}
 \partial_r^k v_h=
 (-hi/2)^k e^{-hi r^2/4}r^{k-12-\nu_h}(d_h+o(1))
 +O(r^{\operatorname{Re}\nu_h-k}).
\end{equation}
The other channel's fast contribution is smaller by $r^{-2}$; this
follows from the off-block reconstruction. The same formula follows
directly by successively applying $D$ to the reduced fast equation,
whose leading multiplier is $-hi r^2/2$.
For $\operatorname{Re}\lambda\ge-1/32$ and $k\ge7$,
$2k>12+2\operatorname{Re}\nu_h$. Consequently the first term in
\eqref{spec:fast-derivative} dominates the slow term and its square is
not integrable against $r^{11}\,dr$. Thus every fast limit vanishes.
For any particular finite $\lambda,\ell$, the argument applies in a
bounded neighborhood of those parameters; it does not require a
uniform bound on the parameters of the mode being tested.
\end{proof}

\begin{lemma}[Upper bound for the real part]\label{spec:real-bound}
For all sufficiently small $a$, there is no nonzero regular outgoing
mode with $\operatorname{Re}\lambda\ge4$.
\end{lemma}

\begin{proof}
Put $V_a=P/a$ and $\sigma=\operatorname{Re}\lambda$. Test
\eqref{spec:real-system} against $(\lambda+W)X$, and take real parts.
The $J$ term has zero real part. Integration by parts, using
\eqref{spec:divergence}, gives
\begin{equation}\label{spec:upper-energy}
 0=\int\left\{(\sigma-c_*/2)(|\nabla X|^2+V_a|f|^2)
 +\nabla X^*(\nabla w)\nabla X
 -\tfrac12 W(V_a)|f|^2\right\}\mu\,dy.
\end{equation}
Here the matrix $\nabla w$ acts on spatial indices and the expression
is summed over the two components. Boundary terms vanish: after
division by $A$, the slow bound is $X=O(r^{-2\sigma})$, with its symbol
derivatives, and $\sigma\ge4$ gives all needed integrability.

For completeness, the weighted Hardy inequality used here is uniform
for small $a$:
\begin{equation}\label{spec:hardy}
 \int r^{-2}|f|^2\mu\,dy\le C\int|\nabla f|^2\mu\,dy.
\end{equation}
Indeed, for $\rho=\langle r\rangle^{-4a}$,
$\operatorname{div}(\rho y/r^2)\ge(10-4a)\rho/r^2$.
Integration against $|f|^2$ and Cauchy--Schwarz prove the inequality
with $\rho$; comparability of $\rho$ and $\mu$ proves
\eqref{spec:hardy}. Cutoffs justify the calculation by density.
Proposition~\ref{profile:existence} gives $\nabla w\ge cI$ and
$W(V_a)\le C a V_a$ on the inner ball, while its exterior positive
part is $o(1)r^{-2}$. Since $\sigma-c_*/2\ge1$, these two negative
contributions in \eqref{spec:upper-energy} are absorbed, respectively,
by the pressure and gradient terms. A nonzero $X$ is impossible.
\end{proof}

\subsection{An outgoing estimate for unbounded parameters}

We next exclude loss of modes to large angular or imaginary spectral
parameters. The main step is to express the outgoing condition as a relation
between a mode and its radial derivative on one fixed exterior shell.
Once this relation is available, interior energy and flux identities will
exclude a regular mode. Consider a sequence with $a\to0$,
$-1/32\le\sigma\le4$, and
$\ell+|\operatorname{Im}\lambda|\to\infty$. Complex conjugation of
\eqref{spec:real-system} allows us to assume
$\omega=\operatorname{Im}\lambda\ge0$. After extraction, one of the
following holds:
\begin{equation}\label{spec:cases}
 \text{I: }\frac{\ell}{\sqrt{1+\omega}}\longrightarrow\infty;
 \qquad
 \text{II: }\omega\longrightarrow\infty,\quad
 \frac{\ell}{\sqrt\omega}\text{ is bounded}.
\end{equation}
Set
\begin{equation}\label{spec:high-definitions}
 \begin{split}
 q_h&=r^{11/2}e^{hi r^2/8}v_h,\qquad
 L_\ell=(\ell+5)^2-\tfrac14,\qquad n_* =\max(1+\ell,\sqrt\omega),\\
 F_h&=r^2/16+b-h\omega-L_\ell/r^2,
 \qquad S=\max(1+\ell,\omega),\qquad E=16\sqrt S.
 \end{split}
\end{equation}
The symbol $q_h$ here denotes an unknown function, not the first
argument of $H$. Direct substitution gives
\begin{equation}\label{spec:schrodinger-radial}
 q_h''+(F_h+hi\gamma)q_h=\sum_j V_{hj}(r)q_j,
 \qquad \gamma=\sigma+a-3,\qquad |V_{hj}(r)|\le Cr^{-2}
 \quad(r\ge r_b).
\end{equation}
The bound, uniform in the sequence, follows from the exterior pressure
estimates; oscillations in $V_{hj}$ cause no difficulty below.

Choose a fixed $B_0>r_b$, and put $B_1=B_0+1$. In Case II take $B_0$
large enough that $L_\ell/B_0^2<\omega/2$ eventually. On this shell
$F_h<0$ for both channels in Case I; in Case II $F_+<0<F_-$.
In either case $|F_h|\asymp n_*^2$, uniformly on the shell.

\begin{lemma}[Outgoing data at the remote endpoint]\label{spec:remote-data}
Every outgoing solution along this sequence satisfies
\begin{equation}\label{spec:remote-robin}
 q_h'(E)-hi\sqrt{F_h(E)}q_h(E)=O(E^{-1})|q(E)|.
\end{equation}
The constant is uniform after discarding finitely many sequence terms.
\end{lemma}

\begin{proof}
For $r\ge E$ use $(v_h,v_h'/r)$ again. Diagonalize the leading matrix
\[
 \begin{pmatrix}
 0&1\\ h\omega/r^2+\ell(\ell+10)/r^4&-hi/2
 \end{pmatrix}.
\]
Its eigenvalues are
\[
 hi\left(-\tfrac14\pm
 \sqrt{\tfrac1{16}-h\omega/r^2-\ell(\ell+10)/r^4}\right).
\]
The square root stays uniformly positive on $r\ge E$. Eigenvectors
with first entry one give a uniformly invertible symbol change.
Group the two plus signs as the outgoing block; take the two minus
signs as separate blocks. These blocks have uniformly positive gaps.
The bounded remainder has norm at most a fixed $C_0$, with uniform
fixed symbol derivatives. The leading diagonal is purely imaginary.

The order of choices is important. First choose $C>C_0+2$, enlarged
if necessary to bound the Hermitian parts of the initial full and
block systems. Next choose one fixed integer $m>2C+4$. Apply
Lemma~\ref{spec:block-reduction} to this order. Finally pass far enough
along the sequence that its finitely many $O_m(r^{-2})$ corrections
have norm at most one on $r\ge E$. Both the full reduced evolution
and each block evolution then have propagator bound
\begin{equation}\label{spec:power-propagator}
 \|\Phi(r,t)\|\le
 \left(\frac{\max(r,t)}{\min(r,t)}\right)^C.
\end{equation}
This is Gronwall in $\log r$ after using skewness of the leading
diagonal. No separation inside the outgoing block is required.

For each individual outgoing solution, all reduced coordinates have
symbol bounds $O(r^\rho)$, where $\rho=-2(a+\sigma)<1$; their constants
may depend on that solution and its parameters. In an incoming
equation the $r^2$ leading coefficient is invertible. Solving for that
coordinate and differentiating gives a gain of two powers of $r$;
iteration gives the same gain for each fixed symbol derivative, up
to the $O_m(r^{-m})$ forcing from the outgoing block. Thus the incoming
coordinate is $o(r^{-C-1})$, with $m$ as chosen. This qualitative
bound is used only to remove its terminal value in variation of
constants at infinity.

After that terminal value has vanished, the quantitative estimate
uses only \eqref{spec:power-propagator}. If $Y$ denotes all reduced
coordinates, its incoming part at $E$ is at most
\[
 C_m|Y(E)|\int_E^\infty(r/E)^{2C}r^{-m}\frac{dr}{r}
 =\frac{C_mE^{-m}}{m-2C}|Y(E)|.
\]
Undoing the changes gives
\[
 \frac{v_h'(E)}E=
 hi\left(-\tfrac14+
 \sqrt{\tfrac1{16}-h\omega/E^2-\ell(\ell+10)/E^4}\right)v_h(E)
 +O(E^{-2})|v(E)|.
\]
The full coordinate size is bounded by $C|v(E)|$, since the incoming
coordinates are small and the first entry of each outgoing
eigenvector is one. Multiplication by the scalar defining $q_h$
proves \eqref{spec:remote-robin}; the $11/(2E)$ term and the difference
between $L_\ell$ and $\ell(\ell+10)$ are within its error.
\end{proof}

\begin{lemma}[Scalar transfer and a turning point]\label{spec:scalar-turn}
Fix one channel and write $F=F_h$. For $R\in[B_0,B_1]$, let $U$ solve
\[
 U''+(F+hi\gamma)U=0,\qquad
 U(E)=F(E)^{-1/4},\quad U'(E)=hi\sqrt{F(E)}U(E).
\]
If $F$ has a zero $r_0\in[R,E]$, define
\begin{equation}\label{spec:turn-scales}
 \begin{split}
 A_0&=r_0^2/16,\qquad H_0=L_\ell/r_0^2,\qquad k_0^2=A_0+H_0,\\
 d_0&=(r_0/(2k_0^2))^{1/3},\qquad
 K(r)=(|F(r)|+d_0^{-2})^{1/2},\\
 H_f(r)&=\int_{\min(r,r_0)}^{r_0}\sqrt{-F(t)}\,dt.
 \end{split}
\end{equation}
If there is no zero, put $K=\sqrt F$ and $H_f=0$.
For $N_r(u)=K(r)^{1/2}|u(r)|+K(r)^{-1/2}|u'(r)|$, after taking a
subsequence common to the two channels, there are $c,C>0$ such that
\begin{align}
 N_r(u)&\le C e^{|H_f(r)-H_f(t)|}N_t(u),\label{spec:scalar-transfer}\\
 |U(R)|&\ge cK(R)^{-1/2}e^{H_f(R)},\label{spec:scalar-value}\\
 U'(R)/U(R)&=\mathcal B_h(R)+o(n_*),\label{spec:scalar-slope}
 \end{align}
where
\[
 \mathcal B_h(R)=
 \begin{cases}-\sqrt{-F_h(R)},&F_h(R)<0,\\
 hi\sqrt{F_h(R)},&F_h(R)>0.
 \end{cases}
\]
The first estimate holds for every scalar homogeneous solution and
all $r,t\in[R,E]$. All left endpoint statements are uniform for
$R\in[B_0,B_1]$.
\end{lemma}

\begin{proof}
The derivative $F'=r/8+2L_\ell/r^3$ is positive, so there is at most
one zero. In the turning cases its defining equation gives
\begin{equation}\label{spec:turn-potential}
 F(r)=A_0((r/r_0)^2-1)+H_0(1-(r_0/r)^2),\qquad
 A_0-H_0=h\omega-b.
\end{equation}
It follows that $k_0^2\asymp S$, $r_0\to\infty$, $r_0<E/2$, and
$d_0\to0$. For the $h=-1$ turn in Case I, for example,
$r_0^2\asymp \ell^2/(\omega+\ell)$, which tends to infinity precisely
because $\ell/\sqrt{1+\omega}\to\infty$. The other turn has
$r_0^2\asymp\omega+\ell$. In the no-turn case, necessarily the
negative channel in Case II, $F\asymp S$ throughout $[R,E]$.

Fix a large number $M$ and remove the interval $|r-r_0|<Md_0$.
On either remaining interval choose $\operatorname{Re}p>0$ with
$p^2=\operatorname{sign}(F)(F+hi\gamma)$. Use the two approximate
solutions
\[
 p^{-1/2}\exp(\pm\!\int p)\quad(F<0),\qquad
 p^{-1/2}\exp(\pm i\!\int p)\quad(F>0).
\]
All integrals in this paragraph are over the outer intervals. Direct
differentiation of \eqref{spec:turn-potential} gives
\begin{align}
 \int |F|^{-1/2}\,dr&\le C,\label{spec:phase-error}\\
 \sup\frac{|F'|}{|F|^{3/2}}+
 \int\left(\frac{|F''|}{|F|^{3/2}}
       +\frac{|F'|^2}{|F|^{5/2}}\right)dr
 &\le CM^{-3/2}+o(1).\label{spec:wkb-error}
\end{align}
Here and below the constants multiplying $M^{-3/2}$ do not depend on
$M$. To verify the estimates, on $[r_0/2,2r_0]$ use
$|F|\asymp k_0^2|r-r_0|/r_0$,
$|F'|\le Ck_0^2/r_0$, $|F''|\le Ck_0^2/r_0^2$ and
$d_0^3=r_0/(2k_0^2)$. Integrating the powers
$|r-r_0|^{-3/2}$ and $|r-r_0|^{-5/2}$ from $Md_0$ gives the stated
error. Off this interval $|F|\gtrsim k_0^2$, with relative derivative
bounds $C/r,C/r^2$, and $E\lesssim k_0$; splitting at $r_0/2$ and
$2r_0$ gives a vanishing error and the bounded integral in
\eqref{spec:phase-error}. For no turn the same computations use
$F\asymp S$ and give $o(1)$ in \eqref{spec:wkb-error} on all of $[R,E]$.

Here is a transfer justification retaining the error that will be
needed below. Each approximate solution has residual potential
\[
 \frac{(p^{-1/2})''}{p^{-1/2}}
 =\frac34(p'/p)^2-\frac12p''/p.
\]
The approximate pair has constant nonzero Wronskian. Its transfer
matrix in the $N$ norms is at most
$C\exp(|H_f(r)-H_f(t)|)$: $K\asymp|p|$, and
$\int|p-\sqrt{|F|}|\,dr\le C$ by \eqref{spec:phase-error}.
Variation of constants for the residual followed by Gronwall bounds
the difference between exact and approximate transfer by
\begin{equation}\label{spec:transfer-error}
 (CM^{-3/2}+o(1))e^{|H_f(r)-H_f(t)|}.
\end{equation}
Indeed the residual divided by $K$ has integral bounded by
\eqref{spec:wkb-error}; when $r,s,t$ occur in order, monotonicity of
$H_f$ gives
$|H_f(r)-H_f(s)|+|H_f(s)-H_f(t)|=|H_f(r)-H_f(t)|$.
This removes the exponential factor from the integral inequality.
With no turn the error coefficient is $o(1)$.

At $E$, the prescribed data select the outgoing approximate branch
up to an error tending to zero. On the positive outer interval,
\eqref{spec:phase-error} bounds its normalized amplitude above and
below, because the absolute imaginary action is bounded. Consequently
\begin{equation}\label{spec:positive-branch}
 c\le F^{1/4}|U|\le C,\qquad
 \frac{U'}{\sqrt F\,U}=hi+O(M^{-3/2})+o(1).
\end{equation}
This proves all claims in the no-turn case.

In a turning case set $\xi=(r-r_0)/d_0$ and
$Y(\xi)=d_0^{-1/2}U(r_0+d_0\xi)$. Since
$F'(r_0)d_0^3=1$ and $\gamma d_0^2\to0$, its equation converges on
every fixed $\xi$ interval to
\begin{equation}\label{spec:airy-equation}
 Y_\infty''+\xi Y_\infty=0.
\end{equation}
The corresponding transfer on $[-M,M]$ is bounded in the scaled data
norm, with a bound depending only on fixed $M$.
Figure~\ref{fig:turning} records these three regions.
\begin{figure}[tbp]
  \centering
  \begingroup
  \definecolor{turnblue}{RGB}{48,91,120}
  \definecolor{turnorange}{RGB}{163,92,43}
  \resizebox{0.95\linewidth}{!}{%
  \begin{tikzpicture}[
    x=1cm,y=1cm,font=\small,
    guide/.style={densely dashed,thin,black!45},
    every node/.style={inner sep=2pt}]
    \fill[turnorange!9] (.6,0) rectangle (5,2.1);
    \fill[black!6] (5,0) rectangle (9,2.1);
    \fill[turnblue!7] (9,0) rectangle (13.4,2.1);
    \draw[guide] (5,0) -- (5,2.1);
    \draw[guide] (9,0) -- (9,2.1);
    \draw[guide] (7,0) -- (7,1.3);
    \node[align=center] at (2.8,1.7) {Forbidden region\\$F<0$};
    \node[align=center] at (7,1.7) {Airy window\\$|r-r_0|\le M d_0$};
    \node[align=center] at (11.2,1.7) {Oscillatory region\\$F>0$};
    \draw[<-,very thick,turnorange] (1.1,.65) -- (4.5,.65);
    \node[above] at (2.8,.65) {growing branch to the left};
    \node[align=center,fill=black!6] at (7,.65) {$\xi=(r-r_0)/d_0$\\$Y_\infty''+\xi Y_\infty=0$};
    \draw[->,very thick,turnblue] (9.6,.65) -- (12.95,.65);
    \node[above] at (11.25,.65) {outgoing branch};
    \draw[->,thin,black!75] (.25,0) -- (14,0) node[right] {$r$};
    \foreach \x in {.6,5,7,9,13.4}
      \draw (\x,.08) -- (\x,-.09);
    \node[below] at (.6,-.1) {$R$};
    \node[below] at (5,-.1) {$r_0-Md_0$};
    \node[below] at (7,-.1) {$r_0$};
    \node[below] at (9,-.1) {$r_0+Md_0$};
    \node[below] at (13.4,-.1) {$E$};
    \node[align=center] at (7,-1.05)
      {Normalized Airy limit: $\displaystyle
       \operatorname{Im}\!\left(\overline{Y_\infty}Y_\infty'\right)\ne0$};
    \node[align=center] at (7,-1.6)
      {$\text{nonzero flux}\ \Longrightarrow\ \text{nonzero growing component on the forbidden side}$};
  \end{tikzpicture}%
  }
  \endgroup
  \caption{Schematic of a channel with one turning point in the
  high-parameter estimate, with $R\in[B_0,B_1]$ as in
  Lemma~\ref{spec:scalar-turn}. The scalar solution is selected at $E$ by
  $U(E)=F(E)^{-1/4}$ and $U'(E)=h i\sqrt{F(E)}\,U(E)$.
  For a fixed large $M$, pass to a subsequence and take the
  high-parameter limit in the window using $Y=d_0^{-1/2}U$
  and $\xi=(r-r_0)/d_0$.
  The resulting Airy solution has constant nonzero flux, so it contains
  a branch growing towards $R$.
  Only after this fixed-window limit is $M$ allowed to tend to infinity
  to obtain the normalized logarithmic derivative
  $U'(R)/(\sqrt{-F(R)}\,U(R))\to-1$.
  Distances are not to scale, and the arrows indicate branch selection,
  not a graph of $U$ or a sign assertion for its complex values.}
  \label{fig:turning}
\end{figure}


Choose one fixed $M_0$ so large that the error in
\eqref{spec:positive-branch} is below its lower amplitude margin.
At $\xi=M_0$, the scaled Cauchy data are bounded and their flux
\[
 \operatorname{Im}(\overline Y Y')
 =\operatorname{Im}(\overline U U')
\]
has sign $h$ and magnitude bounded below. A subsequence of these data
converges, and continuous dependence of the scaled ODE gives one
solution $Y_\infty$ on every fixed interval. Its flux is nonzero and
constant, since \eqref{spec:airy-equation} has real coefficients.

We spell out why this gives the required \emph{relative} margin on
the forbidden side. With $t=-\xi$, the real equation is $y''-ty=0$.
On a sufficiently late half-line it has a positive growing solution
comparable to $t^{-1/4}\exp(2t^{3/2}/3)$. To construct one, prescribe
the growing approximate data at a fixed large $T$ and use the residual
transfer estimate with $p=\sqrt t$: the total error relative to the
growing scale is at most $CT^{-3/2}$. The exact solution and its
derivative are positive when $T$ is large. Its logarithmic slope
$z=y'/(\sqrt t\,y)$ satisfies
\[
 z'=\sqrt t(1-z^2)-\frac{z}{2t}.
\]
This equation implies $z\to1$: for every fixed $\epsilon>0$, its
derivative points into $[1-\epsilon,1+\epsilon]$ at both endpoints
for sufficiently large $t$, and outside that interval it points
inwards with magnitude bounded below by a positive multiple of
$\sqrt t$. A decaying independent solution is obtained by multiplying
the growing one by $\int_t^\infty y(s)^{-2}\,ds$. The growing
comparison bound makes that integral finite and shows that this
second solution and its derivative are negligible relative to the
first. These real solutions have nonzero Wronskian. A complex multiple
of the decaying solution alone has zero flux. Thus the coefficient of
the growing solution in $Y_\infty$ is nonzero.

As $M\to\infty$ its growing contribution dominates its decaying one,
and
\[
 \frac{Y_\infty'(-M)}{\sqrt M\,Y_\infty(-M)}\longrightarrow-1.
\]
In the approximate basis referenced at $r_-=r_0-Md_0$, the coefficient
of growth to the left therefore has magnitude at least a fixed
positive fraction of $N_{r_-}(U)$, for every sufficiently large fixed
$M$ and subsequently late sequence indices. That fraction can be
chosen independently of $M$: the normalized basis vectors have
comparable size one at their reference point, and the displayed
slope is separated from the decaying slope $+1$. For each such fixed
$M$, $N_{r_-}(U)$ is also bounded below by a positive constant.

Transfer to $R$ now gives a growing term of size at least
$cN_{r_-}(U)e^{H_f(R)-H_f(r_-)}$; the decaying term is negligible since
the action difference tends to infinity. The transfer error
\eqref{spec:transfer-error}, relative to this growing term, is still
$CM^{-3/2}+o(1)$ with $C$ independent of $M$, precisely because of the
relative coefficient margin. For the value lower bound, choose one
sufficiently large fixed $M$ and retain it; then $H_f(r_-)=O_M(1)$
and \eqref{spec:scalar-value} follows with a positive constant.
For the slope take an arbitrary sufficiently large fixed $M$, obtain
error $CM^{-3/2}+o(1)$, and then let $M\to\infty$. This proves
\eqref{spec:scalar-slope}. The transfer estimate follows from the two
outer estimates and the scaled central transfer, using one fixed
$M$. Both channels can use a common extracted subsequence. Finally,
all estimates allow a moving $R$ in the fixed shell, where
$|F_h(R)|\asymp n_*^2$. If uniformity there failed, the same argument
at a sequence of offending endpoints would contradict the estimates.
\end{proof}

\begin{lemma}[Coupled outgoing data on a fixed shell]\label{spec:boundary-estimate}
After subselection, every outgoing solution along an escaping
parameter sequence satisfies, uniformly on $B_0\le R\le B_1$,
\begin{equation}\label{spec:shell-robin}
 q_h'(R)=\mathcal B_h(R)q_h(R)+o(n_*)|q(R)|.
\end{equation}
The error is an operator bound on the two-dimensional outgoing space.
\end{lemma}

\begin{proof}
For each scalar channel let $D_h(R)=0$, $D_h'(R)=K_h(R)^{1/2}$.
Lemma~\ref{spec:scalar-turn} gives
\[
 N_r(D_h)\le Ce^{H_h(R)-H_h(r)},\quad
 N_r(U_h)\le Ce^{H_h(r)},\quad
 |\mathcal W(D_h,U_h)|\ge ce^{H_h(R)},
\]
where $H_h=H_f$ in that channel. The Green kernel with left
Dirichlet data and the right Robin condition of $U_h$ is the product
of $D_h$ on the left and $U_h$ on the right, divided by the Wronskian.
It obeys the following two different action bounds:
\begin{equation}\label{spec:green-actions}
 N_r(G_h(r,t))\le C K_h(t)^{-1/2}
 \begin{cases}
 e^{H_h(t)-H_h(r)},&r<t,\\
 e^{H_h(r)-H_h(t)},&r>t.
 \end{cases}
\end{equation}
Each exponential is at most one, by monotonicity within its own
channel. There is no comparison of $H_+$ and $H_-$. A right Robin
derivative error $\beta_h$, with zero left data, similarly gives
$N_r\le CK_h(E)^{-1/2}|\beta_h|$; this follows either from the same
Wronskian or by using $D_h$ for the correction.

Let $\kappa=\min_{h,r\in[R,E]}K_h(r)$. Then $\kappa\to\infty$:
at a turn it is $d_0^{-1}$, and in the no-turn case it is comparable
to $\sqrt S$. Put $\mathcal N(q)=\max_h\sup_rN_r(q_h)$.
The matrix forcing in \eqref{spec:schrodinger-radial} has Green norm
at most
\begin{equation}\label{spec:coupling-norm}
 C\max_{h,j}\int_R^E
 \frac{t^{-2}\,dt}{\sqrt{K_h(t)K_j(t)}}
 \le\frac{C}{R\kappa}=o(1).
\end{equation}
Also $K_h(E)\asymp E$ for both channels, so the right error
\eqref{spec:remote-robin} has norm at most $CE^{-2}\mathcal N(q)$.
The diagonal outgoing extension
$q_h(R)U_h(r)/U_h(R)$ has norm at most
$Cn_*^{1/2}|q(R)|$, since $K_h(R)\asymp n_*$ and
$H_h(r)\le H_h(R)$. Variation of constants and absorption of
\eqref{spec:coupling-norm} and the Robin error give
\[
 \mathcal N(q)\le Cn_*^{1/2}|q(R)|,\qquad
 \mathcal N\left(q-\left(q_h(R)U_h/U_h(R)\right)_h\right)
 =o(1)n_*^{1/2}|q(R)|.
\]
Taking the derivative at $R$ costs $K_h(R)^{1/2}\asymp n_*^{1/2}$.
Together with \eqref{spec:scalar-slope}, this is
\eqref{spec:shell-robin}. Every estimate used is homogeneous in the
left data and uniform over the shell.
\end{proof}

\begin{lemma}[Exclusion of escaping parameters]\label{spec:high-exclusion}
There is no sequence of nonzero regular outgoing modes with $a\to0$,
$-1/32\le\operatorname{Re}\lambda\le4$, and
$\ell+|\operatorname{Im}\lambda|\to\infty$.
\end{lemma}

\begin{proof}
Use \eqref{spec:cases} and Lemma~\ref{spec:boundary-estimate}.
All volume pairings in this proof have weight $\mu$, and $[\cdot]$
denotes its surface integral at the current sphere, including the
factor $r^{11}$. On a fixed ball $\mu$ is uniformly bounded above and
below, and the coefficients of $W$ and $L+\Delta$ are uniformly
bounded.

In Case I pair \eqref{spec:real-system} with $X$ on the ball $r<B_1$
and take real parts. The second-order term is
\[
 \int(|\nabla X|^2+(P/a)|f|^2)\mu-\operatorname{Re}[X^*X'].
\]
The boundary term has nonpositive real part, by
\eqref{spec:shell-robin} and both negative signs of $F_h$.
More precisely its real part is
$\frac12\operatorname{Re}\sum_h\overline q_hq_h'+O(|q|^2)$, so its
leading value is $-c n_*|q|^2$ and absorbs its errors.
The absolute value of the remaining real terms is at most
$C\|X\|_2\|\nabla X\|_2+C(1+\omega)\|X\|_2^2$.
The angular barrier
\[
 \|\nabla X\|_2^2\ge
 \ell(\ell+10)B_1^{-2}\|X\|_2^2
\]
and $\omega=o(\ell^2)$ absorb these terms into an arbitrarily small
fraction of the gradient term. Nonnegative pressure then gives a
contradiction.

In Case II normalize
\begin{equation}\label{spec:case-II-normalization}
 \int_{r<B_1}(|X|^2+|\nabla X|^2)\mu=1.
\end{equation}
Write the shell error in \eqref{spec:shell-robin} as
$\varepsilon n_*|q|$, with $\varepsilon\to0$. Its leading diagonal
has modulus at least $c n_*$, so $n_*|q|\le C|q'|$ on the shell.
There is a point in its first half at which both $|q'|$ and
$n_*|q|$ are bounded, by the shell integral bound furnished by
\eqref{spec:case-II-normalization}.
For the channel $h=-1$, the energy
$|q_-'|^2+F_-|q_-|^2$ propagates this bound to $B_1$.
Indeed $|F_-'|/F_-\le C$, $F_-\asymp n_*^2$ on the shell, and the
forcing in its equation has bounded integral, since it is at most
$C|q|$ there. Differentiating the square root of this energy and
applying Gronwall proves
\[
 |q_-'|+n_*|q_-|\le C
\]
from that point to $B_1$.
The positive channel then obeys, in the upper-derivative sense,
\[
 \frac{d}{dr}|q_+|
 \le-(\sqrt{-F_+}-\varepsilon n_*)|q_+|+C\varepsilon.
\]
The distance from the chosen point to $B_1$ is at least $1/2$.
Integrating this inequality with $\sqrt{-F_+}\ge c n_*$ gives
\begin{equation}\label{spec:plus-vanishing}
 n_*q_+(B_1)\longrightarrow0,\qquad q_+'(B_1)\longrightarrow0.
\end{equation}

We next remove the other boundary channel by a flux identity.
The imaginary part of the pairing of \eqref{spec:real-system} with
$X$ is
\begin{equation}\label{spec:flux}
 \operatorname{Im}[X^*X']
 =\frac{w(B_1)}2\operatorname{Im}[X^*JX]
 +(\sigma-c_*/2)\int\operatorname{Im}(X^*JX)\mu.
\end{equation}
To check the first-order term, integration by parts and
\eqref{spec:divergence} give
\[
2\operatorname{Im}\int X^*JW X\,\mu
=w(B_1)\operatorname{Im}[X^*JX]-c_*\int\operatorname{Im}(X^*JX)\mu.
\]
Also the pressure-free equation $Lg=(\lambda+W)f$, tested against $g$,
gives
\begin{equation}\label{spec:cross-small}
 \left|\int\overline g f\,\mu\right|\le C/\omega.
\end{equation}
Its integrated gradient term, $W$ term, and boundary term are bounded
by the normalization and the boundary bounds already obtained.
The change to circular coordinates gives the exact cancellation
\[
 \operatorname{Im}[X^*X']
 -\frac{w(B_1)}2\operatorname{Im}[X^*JX]
 =\frac12\sum_h\operatorname{Im}(\overline q_hq_h')(B_1).
\]
Multiplying \eqref{spec:flux} by $n_*\asymp\sqrt\omega$, and using
\eqref{spec:cross-small}, makes its right volume term tend to zero.
The positive channel contributes $o(1)$ by
\eqref{spec:plus-vanishing}. The outgoing relation for the other
channel gives
$\operatorname{Im}(\overline q_-q_-')=-\sqrt{F_-}|q_-|^2+
o(n_*)|q_-||q|$. Its sign and $F_-\asymp n_*^2$ imply
\begin{equation}\label{spec:all-boundary-small}
 n_*|q(B_1)|+|q'(B_1)|\longrightarrow0.
\end{equation}

Rellich compactness on the ball, including the full angular fields,
gives a strongly $L^2$ convergent subsequence of $f$. The equation
$Lg-Wf=\lambda f$ has uniformly bounded left side in $H^{-1}$ by
the normalization. Division by $\lambda$, whose modulus tends to
infinity, shows that the limit is zero. Thus
\begin{equation}\label{spec:f-small}
 f\longrightarrow0\quad\hbox{strongly in }L^2(r<B_1).
\end{equation}

Finally put $S_d=W+c_*/2$ and test \eqref{spec:real-system} against
$S_dX$. The real part of the second-order and pressure terms is
\begin{align*}
 &\int\left(\nabla X^*(\nabla w)\nabla X
       -\tfrac12W(P/a)|f|^2\right)\mu\\
 &\hspace{8mm}
 +\frac{w(B_1)}2[|\nabla X|^2+(P/a)|f|^2]
 -\operatorname{Re}[(S_dX)^*X'].
\end{align*}
All boundary terms vanish by \eqref{spec:all-boundary-small}; the
tangential derivatives cost $\ell\le n_*$, and $P/a$ is bounded on
this exterior sphere. The $JS_dX$ pairing has zero real part. For
$I=\int(S_dX)^*JX\,\mu$, \eqref{spec:f-small} and one integration
by parts in the term containing $Wf$ give $I\to0$. Further,
\[
 2\operatorname{Im}I=w(B_1)\operatorname{Im}[X^*JX].
\]
Hence $\omega\operatorname{Im}I\to0$, since
$\omega=O(n_*^2)$ and the boundary values are $o(n_*^{-1})$.
The remaining pairing $(\lambda-c_*/2)I$ therefore has real part
tending to zero. We conclude that the displayed bulk integral tends
to zero. The profile estimates $\nabla w\ge cI$ and $W(P/a)\le C$,
together with \eqref{spec:f-small}, force $\|\nabla X\|_2\to0$.
The vanishing boundary trace and Poincar\'e's inequality then force
$\|X\|_2\to0$, contradicting \eqref{spec:case-II-normalization}.
\end{proof}

\subsection{The constrained compact limit}

The preceding estimates keep every possible mode in a bounded set of
angular and spectral parameters as $a\to0$. We now fix an angular degree
and analyze that bounded region. The pressure coefficient $P_a/a$ becomes
singular in the core; its limiting effect is the constraint $f=0$ there.
We first prove convergence of the inverses with this constraint, then
use an analytic compact pencil to transfer both the limiting roots and
their algebraic multiplicities to positive $a$.

\begin{lemma}[Convergence of the constrained inverses]\label{spec:inverse-limit}
Fix $\ell$ and a convergent profile subsequence. Let $R=2\sqrt Z$ be
its limiting core radius, and fix a sufficiently large outer radius
$B_2>r_b$. Let $E$ be the Hilbert space of $H^1$ pairs in this harmonic
on the ball $r<B_2$, and let
\[
 \mathscr Z=L^2(r<B_2;\mathbb C^2)\oplus\mathbb C^2,\qquad
 j_E X=(X,X(B_2)).
\]
The map $j_E:E\to\mathscr Z$ is compact. Define $T_a:E'\to E$ by the
coercive form
\begin{equation}\label{spec:base-form}
 \mathfrak b_a(X,\Theta)=
 \int\left(\nabla\Theta^*\nabla X+\Theta^*X
                         +(P/a)\overline{\Theta_f}f\right)\mu.
\end{equation}
The dual $E'$ consists of conjugate-linear functionals on test vectors.
Let
\[
 E_0=\{(f,g)\in E:f=0\text{ a.e. on }r<R\},
\]
and let $T_0:E'\to E_0$ be the inverse of the limiting form, with no
pressure term and tests restricted to $E_0$. Then
\begin{equation}\label{spec:inverse-norm}
 \|j_ET_a-j_ET_0\|_{E'\to\mathscr Z}\longrightarrow0.
\end{equation}
\end{lemma}

\begin{proof}
Rellich compactness gives the volume part of $j_E$; in a fixed harmonic
the trace has only two scalar entries and is bounded, hence compact.
Uniform upper and lower bounds on $\mu$ give uniform coercivity of
\eqref{spec:base-form}. Inverses exist by the Hilbert-space coercive
form argument, also on the closed subspace $E_0$.

We prove a stronger sequential statement. Let $F_a\rightharpoonup F$
be any bounded weakly convergent sequence in $E'$, and put
$X_a=T_aF_a$. Coercivity gives a bounded sequence in $E$ and
\[
 \int P_a|f_a|^2\,dy\le Ca.
\]
After extraction, $X_a\rightharpoonup X$ in $E$ and $j_EX_a\to j_EX$
strongly. Uniform boundedness of $P_a$ permits replacing $|f_a|^2$
by $|f|^2$ in the last integral, with error $o(1)$. Its weak-star
limit $b\mathbf1_{r<R}$ therefore gives
$b\int_{r<R}|f|^2=0$. This proves $X\in E_0$ without any pointwise
lower bound on pressure near its moving transition.

Test first with $\Theta\in E_0$ whose first component is supported a
fixed positive distance outside $r=R$. On that support the limiting
amplitude is strictly below one, so $P_a/a\to0$ uniformly. The other
coefficients converge uniformly on the ball. Weak convergence in
$E$, and weak convergence of the inputs against this fixed test,
therefore give the limiting inverse equation.
Such tests are dense in $E_0$. Indeed the first component has zero
outer trace at $R$; cutting it off across $[R+\epsilon,R+2\epsilon]$
costs, in addition to its shrinking collar norm, at most
\[
 C\epsilon^{-2}\int_R^{R+2\epsilon}|\Theta_f|^2\,dr
 \le C\int_R^{R+2\epsilon}|\partial_r\Theta_f|^2\,dr
 \longrightarrow0.
\]
The radial weight is comparable to a positive constant on this
collar. The second component needs no cutoff. Density and uniqueness
identify $X=T_0F$. Every subsequence has the same possible strong
limit after $j_E$, proving
$j_ET_aF_a\to j_ET_0F$ for weakly varying inputs.

If \eqref{spec:inverse-norm} failed, choose unit witnessing inputs
and extract a weakly convergent subsequence in the Hilbert dual.
The sequential assertion just proved and compactness of $j_ET_0$
would make both images converge strongly to $j_ET_0F$, a
contradiction. This proves exactly the stated norm convergence; no
operator-norm convergence into the full $H^1$ space is asserted.
\end{proof}

\begin{lemma}[The analytic compact pencil]\label{spec:compact-pencil}
Near each bounded candidate $\lambda$, choose $B_2$ sufficiently large.
The regular outgoing mode equation is equivalent to
\begin{equation}\label{spec:compact-equation}
 (I-\mathcal C_a(\lambda))z=0\quad\hbox{on }\mathscr Z,\qquad
 \mathcal C_a(\lambda)=j_ET_aK_a(\lambda).
\end{equation}
These are compact holomorphic families, and
$\mathcal C_a\to\mathcal C_0$ locally uniformly in operator norm.
\end{lemma}

\begin{proof}
Lemma~\ref{spec:slow-basis} and its large-radius normalization show
that the values of a slow basis at $B_2$ form an invertible matrix,
uniformly in a small neighborhood of the candidate parameter and
for sufficiently small $a$. Thus the outgoing condition is
\[
 X'(B_2)=\mathcal M_a(\lambda)X(B_2),
\]
with holomorphic matrices converging locally uniformly to
$\mathcal M_0$. Define $K_a$ on volume values and boundary entries
by its pairing with $\Theta\in E$:
\begin{equation}\label{spec:K-form}
 \begin{split}
 \langle K_a(\lambda)j_EX,\Theta\rangle
 ={}&\int\left\{\Theta^*(I+(c_*-\lambda)J)X
                         +(W\Theta)^*JX\right\}\mu\\
 &+[\Theta^*(\mathcal M_a(\lambda)-w(B_2)J)X].
 \end{split}
\end{equation}
The boundary bracket is at $B_2$. This is the weak equation obtained
by integrating $-\Theta^*JW X$ by parts. It uses only the two types
of data in $j_EX$. Its coefficients converge uniformly, and hence
$K_a:\mathscr Z\to E'$ converges locally uniformly in operator norm.
Lemma~\ref{spec:inverse-limit} proves the asserted convergence of
$\mathcal C_a$, and compactness follows by factoring through $E$.
For a kernel vector $z$ its lift is $X=T_aK_a z$, with $j_EX=z$.
For positive $a$, smooth elliptic regularity gives a regular solution
through the origin; its Robin condition extends it uniquely into the
slow space. Conversely any regular outgoing mode satisfies the weak
equation. These two maps are inverse on the kernels.
\end{proof}

\begin{lemma}[Interface conditions and the limiting determinant]
\label{spec:interface}
The limiting kernel in \eqref{spec:compact-equation} consists of free
slow exterior modes satisfying
\begin{equation}\label{spec:matching}
 f(R)=0,\qquad g(R+)=g(R-),\qquad g'(R+)=g'(R-),\qquad
 g(r)=c r^\ell\quad(r<R).
\end{equation}
It is nonzero exactly when the two-column determinant of
Proposition~\ref{free:main} vanishes.
\end{lemma}

\begin{proof}
The form domain gives $f=0$ in the core and continuity of both traces
at $R$. Testing the second equation without a constraint on its test
component gives
$Lg=(\lambda+W)f$ across the interface. Its right side is locally
$L^2$ and has no surface delta, since $f$ is $H^1$. Radial integration
therefore makes $\mu r^{11}g'$ continuous. Positivity and continuity
of $\mu$ give continuity of $g'$. Inside the core $\mu=1$ and $f=0$,
so regularity and harmonicity give $g=c r^\ell$.
Piecewise ODE regularity supplies the one-sided derivatives involved.
The first equation may contain a constraint reaction at $R$; the form
domain does not impose any condition on $f'(R+)$. In particular we
do not impose the generally false condition $f'(R+)=0$.

Since $Q(R)=1$ and $Q'(R)=0$, the exterior circular conditions are
\[
 v_++v_-=0,\qquad
 (\partial_r-\ell/R)(v_+-v_-)=0.
\]
Using \eqref{spec:free-bases}, the derivative of the argument
$-hi r^2/4$ is $-hi r/2$. Consequently these two equations, after
nonzero row and column scalings and a sign change in one unknown
coefficient, have columns
\begin{equation}\label{spec:free-determinant}
 \binom{H(\lambda+\ell/2-hib,6+\ell,-hiZ)}
       {-H'(\lambda+\ell/2-hib,6+\ell,-hiZ)},\qquad h=\pm1.
\end{equation}
This is exactly the determinant of Proposition~\ref{free:main}.
The columns are nonzero by Lemma~\ref{free:columns}.
\end{proof}

\begin{lemma}[Symmetry roots and first-order chains]\label{spec:simplicity}
At the limiting disk root $j(b,Z)=0$, the only roots of
\eqref{spec:free-determinant} in $\operatorname{Re}\lambda\ge-1/32$
are $0,1$ for $\ell=0$ and $1/2$ for $\ell=1$, each simple.
The compact pencil $I-\mathcal C_0$ has a one-dimensional kernel at
each such root and no first-order analytic chain.
\end{lemma}

\begin{proof}
For every positive $a$ the first circular components
$iQ$, $\partial_jQ$, and $(a-ib+D/2)Q$ solve the linearized equation
at $0,1/2,1$, respectively, by differentiating phase, center, and
blow-up time. These identities can also be checked without referring
to a family of solutions. The first-component generator is
\[
 \mathcal A v=i\Delta v-Dv/2-av+ibv-iD_Q(|Q|^{1/a}Q)v.
\]
Phase covariance gives $\mathcal A(iQ)=0$; differentiating the
stationary equation gives $\mathcal A(\partial_jQ)=\partial_jQ/2$.
Commuting $D$ with $\Delta$ and using homogeneity gives
$\mathcal A(DQ/2)=i\Delta Q$ and
$\mathcal A(aQ)=-i|Q|^{1/a}Q$, whose sum equals
$(a-ib+D/2)Q$ by the stationary equation. The term $-ibQ$ is a phase
multiple and contributes zero.

At $a=0$ the same exterior identities hold with the power term
deleted. The limiting phase, time, and translation fields satisfy
\eqref{spec:matching}. For the time mode its first circular component
at $R$ is $-ib$, while its derivative is $-Rb/2$, which is real;
for a translation it has value zero and derivative $-b$ times the
appropriate first harmonic. These facts follow from the exterior
profile equation $Q''(R)=-b$. The phase mode is immediate.
They are slow by the free expansions; in the time mode the leading
tail exponent cancels. They are nonzero: $b>0$, and the exterior
profile is nonconstant. Proposition~\ref{free:main} gives exactly
two, one, and zero roots with multiplicity in degrees $0,1,$ and
$\ell\ge2$. Thus the displayed distinct roots exhaust the counts
and are simple.

Suppose the compact pencil has a first-order chain, that is,
$z_0\ne0$ and
\[
 (I-\mathcal C_0(\lambda_0))z_0=0,\qquad
 (I-\mathcal C_0(\lambda_0))z_1
       -\mathcal C_0'(\lambda_0)z_0=0.
\]
Lift it by $T_0K_0$, including the derivative term in the lift of
$z_1$. Both first components vanish in the core. Differentiating the
second equation shows that both $g$ components are harmonic there,
since its new forcing is the leading $f$, which vanishes in the core.
Both obey the interface trace and derivative requirements. At $B_2$
the correction satisfies the differentiated Robin condition
\[
 X_1'=\mathcal M_0(\lambda_0)X_1
                 +\mathcal M_0'(\lambda_0)X_0.
\]
The analytic slow basis extends these data to a first-order exterior
solution. Therefore its two coefficients satisfy
$D_0c_0=0$ and $D_0c_1+D_0'c_0=0$, where $D_0(\lambda)$ is the
matrix in \eqref{spec:free-determinant}. A simple determinant zero
has rank one and nonzero derivative pairing between a null vector
and a conull vector. Applying that conull vector to the second
equation is a contradiction. Kernel dimension one follows from the
same matrix and the unique harmonic continuation into the core.
\end{proof}

\begin{proof}[Proof of Theorem~\ref{spectrum:classification}]
First work in a fixed harmonic and a compact spectral neighborhood.
The compact pencil is Fredholm of index zero. At a limiting simple
root split its domain into the kernel and a closed complement, and
its range into the image and a one-dimensional complement. The
complement block is invertible, so elimination reduces the local
kernel problem to a scalar holomorphic function. Its derivative is
nonzero: vanishing would produce exactly the first-order chain
excluded by Lemma~\ref{spec:simplicity}. Locally uniform norm
convergence in Lemma~\ref{spec:compact-pencil} gives locally uniform
convergence of this scalar function and its derivative, by Cauchy's
formula. The argument principle on a small circle therefore gives
one zero counted with multiplicity for small $a$. The exact symmetry
mode places that zero at $0,1/2,$ or $1$, so it remains simple. Away
from the limiting roots invertibility persists by a Neumann series.

If the classification failed for arbitrarily small $a$, select
offending modes with $a\to0$. Lemma~\ref{spec:real-bound} bounds
their real parts, and Lemma~\ref{spec:high-exclusion} bounds their
angular and imaginary parameters. Extract a fixed harmonic, a
convergent spectral parameter, and convergent profile parameters.
The preceding compact-pencil argument gives a contradiction.
This also gives the assertion uniformly over all the profiles in
Proposition~\ref{profile:existence}. Lemma~\ref{spec:slow-basis}
then proves the homogeneous top-derivative assertion.

Finally suppose a regular vector $v_1$ with the top-derivative
condition satisfies $(\mathcal A-\lambda_0)v_1=v_0$, where $v_0$ is
one of the surviving modes. Choose an analytic exterior slow family
$v(\lambda)$ with $v(\lambda_0)=v_0$. Differentiating its mode
equation gives
$(\mathcal A-\lambda_0)\partial_\lambda v=v_0$ on the exterior.
Its differentiated slow tail has bounds
\[
 |\partial_r^k\partial_\lambda v|
 \le C r^{-2(a+\operatorname{Re}\lambda_0)-k}(1+\log r).
\]
For example, Cauchy's formula on parameter circles of radius
$(1+\log r)^{-1}$ gives these bounds from the locally uniform
analytic slow construction. This tail is square integrable in
order $k\ge7$ at the three nonnegative eigenvalues. Subtracting it
from $v_1$ leaves a homogeneous exterior solution with the same
top-derivative condition, so Lemma~\ref{spec:slow-basis} makes that
difference outgoing. The original $v_1$ consequently satisfies the
differentiated Robin condition and produces a forbidden first-order
chain of the finite-$a$ compact pencil. This proves the claimed
algebraic simplicity and completes the proof.
\end{proof}

\section{Linear decay and transfer to expanding tori}\label{sec:linear}

Fix one sufficiently small admissible $a>0$ and a profile $Q=Ae^{i\phi}$
from Proposition~\ref{profile:existence} for which
Theorem~\ref{spectrum:classification} holds.  All constants in this section
may depend on this fixed profile and on $a$.  No uniformity as $a\downarrow0$
is needed here.  Put $P=|Q|^{1/a}$ and write the real-linear potential as
\begin{equation}\label{lin:potential}
 \mathcal Vv=P\left(1+\frac1{2a}\right)v
       +\frac{P}{2a}e^{2i\phi}\overline v,
 \qquad
 \mathcal A=i\Delta-\frac12y\cdot\nabla-a+ib-i\mathcal V.
\end{equation}
The profile estimates give, for every nonnegative integer $j$,
\begin{equation}\label{lin:symbol-hypotheses}
 |\nabla^j Q(y)|\le C_j\langle y\rangle^{-2a-j},\qquad
 |\nabla^j\mathcal V(y)|\le C_j\langle y\rangle^{-2-j}.
\end{equation}
The second inequality applies to each of the two coefficients in
\eqref{lin:potential}.  We take $0<a<1$.  The argument below uses the
all-order estimates only up to a finite order, after that order has been
chosen.

\subsection{Spaces, products, and localization}

Set $s_0=6-a$.  For an integer $k>8$ to be fixed below, let $Y$ be the real
Hilbert space obtained by completing complex-valued Schwartz functions in
\begin{equation}\label{lin:Y}
 \|v\|_Y^2=\|v\|_{\dot H^{s_0}(\mathbb R^{12})}^2
                      +\|v\|_{\dot H^k(\mathbb R^{12})}^2.
\end{equation}
Homogeneous Fourier multipliers have their usual Plancherel normalization.
In particular, no global $L^2$ norm is included in \eqref{lin:Y}.
The choice $6-2a<s_0=6-a<6$ serves two purposes: the lower inequality
admits the profile tail of order $r^{-2a}$, while the upper inequality
makes the inverse Fourier weight integrable at low frequency. Under the
free similarity evolution this norm component will have squared-norm
rate $6-s_0-2a=-a<0$. The high order $k$ supplies the local regularity
and the remaining dissipation.

For $L\ge1$ put $\mathbb T_L^{12}=\mathbb R^{12}/(2\pi L\mathbb Z)^{12}$.
If $v(y)=\sum_{n\in\mathbb Z^{12}}v_ne^{in\cdot y/L}$, define
\begin{equation}\label{lin:YL}
 \|v\|_{Y_L}^2=(2\pi L)^{12}\sum_n
 \left[(L^{-2}+|n/L|^2)^{s_0}+|n/L|^{2k}\right]|v_n|^2.
\end{equation}
Thus $Y_L$ is $H^k(\mathbb T_L^{12})$ with a scale-dependent norm.
The regularization in its first summand is essential: a constant of value
$z$ has norm $(2\pi)^6L^a|z|$.
Throughout this section $\nabla^k$ is the \emph{ordered} derivative tensor:
its components are $\partial_{j_1}\cdots\partial_{j_k}$ for all ordered
$k$-tuples.  Consequently
$\|\nabla^kv\|_2^2=\|v\|_{\dot H^k}^2$ exactly.

\begin{lemma}[Uniform spatial estimates]\label{lin:spatial}
The embeddings $Y\hookrightarrow C_0(\mathbb R^{12})$ and
$Y_L\hookrightarrow L^\infty(\mathbb T_L^{12})$ are bounded, uniformly for
$L\ge1$.  Both spaces are algebras, with a uniform algebra constant for
$Y_L$.  For a ball or annulus of radius $R\ge1$, contained in one torus
cell when appropriate and with $R\le C_0L$, one has
\begin{equation}\label{lin:local-derivatives}
 \|\nabla^jv\|_{L^2(\text{ball or annulus})}
 \le C R^{\max(s_0-j,0)}\|v\|,
 \qquad 0\le j\le k.
\end{equation}
Here and below an undecorated norm denotes the applicable $Y$ or $Y_L$
norm, and constants may depend on the fixed geometric constant $C_0$.

If $F$ satisfies $|\nabla^jF|\le C_j\langle y\rangle^{-\beta-j}$, then
a smooth dyadic piece supported where $|y|\asymp R$ satisfies
\begin{equation}\label{lin:symbol-norm}
 \|\psi_RF\|\le C\big(R^{6-\beta-s_0}+R^{6-\beta-k}\big),
\end{equation}
also on the tori for $R\le C_0L$.  In particular $Q\in Y$, and for every
fixed $\chi\in C_c^\infty(B_1)$ the cutoff profiles $\chi(y/L)Q(y)$
have uniformly bounded $Y_L$ norms.
\end{lemma}

\begin{proof}
For the continuous Fourier transform, Cauchy--Schwarz uses
\[
 \int_{|\xi|\le1}|\xi|^{-2s_0}\,d\xi<\infty,
 \qquad
 \int_{|\xi|\ge1}|\xi|^{-2k}\,d\xi<\infty.
\]
Thus $\|\widehat v\|_1\le C\|v\|_Y$.  This identifies the completion in
\eqref{lin:Y} with functions vanishing at infinity, rather than with
distributions modulo polynomials.  In the periodic case the corresponding
inverse-weight sum, after division by $(2\pi L)^{12}$, is uniformly
bounded.  Indeed its portion $|n|\le L$ is bounded by
\[
 C L^{-12+2s_0}\sum_{|n|\le L}\langle n\rangle^{-2s_0}\le C,
\]
and the portion $|n|>L$ is bounded using $k>6$.
The convolution inequality
$|\xi+\zeta|^s\le C_s(|\xi|^s+|\zeta|^s)$, and its version for
$(L^{-2}+|\xi|^2)^{s/2}$, now prove the algebra estimates by Young's
inequality.  The volume factors in \eqref{lin:YL} cancel against the
Fourier $\ell^1$ estimate just proved.

For $j<s_0$, split at frequency $R^{-1}$.  The low-frequency part has
$j$th derivative bounded pointwise by $CR^{-a-j}\|v\|$.  This follows
by the same inverse-weight integral or sum; it includes the periodic
zero mode.  Multiplication by the square root of the volume of the
ball gives $CR^{s_0-j}\|v\|$.  On higher frequencies, Plancherel gives
that same bound because
$|\xi|^j\le R^{s_0-j}|\xi|^{s_0}$ there.  For $s_0\le j\le k$,
the global derivative norm is bounded by the two defining Sobolev
norms.  This proves \eqref{lin:local-derivatives}.

For a localized symbol, the integer derivative norms of order $m$ are
$O(R^{6-\beta-m})$, with the same bounds for derivatives falling on
$\psi_R$.  Fourier interpolation gives the fractional order in
\eqref{lin:symbol-norm}.  The additional periodic term is bounded by
\[
 L^{-s_0}\|\psi_RF\|_2
 \le C(R/L)^{s_0}R^{6-\beta-s_0},
\]
since the first norm in \eqref{lin:YL} is uniformly equivalent to
$\|\cdot\|_{\dot H^{s_0}}+L^{-s_0}\|\cdot\|_2$.
For $\beta=2a$ the first exponent in \eqref{lin:symbol-norm} is $-a$;
both exponents are negative.  Summing dyadic pieces proves the last
assertions, and also convergence of smooth compactly supported
approximations to $Q$ in $Y$.
\end{proof}

Fix a real $\chi\in C_c^\infty(B_1)$ with $0\le\chi\le1$ and
$\chi=1$ on $B_{1/2}$.  Define
\begin{equation}\label{lin:localization}
 J_Lv(y)=\chi(y/L)v(y),
\end{equation}
using the central-cell representative on its support and extending by
zero to $\mathbb R^{12}$.  There is a constant $C_J$, independent of
$L\ge1$, such that
\begin{equation}\label{lin:J-bound}
 \|J_Lv\|_Y\le C_J\|v\|_{Y_L}.
\end{equation}
Here is a direct verification, useful because a fractional derivative
does not localize pointwise.  Up to fixed Fourier normalization constants,
\[
 \widehat{J_Lv}(\xi)=L^{12}\sum_nv_n\widehat\chi(L\xi-n).
\]
For $s=s_0,k$, set $\zeta=L\xi$ and use
$|\zeta|^s\le C_s\langle n\rangle^s\langle\zeta-n\rangle^s$.
The mixed discrete-to-continuous convolution map with kernel
$\langle\zeta\rangle^s|\widehat\chi(\zeta)|$ is bounded from
$\ell^2$ to $L^2$: Cauchy--Schwarz uses the uniformly bounded sum of
its translates, followed by integration of the kernel.  It follows that
\[
 \|J_Lv\|_{\dot H^s}^2
 \le C L^{12-2s}\sum_n\langle n\rangle^{2s}|v_n|^2.
\]
For $s=k$ the extra term $L^{-2k}\|v\|_2^2$ is controlled by the
regularized low-order norm since $L^{-k}\le L^{-s_0}$.  This proves
\eqref{lin:J-bound}.

\subsection{Evolution and dissipation with a local observation}

The free operator $\mathcal A_0=\mathcal A+i\mathcal V$ has the group
\begin{equation}\label{lin:free-flow}
 S_0(s)v(y)=e^{(-a+ib)s}
 \left[e^{i(1-e^{-s})\Delta}v\right](e^{-s/2}y),\qquad s\in\mathbb R.
\end{equation}
Fourier approximation proves strong continuity on $Y$.  The potential
is bounded on $Y$: since $p=1+1/a$ is an odd integer, its coefficients
are polynomials in $Q,\overline Q$, and Lemma~\ref{lin:spatial} applies.
The integral equation
\begin{equation}\label{lin:mild}
 S(s)v=S_0(s)v-i\int_0^s S_0(s-t)\mathcal V S(t)v\,dt
\end{equation}
therefore constructs a strongly continuous real-linear group on $Y$.
Successive approximation on short intervals proves existence, and
Gronwall's inequality proves uniqueness and boundedness on every compact
time interval.

For the torus evolution, put $L_s=Le^{s/2}$ and
\begin{equation}\label{lin:torus-coefficients}
 Q_s(y)=\chi(y/L_s)Q(y),\qquad
 \mathcal V_s=D\big(|U|^{1/a}U\big)\big|_{U=Q_s}.
\end{equation}
Both are defined cellwise and periodically.  In particular
$\mathcal V_s=\chi(y/L_s)^{p-1}\mathcal V$ on the central cell.
The equation is
\begin{equation}\label{lin:torus-evolution}
 (\partial_s+\tfrac12y\cdot\nabla)v
       =i\Delta v-a v+ibv-i\mathcal V_s v.
\end{equation}
It is interpreted on the expanding torus: the left side differentiates
at fixed $y/L_s$.  Formula~\eqref{lin:free-flow} and a nonautonomous
version of \eqref{lin:mild} apply.  The uniform product bounds give a
unique propagator with bounds depending on a fixed slab length, but
independent of the starting scale $L\ge1$.

Let $B$ bound the pointwise real operator norms of $\mathcal V$ and all
$\mathcal V_s$.  It is finite and independent of $k$; one may take
$B=(1+1/a)\|P\|_\infty$.  We now choose and fix an integer
\begin{equation}\label{lin:k-choice}
 k>\max\{8,\,6-2a+2B+1\}.
\end{equation}
Subsequent constants may depend on this $k$.

\begin{lemma}[Compact errors]\label{lin:compact-error}
For every $\varepsilon>0$ there are $R_\varepsilon,C_\varepsilon<\infty$
such that
\begin{equation}\label{lin:compact-bound}
 \|\mathcal V_*v\|_{s_0,*}
 +\|\nabla^k(\mathcal V_*v)-\mathcal V_*\nabla^kv\|_2
 \le\varepsilon\|v\|+C_\varepsilon\|v\|_{L^2(B_{R_\varepsilon})}.
\end{equation}
Here $*$ denotes either the whole-space potential and homogeneous
low-order norm, or a torus potential at its current scale and the
regularized low-order norm in \eqref{lin:YL}.  The constants are uniform
for all sufficiently large current torus scales.  The potential acts
componentwise on the ordered tensor in the second term.
\end{lemma}

\begin{proof}
All coefficient derivatives have the bounds in
\eqref{lin:symbol-hypotheses}, uniformly for the cutoff potentials.
On an annulus of radius $R$, a lower Leibniz term with $j<k$ derivatives
on $v$ has norm at most
\begin{equation}\label{lin:top-tail}
 C_kR^{-2-k+j+\max(s_0-j,0)}\|v\|.
\end{equation}
Every exponent is negative: it is $-2-k+s_0$ for $j<s_0$ and at most
$-3$ otherwise.  Dyadic summation therefore makes these tail terms
arbitrarily small.

For the fractional part, use a smooth dyadic partition and apply
\emph{global} Fourier interpolation to the actual localized product
$z_R=\psi_R\mathcal V_*v$.  Leibniz' rule, including derivatives of
$\psi_R$, gives
\[
 \|\nabla^5z_R\|_2\le CR^{-2+s_0-5}\|v\|
                     =CR^{-1-a}\|v\|,
 \qquad
 \|\nabla^6z_R\|_2\le CR^{-2}\|v\|.
\]
In the second estimate, terms with at most five derivatives on $v$ are
actually $O(R^{-2-a})$; the sixth derivative term is $O(R^{-2})$.
Since $s_0=5a+6(1-a)$, interpolation yields the explicit bound
\begin{equation}\label{lin:fractional-tail}
 \|z_R\|_{\dot H^{s_0}}
 \le\|\nabla^5z_R\|_2^a\|\nabla^6z_R\|_2^{1-a}
 \le CR^{-2+a-a^2}\|v\|.
\end{equation}
The identical proof uses Fourier sums on a torus.  Its remaining
regularized term is bounded separately by
\begin{equation}\label{lin:zero-mode-tail}
 L_s^{-s_0}\|z_R\|_2
 \le CR^{-2}(R/L_s)^{s_0}\|v\|
 \le CR^{-2}\|v\|,
\end{equation}
with a harmless geometric constant when $R\le C_0L_s$.
Both exponents are summable over dyadic annuli; indeed
$-2+a-a^2\le-7/4$ for $0<a<1$.  Summation uses the Sobolev triangle
inequality, not a locality assertion for fractional derivatives.

It remains to estimate a fixed interior region.  Its low-order product
is controlled by a local $H^6$ norm, and the lower top-order Leibniz
terms by a local $H^{k-1}$ norm.  Choose a slightly larger fixed ball
and a cutoff equal to one on the region.  The usual interpolation
inequality on these two balls gives, for each $m<k$,
\[
 \|v\|_{H^m(B_R)}
 \le\delta\|v\|_{H^k(B_{2R})}
                 +C_{R,m,\delta}\|v\|_{L^2(B_{2R})}.
\]
For completeness, this follows from compactness of the embedding
$H^k(B_{2R})\to H^m(B_R)$: a failure after normalizing the left side
would give a bounded sequence whose $L^2$ norm tends to zero but whose
$H^m$ norm does not.  Local extension and Fourier truncation give this
compactness.  By \eqref{lin:local-derivatives}, the $H^k$ norm on the
fixed larger ball is bounded by $C_R\|v\|$.  Choose the tail radius
first and $\delta$ next.  This proves \eqref{lin:compact-bound}, including
the interior contribution of the regularized torus term.
\end{proof}

\begin{proposition}[Energy observation inequality]\label{lin:observation}
There are $c>0$, $C<\infty$, and $R_0<\infty$ such that every whole-space
linear solution, and every torus linear solution with sufficiently large
starting scale, satisfies
\begin{equation}\label{lin:energy-observation}
 \frac{d}{ds}\|v(s)\|^2
 \le-c\|v(s)\|^2+C\|v(s)\|_{L^2(B_{R_0})}^2.
\end{equation}
The derivative can equivalently be understood through the integrated
inequality.  The constants are independent of the torus scale and of
the slab length.  In particular, for $t\ge0$,
\begin{equation}\label{lin:integrated-observation}
 \|v(t)\|^2\le e^{-ct}\|v(0)\|^2
   +C\int_0^t e^{-c(t-s)}\|v(s)\|_{L^2(B_{R_0})}^2\,ds.
\end{equation}
\end{proposition}

\begin{proof}
The two squared free norm rates, on both kinds of space, are
\begin{equation}\label{lin:free-rates}
 6-s_0-2a=-a,\qquad 6-k-2a.
\end{equation}
For the first periodic norm this also uses
$L_s^{-2}=e^{-s}L^{-2}$; thus the zero mode has exactly the same first
rate.  At order $k$, the leading potential term contributes at most
$2B\|\nabla^kv\|_2^2$.  This constant does not grow with $k$, because
the same pointwise operator acts on each ordered derivative component.
All combinatorial constants occur only in the lower terms covered by
Lemma~\ref{lin:compact-error}.  The entire fractional potential term is
also covered there.  Consequently, writing $N_0,N_k$ for the two norm
components, one obtains
\[
 \frac{d}{ds}(N_0^2+N_k^2)
 \le-aN_0^2-(k+2a-6-2B)N_k^2
       +C\|v\|\big(\varepsilon\|v\|
                         +C_\varepsilon\|v\|_{L^2(B_{R_\varepsilon})}\big).
\]
Choose $\varepsilon$ after the fixed $k$, and apply Young's inequality.
Condition~\eqref{lin:k-choice} leaves a positive $c$.

This computation does not demand two additional derivatives of the
initial datum.  On a short step pull the mild solution back by its free
propagator.  Strong continuity and the bounded potential give first
variation $-i\mathcal V_*v$ in the relevant space, while the two free
squared norms acquire exactly the exponential factors in
\eqref{lin:free-rates}.  Taking the norm variation gives the displayed
inequality for arbitrary mild solutions.  Equivalently one can first
use smooth data, integrate the inequality, and pass to general data by
the uniform mild-solution bounds.  This proves
\eqref{lin:integrated-observation} as well.
\end{proof}

\subsection{Actual decay on the whole-space stable subspace}

Let $E:\mathbb R^{14}\to Y$ have columns
\begin{equation}\label{lin:E}
 iQ,\quad \partial_1Q,\ldots,\partial_{12}Q,\quad
 (a-ib+\tfrac12y\cdot\nabla)Q,
 \qquad
 \Lambda=\operatorname{diag}(0,\tfrac12,\ldots,\tfrac12,1).
\end{equation}
These columns are independent symmetry vectors, as established in
Theorem~\ref{spectrum:classification}.

The passage from dissipation modulo a compact term to stable semigroup
decay also appears in \cite[Section~4.2]{BuckEtAl2026}. Here the full
mode classification identifies the complementary spectral space with
the fourteen symmetry directions above.

\begin{theorem}[Stable whole-space evolution]\label{linear:decay}
There is a bounded real-linear map $\pi:Y\to\mathbb R^{14}$ such that
\begin{equation}\label{lin:spectral-coordinates}
 \pi E=I,\qquad S(s)E=Ee^{s\Lambda},\qquad
 \pi S(s)=e^{s\Lambda}\pi\quad(s\in\mathbb R).
\end{equation}
For some $C_S<\infty$ and $\eta>0$,
\begin{equation}\label{lin:stable-decay}
 \|S(s)v\|_Y\le C_Se^{-\eta s}\|v\|_Y,
 \qquad s\ge0,\quad \pi v=0.
\end{equation}
Thus $E\pi$ is the symmetry spectral projection.  Constants depend on
the fixed profile, $a$, and $k$.
\end{theorem}

\begin{proof}
We first obtain a statement about the bounded time-one operator
$T=S(1)$, rather than deducing decay from a generator spectral gap.
If $v_j\rightharpoonup0$ in $Y$ and $\|v_j\|\le1$, then for every fixed
$s$ the functions $S(s)v_j$ converge strongly to zero in
$L^2(B_{R_0})$.  Indeed they converge weakly in $Y$, and the restriction
from $Y$ to that $L^2$ space is compact by
\eqref{lin:local-derivatives} and local Sobolev compactness.  Slab bounds
and dominated convergence in \eqref{lin:integrated-observation} give
\begin{equation}\label{lin:essential-contraction}
 \limsup_j\|Tv_j\|\le e^{-c/2}.
\end{equation}
The same conclusion holds on the complexification, by adding the two
real energy inequalities.

Let $F_N$ be increasing finite-dimensional orthogonal projections
converging strongly to the identity.  Formula~\eqref{lin:essential-contraction}
implies $\limsup_N\|T(I-F_N)\|\le e^{-c/2}$: otherwise unit vectors
in the complements of $F_N$ would give a weakly null counterexample.
For some $N$ we therefore have
\[
 T=B_0+K_0,\qquad \|B_0\|=q<1,\qquad \operatorname{rank}K_0<\infty.
\]
Outside $|z|\le q$, invert $z-B_0$ and write
$K_0=UV$ through a finite-dimensional space.  The resolvent of $T$
is determined by the finite matrix
$I-V(z-B_0)^{-1}U$.  Its determinant is analytic and tends to one as
$|z|\to\infty$.  Thus the spectrum there consists of isolated
eigenvalues with finite-dimensional generalized eigenspaces; there
are only finitely many outside any larger disk.  This also proves the
usual spectral decomposition in this region by the finite-matrix
resolvent formula.

Choose $r$ with
\[
 \max\{q,e^{-1/32}\}<r<1
\]
and with $|z|=r$ disjoint from the spectrum.  Every spectral subspace
of $T$ outside this circle is preserved by $S(s)$, because $S(s)$
commutes with the time-one resolvent and its contour integrals.
The restricted continuous group on a finite-dimensional space is a
matrix exponential.  Each of its generator eigenvalues $\lambda$
satisfies $\operatorname{Re}\lambda=\log|z|>\log r>-1/32$ for the corresponding
time-one eigenvalue $z$.

Its generator equations are the distributional equations for
\eqref{lin:potential}, as follows by differentiating
\eqref{lin:mild} against compactly supported smooth tests.  Local
elliptic regularity gives smoothness of eigenvectors and chain
vectors.  Projection onto any spherical harmonic gives regular
radial solutions of the mode equation in
Theorem~\ref{spectrum:classification}.  The unweighted radial $k$th
derivative is in $L^2(r^{11}dr)$: along each ray $\partial_r^k$ is a
contraction of the ordered tensor $\nabla^k$, and angular projection
is bounded on $L^2$ of each sphere.  Smoothness of the projection at
the origin follows by projecting the Taylor polynomials, using the
harmonic decomposition of homogeneous polynomials.  The
top-derivative criterion in Theorem~\ref{spectrum:classification}
excludes the fast exterior branches, including in a generalized
eigenvector after subtraction of the differentiated slow basis.
That theorem therefore leaves precisely the symmetry vectors
\eqref{lin:E} and excludes all nontrivial Jordan chains.  Completeness
of spherical harmonics gives the statement for the original vectors;
a nontrivial chain would have a nonzero harmonic projection of its
leading vector and hence would contradict the same classification.

Conversely all columns of $E$ lie in $Y$ by
\eqref{lin:symbol-norm}; the time column even has symbol order
$-2a-2$, since its leading slow power cancels.  They are generator
domain vectors with the asserted eigenvalues.  To check the domain
claim directly, multiply a column by a smooth cutoff at radius $R$.
Both the discarded symbol tail and the commutator
\[
 i[\Delta,\chi(y/R)]E-\tfrac12(y\cdot\nabla\chi(y/R))E
\]
tend to zero in $Y$ by the same dyadic estimates.  The smooth compactly
supported approximations thus converge in the generator graph norm.
It follows that the entire spectral subspace outside $|z|=r$ is
$\operatorname{ran}E$, with the action in
\eqref{lin:spectral-coordinates}.  Complex conjugation in the
complexification preserves its projection, so this projection is the
complexification of a bounded real projection $E\pi$.

On its complement, the spectrum of $T$ lies in $|z|\le r$.  Choose
$r<r_1<1$.  The contour formula for powers on $|z|=r_1$ gives
\[
 \|T^n|_{\ker\pi}\|
 \le r_1^{n+1}\max_{|z|=r_1}
                  \|(z-T|_{\ker\pi})^{-1}\|.
\]
Bounded propagation for the remaining time in $[0,1]$ proves
\eqref{lin:stable-decay}, with $\eta=-\log r_1>0$ and an enlarged
$C_S$.  This derivation supplies actual semigroup decay, not only a
description of the generator spectrum.
\end{proof}

\subsection{A transfer proposition for expanding tori}

Define the cellwise cutoff symmetry map $E_L=\chi(y/L)E$.
The bounds already proved imply
\begin{equation}\label{lin:EL-convergence}
 \sup_{L\ge1}\|E_L\|<\infty,\qquad J_LE_L\longrightarrow E
       \quad\hbox{in }\mathcal L(\mathbb R^{14},Y).
\end{equation}
Indeed $J_LE_L=\chi(y/L)^2E$, and its omitted dyadic tail tends to
zero.  For all sufficiently large $L$ put
\begin{equation}\label{lin:piL}
 C_L=\pi J_LE_L,\qquad
 \pi_L=C_L^{-1}\pi J_L,\qquad P_L=E_L\pi_L.
\end{equation}
Then $C_L\to I$, $\pi_LE_L=I$, and $P_L$ is a projection.  The
operator norms of $E_L,\pi_L,P_L$, and $I-P_L$ are uniformly bounded.
Let $\mathcal F_L^{(M)}:Y_L\to Y_{Le^{M/2}}$ denote evolution by
\eqref{lin:torus-evolution} from time zero to time $M$.

\begin{proposition}[Uniform expanding-torus transfer]\label{linear:torus}
For each fixed $M>0$ the step maps $\mathcal F_L^{(M)}$ are uniformly
bounded for all sufficiently large $L$.  Writing $L'=Le^{M/2}$,
\begin{equation}\label{lin:coordinate-defect}
 \lim_{L\to\infty}
 \|\pi_{L'}\mathcal F_L^{(M)}-e^{M\Lambda}\pi_L\|
 _{\mathcal L(Y_L,\mathbb R^{14})}=0.
\end{equation}
There exists $M_0<\infty$ such that for every fixed $M\ge M_0$
there is $L_0(M)<\infty$ for which
\begin{equation}\label{lin:torus-stable-contraction}
 \|(I-P_{L'})\mathcal F_L^{(M)}|_{\ker\pi_L}\|
       _{\mathcal L(Y_L,Y_{L'})}\le\frac18,
 \qquad L\ge L_0(M).
\end{equation}
Moreover $\|(I-P_{L'})\mathcal F_L^{(M)}E_L\|\le C_M$.
Any prescribed positive bound on the coordinate defect can be imposed
by enlarging $L_0(M)$ after $M$ has been fixed.  These bounds persist
after restarting a step at any larger scale.

After identifying torus spaces by dilation, the maps above and the
step propagator are continuous when the scale and the input vary
continuously.  This last assertion is strong continuity for varying
inputs, and does not assert operator-norm continuity of Schr\"odinger
time evolution.
\end{proposition}

\begin{proof}
We give the compactness and uniqueness argument in detail.  Fix $M$
and a sequence $L_j\to\infty$ with initial data $v_j(0)$ bounded in
$Y_{L_j}$.  Write $L_{j,s}=L_je^{s/2}$.  Uniform slab bounds and
\eqref{lin:local-derivatives} give boundedness in local $H^k$, uniformly
for $0\le s\le M$.  On each fixed ball the equation also bounds
$\partial_sv_j$ in $H^{k-2}$; the drift coefficient is bounded on
that ball.  Local compactness, interpolation, and Arzel\`a--Ascoli
therefore give a subsequence converging in
$C([0,M];L^2(B_R))$ for every fixed $R$, to a function $v$.
One may obtain the stronger convergence in local $H^{k-1}$, but it
is unnecessary here.

The localized functions $J_{L_{j,s}}v_j(s)$ are bounded in $Y$ by
\eqref{lin:J-bound}.  At each $s$, every weakly convergent subsequence
has the same distributional limit $v(s)$, since $\chi(y/L_{j,s})=1$
on each fixed ball for all large $j$.  Thus
\begin{equation}\label{lin:weak-local-limit}
 J_{L_{j,s}}v_j(s)\rightharpoonup v(s)\quad\hbox{in }Y
 \quad(0\le s\le M),\qquad
 \|v(0)\|_Y\le C_J\liminf_j\|v_j(0)\|_{Y_{L_j}}.
\end{equation}
The limit is weakly continuous in $Y$: boundedness and its local
continuity identify the weak limit along every converging time
sequence.  As $Y$ is separable, it is strongly measurable and bounded
on this slab.  The torus coefficients agree with the whole-space
coefficients on each fixed ball eventually, so local passage in the
equation gives $\partial_sv=\mathcal Av$ in distributions.

This distributional solution is the whole-space mild solution.
To see this without a propagation assertion, test against Schwartz
functions transported by the free adjoint.  Formula~\eqref{lin:free-flow}
preserves Schwartz functions, uniformly on compact time intervals.
The global $L^\infty$ bound supplied by $Y$, and the at-most-linear
growth of the drift, justify removal of spatial cutoffs in these
tests.  Integration in time gives
\[
 S_0(-s)v(s)=v(0)-i\int_0^s S_0(-t)\mathcal Vv(t)\,dt
\]
as tempered distributions.  The integral is a Bochner integral in
$Y$ by boundedness of $\mathcal V$ and strong measurability.  The
distributional identity is therefore an identity in $Y$ and is
equivalent to \eqref{lin:mild}.  Uniqueness gives
\begin{equation}\label{lin:limit-identification}
 v(s)=S(s)v(0).
\end{equation}

Finite dimensionality, \eqref{lin:weak-local-limit}, and $C_L\to I$
now imply convergence of the initial and final coordinates to
$\pi v(0)$ and $\pi S(M)v(0)$, respectively.  Their difference in
\eqref{lin:coordinate-defect} tends to zero by
\eqref{lin:spectral-coordinates}.  This proves convergence in
operator norm: failure would furnish a sequence of unit initial
vectors witnessing a fixed positive lower bound, to which exactly
the preceding subsequence argument applies.

For the contraction, take such a sequence with
$\|v_j(0)\|_{Y_{L_j}}=1$ and $\pi_{L_j}v_j(0)=0$.  The limit satisfies
$\pi v(0)=0$, so \eqref{lin:stable-decay} gives
\[
 \|v(s)\|_Y\le C_SC_Je^{-\eta s}.
\]
The constants in this estimate are independent of the fixed slab
length $M$.  Apply the integrated observation inequality
\eqref{lin:integrated-observation} \emph{before} passing to a local
limit.  The full initial norm is still one in its first term.  Local
uniform convergence passes only the observation integral to the
limit, giving
\begin{equation}\label{lin:escaped-norm}
 \limsup_j\|\mathcal F_{L_j}^{(M)}v_j(0)\|_{Y_{L_je^{M/2}}}^2
 \le D(M):=e^{-cM}+C_1\int_0^M
                         e^{-c(M-s)}e^{-2\eta s}\,ds.
\end{equation}
Here $C_1$ uses only the fixed observation embedding, $C_J$, $C_S$,
and the energy constant; it does not contain the slab bound $C_M$.
Any norm that escaped to spatial infinity remains charged in the
first term $e^{-cM}$.  In particular a zero local limit causes no loss
in this estimate.

The quantity $D(M)$ tends to zero as $M\to\infty$; when $c=2\eta$
its integral is $Me^{-cM}$, and otherwise it is
$(e^{-2\eta M}-e^{-cM})/(c-2\eta)$.  Let $C_P$ uniformly bound
$\|I-P_L\|$ for sufficiently large $L$.  Choose $M_0$ so that
$C_P\sqrt{D(M)}<1/16$ for every $M\ge M_0$.  For each such fixed $M$,
a failure of \eqref{lin:torus-stable-contraction} along arbitrarily
large $L$ would supply unit stable vectors contradicting
\eqref{lin:escaped-norm}.  This proves the required order of choices:
first $M$, then $L$.  The stated bound on the map from $E_L$ follows
from the uniform step and projection bounds.  Restarting a step uses
the same formulas with its current scale as the initial scale.

Finally identify $v\in Y_L$ with $v(Lx)$ on $\mathbb T^{12}$.
On a compact range of positive $L$, these norms are equivalent to
$H^k$, continuously in $L$.  The free step on this fixed torus has
Fourier phases
$\exp(-i|n|^2L^{-2}(1-e^{-s}))$ and scalar factor $e^{(-a+ib)s}$.
Dominated convergence of the weighted Fourier sum proves strong
continuity for continuously varying inputs.  The cutoff coefficients
are smooth and continuous in $L$ on each finite slab; successive
approximation in the integral equation transfers this continuity to
the full propagator.  Dilation is strongly continuous on $Y$, which
gives the same assertion for $J_L,E_L$ and then for $\pi_L,P_L$ by
\eqref{lin:piL}.  This completes the proof.
\end{proof}

\section{Nonlinear stability, open sets, and Gaussian data}
\label{sec:nonlinear}

We first construct a decaying graph for fixed similarity parameters.
By varying phase, center, and blowup time, we then show that every datum
in a small physical $H^k$ neighborhood admits parameters placing its
transformed perturbation on that graph.
The graph construction is a discrete Lyapunov--Perron argument,
related to \cite[Proposition~4.8]{BuckEtAl2026}. Selecting symmetry
parameters by a finite-dimensional topological argument has a focusing
antecedent in \cite[Section~6.3]{Li2023}; here the graph also depends
on the expanding-torus scale.

Fix one sufficiently small admissible $a=1/(p-1)$, with $p$ odd, and one
profile $Q$ given by Proposition~\ref{profile:existence}. Fix the integer
$k>8$ and the spaces, cutoff, and spectral coordinates from
Section~\ref{sec:linear}. Every constant in this section may depend on this
fixed profile and on $k$. In particular, no estimate uniform in $p$ is needed.
We use the Euclidean norm on the fourteen spectral coordinates.

\subsection{Uniform nonlinear step maps}

Let $L_s=L e^{s/2}$ and, on the central cell of the expanding torus, set
\[
 Q_s(y)=\chi(y/L_s)Q(y).
\]
Extend this function periodically. The operator
$\partial_s+y\cdot\nabla/2$ differentiates at fixed $y/L_s$, so the
transport derivative of the cutoff is zero. Write
$\mathcal N(z)=|z|^{p-1}z$. The defect of $Q_s$ in the similarity equation is
\begin{equation}\label{nl:residual-formula}
 \mathcal R_s
 =2\nabla\chi(y/L_s)\cdot\nabla Q
   +\Delta\chi(y/L_s)Q
   +\bigl(\chi(y/L_s)-\chi(y/L_s)^p\bigr)\mathcal N(Q),
\end{equation}
where the derivatives act on the displayed functions of $y$.
Every term is supported in an annulus of radius comparable to $L_s$.
The profile symbol estimates imply, for every fixed integer $j\geq0$,
\[
 |\nabla^j\mathcal R_s|\leq C_j L_s^{-2-2a-j},
 \qquad
 \|\mathcal R_s\|_{Y_{L_s}}\leq C L_s^{-2-a}.
\]
Indeed the low Sobolev order $6-a$ contributes the factor
$L_s^{6-(6-a)}=L_s^a$, and the high order contributes a smaller power.
The regularized zero-frequency part has the same bound. On another cell
write $z=y-2\pi L_s n$. Then
$(\partial_s+y\cdot\nabla/2)z=z/2$, so this computation is compatible
with the periodic lift.

Because $p$ is odd, $\mathcal N$ is a polynomial in $z$ and $\overline z$.
The uniform algebra bounds from Section~\ref{sec:linear} and the uniform
bound for $Q_s$ therefore give, with
\[
 \mathcal E_s(v)=\mathcal N(Q_s+v)-\mathcal N(Q_s)
                    -D\mathcal N(Q_s)v,
\]
the estimates
\begin{equation}\label{nl:quadratic}
 \|\mathcal E_s(v)\|\leq C\|v\|^2,
 \qquad
 \|\mathcal E_s(v)-\mathcal E_s(\widetilde v)\|
 \leq C(\|v\|+\|\widetilde v\|)\|v-\widetilde v\|
\end{equation}
on a fixed small ball. All norms in this display are at the current scale.

Choose the step length $M$ in Proposition~\ref{linear:torus} sufficiently
large that also
\begin{equation}\label{nl:step-condition}
 2e^{-(2+a)M/2}<1.
\end{equation}
For a step beginning at scale $L$, put $L'=Le^{M/2}$ and denote its linear
evolution by $\mathcal F_L$.

\begin{lemma}[Nonlinear step estimate]\label{nl:step}
There are $\delta>0$, $L_0<\infty$ and $C_M<\infty$ such that, for
$L\geq L_0$, the perturbation equation $U=Q_s+v$ has a unique solution
through the step whenever $\|v(0)\|_{Y_L}\leq\delta$. Its endpoint is
\[
 v(M)=\mathcal F_Lv(0)+h_L(v(0)),
\]
where
\begin{align}
 \|h_L(0)\|&\leq C_M L^{-2-a},\label{nl:step-zero}\\
 \operatorname{Lip}_{\|v\|\leq d}h_L
   &\leq C_M(d+L^{-2-a})\qquad(0<d\leq\delta).
 \label{nl:step-lip}
\end{align}
During the step,
\[
 \sup_{0\leq s\leq M}\|v(s)\|_{Y_{L_s}}
 \leq C_M\bigl(\|v(0)\|_{Y_L}+L^{-2-a}\bigr).
\]
\end{lemma}

\begin{proof}
Identify the expanding spaces by dilation and use the free Fourier
propagator from Section~\ref{sec:linear}. On a fixed slab its norm is
bounded uniformly in the starting scale, as are the linearized
multiplication operators. On sufficiently short substeps, the integral
equation is a contraction on a small slab-norm ball by
\eqref{nl:quadratic}; the forcing is bounded by
\eqref{nl:residual-formula}. The short-substep length can be chosen
independently of $L$. Iteration and Gronwall's inequality give the displayed
slab bound for sufficiently small $\delta$ and large $L_0$.

At zero initial perturbation the same integral equation gives
$O_M(L^{-2-a})$, proving \eqref{nl:step-zero}. Differences of two
solutions are at most $C_M$ times their initial difference. Subtract the
linearized solution from this difference. Its remaining source is
$\mathcal E_s(v)-\mathcal E_s(\widetilde v)$; use
\eqref{nl:quadratic} and the slab bound. Integration over the fixed slab
gives \eqref{nl:step-lip}. These arguments also prove uniqueness.
\end{proof}

\subsection{A graph of decaying perturbations}

Set $L_n=Le^{nM/2}$ and abbreviate $E_{L_n},\pi_{L_n},P_{L_n}$ by
$E_n,\pi_n,P_n$. The normalized coordinates satisfy $\pi_nE_n=I$.
Let $D_*=e^{M\Lambda}$, where $\Lambda$ has entries $0$, twelve copies
of $1/2$, and $1$. Thus $\|D_*^{-1}\|\leq1$.
Split
\[
 v_n=w_n+E_nu_n,\qquad w_n\in\ker\pi_n,\quad u_n\in\mathbb R^{14}.
\]
Proposition~\ref{linear:torus} and Lemma~\ref{nl:step} give the recursions
\begin{align}
 w_{n+1}&=A_nw_n+B_nu_n+(I-P_{n+1})h_{L_n}(v_n),\label{nl:forward}\\
 u_{n+1}&=D_*u_n+e_n(v_n)+\pi_{n+1}h_{L_n}(v_n),\label{nl:coordinates}
\end{align}
where
\[
 \|A_n\|\leq\tfrac18,\qquad \|B_n\|\leq C_M,\qquad
 \sup_{n\geq0}\|e_n\|\longrightarrow0\quad\text{as }L\to\infty.
\]
Here $e_n=\pi_{n+1}\mathcal F_{L_n}-D_*\pi_n$. All coordinate maps
and projections are uniformly bounded.

\begin{proposition}[Continuous stable graph]\label{nl:graph}
For all sufficiently large $L$ there is a continuous map $G_L$, defined
on a fixed small ball in $\ker\pi_L$, such that the initial perturbation
\[
 v_0=w_0+E_LG_L(w_0)
\]
generates a solution on every similarity-time slab and
\[
 \|v(s)\|_{Y_{L_s}}\longrightarrow0\quad(s\to\infty).
\]
Uniformly in $L$,
\begin{equation}\label{nl:graph-size}
 |G_L(w_0)|\leq C\bigl(\|w_0\|+L^{-2-a}\bigr).
\end{equation}
The coordinate $G_L(w_0)$ is jointly continuous in $L$ and $w_0$ after
identifying periodic functions by dilation.
\end{proposition}

\begin{proof}
For a constant $K\geq1$ to be fixed, use the sequence norm
\[
 \|(w,u)\|_{\mathcal X}
 =\sup_{n\geq0}2^n\bigl(\|w_n\|+K|u_n|\bigr).
\]
Define a map by the right side of \eqref{nl:forward} for $w_{n+1}$,
resetting its zeroth coordinate to the prescribed $w_0$, and by
\begin{equation}\label{nl:backward}
 u_n=-\sum_{j=n}^{\infty}D_*^{n-j-1}
       \bigl(e_j(v_j)+\pi_{j+1}h_{L_j}(v_j)\bigr)
\end{equation}
for the coordinate sequence. The geometric weights make this sum
convergent even in the phase coordinate, where $D_*$ equals one.

For differences of sequences, the linear terms in the forward map cost
at most $2(1/8+C_M/K)$ in the sequence norm. Choose $K$ so that this is
less than $1/2$. The weighted sum in \eqref{nl:backward} costs at most
$2K$ times the weighted supremum of the forcing difference. By
\eqref{nl:step-lip}, the remaining Lipschitz costs tend to zero as the
sequence ball is made small and then $L$ is made large. The same is true
of $e_j$, uniformly in $j$. Consequently the complete map has a
contraction constant strictly below one, uniformly for all sufficiently
large $L$.

At the zero sequence the residual forcing is $O_M(L_n^{-2-a})$.
Condition~\eqref{nl:step-condition} bounds its weighted sequence norm by
$C_M L^{-2-a}$. Thus a fixed small ball is invariant if $w_0$ is small,
and the contraction theorem yields a unique sequence in that ball.
Its norm is at most $C(\|w_0\|+L^{-2-a})$. The fixed point satisfies
\eqref{nl:coordinates}, since consecutive backward sums telescope.
Define $G_L(w_0)=u_0$. This proves \eqref{nl:graph-size} and the decay
at step endpoints. Lemma~\ref{nl:step} gives the decay between endpoints.

We spell out the continuity assertion because the free Schr\"odinger
flow is not continuous in operator norm as its time varies. At every
fixed index, and after dilation to a fixed torus, the free Fourier
formula is strongly continuous in the starting scale and in the input.
The polynomial integral equation preserves this continuity on each
finite slab. Cutoffs and coordinate maps are likewise strongly
continuous: prove this first for smooth functions, then use density and
their locally uniform bounds. The finite matrix normalizing $\pi_L$
stays invertible and varies continuously.

Start the sequence contraction at zero. Each iterate is continuous at
every fixed coordinate: in \eqref{nl:backward}, the summands are
continuous and the geometric tail is uniform. The iterates converge
uniformly in the sequence norm on the allowed parameter set. Their
zeroth coordinate is therefore continuous. This proves the assertion
without differentiating the graph or taking an operator-norm limit of
Schr\"odinger propagators.
\end{proof}

\subsection{Physical parameters fill an open neighborhood}

For a small $T_0>0$, define the smooth periodic reference datum by
\begin{equation}\label{nl:reference}
 u_{\mathrm{ref}}(x)=T_0^{-a}\chi(x)Q(x/\sqrt{T_0})
\end{equation}
in the central cell. Its support is contained strictly inside that cell.
For physical data $u_{\mathrm{in}}$ near this datum, vary
\[
 T=T_0e^\beta,\qquad x_*=\sqrt{T_0}\,\zeta,\qquad \theta\in\mathbb R,
 \qquad d'=(\theta,\zeta,\beta)\in\mathbb R^{14}.
\]
Put $L=T^{-1/2}$. The required starting perturbation is
\begin{equation}\label{nl:starting-perturbation}
 v_0(d';u_{\mathrm{in}})
 =e^{-i\theta}T^a u_{\mathrm{in}}(x_*+\sqrt T\,y)
      -\chi(y/L)Q(y).
\end{equation}

\begin{lemma}[Parameter expansion]\label{nl:parameter-expansion}
For $|d'|\leq\rho$ with $\rho$ fixed and small, uniformly as $T_0\to0$,
\begin{equation}\label{nl:leading-expansion}
 v_0(d';u_{\mathrm{ref}})
 =E_L C d'+r(d'),\qquad
 C d'=(-\theta+b\beta,\zeta,\beta),
\end{equation}
and
\[
 \|r(d')\|_{Y_L}\leq C\bigl(|d'|^2+\sqrt{T_0}|d'|\bigr).
\]
For fixed $T_0$, replacing $u_{\mathrm{ref}}$ by $u_{\mathrm{in}}$
adds at most $C(T_0)\|u_{\mathrm{in}}-u_{\mathrm{ref}}\|_{H^k}$,
uniformly on the parameter ball. All these maps are continuous.
\end{lemma}

\begin{proof}
Substitution of \eqref{nl:reference} into
\eqref{nl:starting-perturbation} gives
\[
 e^{-i\theta}e^{a\beta}
 \chi(\sqrt{T_0}\zeta+\sqrt T\,y)
 Q(\zeta+e^{\beta/2}y)-\chi(\sqrt T\,y)Q(y).
\]
Taylor expansion of the smooth reference profile yields
\[
 \chi(y/L)\bigl(-i\theta Q+\zeta\cdot\nabla Q
                    +\beta(a+y\cdot\nabla/2)Q\bigr).
\]
The profile symbol bounds and their dyadic norm estimates apply to its
first two parameter derivatives: translations preserve the tail bounds,
and the dilation derivative $y\cdot\nabla$ preserves their order.
Hence the Taylor remainder is $O_{Y_L}(|d'|^2)$ uniformly in $L$.
The translated cutoff differs by $O(\sqrt{T_0}|\zeta|)$, with the same
scaled derivative estimates on its annulus; multiplication by the
profile gives the additional stated error. For small parameters the
supports remain strictly inside one physical cell.

The spectral time vector is $(a-ib+y\cdot\nabla/2)Q$.
The missing $ibQ$ contributes $b\beta$ to the phase coordinate, proving
\eqref{nl:leading-expansion}. On the compact parameter ball with $T_0$
fixed, translation and rescaling from physical $H^k$ to $Y_L$ have
uniformly bounded norms. They are strongly continuous, which proves the
remaining assertions. Only the smooth reference has been differentiated;
no extra derivative is required of nearby $H^k$ data.
\end{proof}

\begin{proof}[Proof of Theorem~\ref{thm:blowup}]
Choose $a$ and $k$ as above. We seek a zero of the continuous map
\begin{equation}\label{nl:degree-map}
 \Phi(d')=\pi_Lv_0(d';u_{\mathrm{in}})
       -G_L\bigl((I-P_L)v_0(d';u_{\mathrm{in}})\bigr).
\end{equation}
The matrix $C$ in Lemma~\ref{nl:parameter-expansion} is invertible.
At every current scale, $\pi_LE_L=I$ and $(I-P_L)E_L=0$ exactly.
Consequently the stable argument of $G_L$ contains only the error in
that lemma; varying $L$ introduces no leading stable component.
By \eqref{nl:graph-size}, on $|d'|\leq\rho$ we have
\begin{equation}\label{nl:degree-error}
 |\Phi(d')-Cd'|
 \leq C_1\bigl(\rho^2+\sqrt{T_0}\rho
                   +\varepsilon_{\mathrm{data}}+T_0^{1+a/2}\bigr),
\end{equation}
where
$\varepsilon_{\mathrm{data}}=C(T_0)
\|u_{\mathrm{in}}-u_{\mathrm{ref}}\|_{H^k}$.

Let $c_C>0$ be the smallest singular value of $C$.
First choose $\rho$ small enough for the quadratic term in
\eqref{nl:degree-error} to be less than $c_C\rho/4$ and for the graph
domain requirements to hold. Next choose $T_0$ sufficiently small that
all scales are in the range of Proposition~\ref{nl:graph}, and the two
$T_0$ terms together are less than $c_C\rho/4$.
Finally choose a positive physical $H^k$ radius so that the data term
is less than $c_C\rho/4$. The graph is then defined throughout the
closed parameter ball, and on its boundary
\[
 |\Phi(d')-Cd'|<c_C\rho\leq |Cd'|.
\]
The straight-line homotopy to $C$ has no boundary zero. Brouwer degree
therefore gives a zero of \eqref{nl:degree-map} for every datum in this
open $H^k$ ball.

At such a zero the initial perturbation lies on the graph. The resulting
similarity solution exists for all $s\geq0$ and its perturbation tends
uniformly to zero, by Proposition~\ref{nl:graph} and the uniform
$Y_{L_s}\to L^\infty$ embedding. Undoing the similarity variables gives
the physical solution on $0\leq t<T$:
\[
 u(t,x)=e^{i\theta}(T-t)^{-a}
       \left(\frac{T-t}{T}\right)^{ib}
       \bigl(Q_s(y)+v(s,y)\bigr),
 \quad y=\frac{x-x_*}{\sqrt{T-t}},\quad
 s=\log\frac{T}{T-t}.
\]
On each finite similarity slab, the norms are equivalent after dilation
to the physical $H^k$ norm. The integral equation and polynomial algebra
therefore give $u\in C([0,T);H^k)\cap C^1([0,T);H^{k-2})$.
The usual local uniqueness follows directly from the physical
Schr\"odinger integral equation and its $H^k$ difference estimate.
At the center, $Q_s(0)=Q(0)\ne0$, and hence
\[
 (T-t)^a|u(t,x_*)|\longrightarrow |Q(0)|>0.
\]
Since $H^k(\mathbb T^{12})$ embeds continuously into $C^0$, the solution
cannot extend continuously in $H^k$ through $T$. This proves the theorem.
\end{proof}

\paragraph{Higher-dimensional deterministic examples.}
For the fixed odd power $p$ in Theorem~\ref{thm:blowup}, smooth
finite-time blowup data also exist on every standard product torus
$\mathbb T^d$ with $d\ge12$. To see this, choose the smooth reference
datum \eqref{nl:reference} at the center of the blowup ball just
constructed. Its solution $u$ remains smooth before its blowup time $T$.
Indeed, for each integer $m\ge k$, the Fourier convolution estimate gives
\[
 \big\||v|^{p-1}v\big\|_{H^m(\mathbb T^{12})}
 \le C_{m,p,k}\|v\|_{H^k(\mathbb T^{12})}^{p-1}
                 \|v\|_{H^m(\mathbb T^{12})}.
\]
The frequency weight of order $m$ is bounded by a sum of weights on
single factors, and the remaining factors are controlled by
$\|\widehat v\|_{\ell^1}\lesssim\|v\|_{H^k}$ since $k>6$.
On every $[0,\tau]\subset[0,T)$, the $H^k$ norm of $u$ is bounded.
Duhamel's formula and Gronwall therefore bound each higher Sobolev norm;
local $H^m$ existence, continuation, and $H^k$ uniqueness identify these
solutions with $u$ throughout $[0,\tau]$. The equation then supplies all
time derivatives. This persistence argument takes place entirely in
dimension twelve.

Now set $U(t,x,y)=u(t,x)$ on
$\mathbb T^{12}\times\mathbb T^{d-12}$. The extra Laplacian derivatives
vanish, so $U$ is a smooth solution of \eqref{eq:nls}.
The polynomial nonlinearity is Lipschitz on bounded sets, so the $L^2$
difference estimate and Gronwall give uniqueness among smooth solutions
on compact time intervals.
For every $y$, the center limit from Theorem~\ref{thm:blowup} becomes
$(T-t)^{1/(p-1)}|U(t,x_*,y)|\to c_0>0$.
Thus the same time and pointwise blowup rate persist along the product
fiber. No higher-dimensional stability or Gaussian-probability conclusion
is asserted.

\subsection{The Gaussian counterexample}

\begin{proof}[Proof of Corollary~\ref{cor:gaussian}]
Fix the pair $p,k$ supplied by Theorem~\ref{thm:blowup} and take
$\alpha>k+6$. Independence and the finite Gaussian second moment give
\[
 \mathbb E\|u_0^\omega\|_{H^k}^2
 =c_g\sum_{n\in\mathbb Z^{12}}\langle n\rangle^{2k-2\alpha}<\infty,
\]
up to the fixed Fourier normalization. Thus the series converges as an
$H^k$-valued Gaussian random variable. Its finite Fourier segment has a
strictly positive density on the whole finite-dimensional complex
coefficient space.

For completeness, let $B_{H^k}(f,\varepsilon)$ be any nonempty ball.
Choose a finite Fourier projection $P_N$ such that
$\|(I-P_N)f\|_{H^k}<\varepsilon/4$ and
\[
 \mathbb E\|(I-P_N)u_0^\omega\|_{H^k}^2<\varepsilon^2/32.
\]
Markov's inequality gives positive probability that the random tail
has norm less than $\varepsilon/4$ after increasing $N$ further if
necessary. Independently, the finite segment belongs to the ball of
radius $\varepsilon/2$ about $P_Nf$ with positive probability.
The intersection lies in $B_{H^k}(f,\varepsilon)$, so the law has full
support in $H^k$.

Apply this to the open blowup neighborhood in Theorem~\ref{thm:blowup}.
It has positive probability for every $\alpha>k+6$. Since
$8<k<\alpha-6$, the exponent $m=k$ lies in the interval in the universal assertion of
Corollary~\ref{cor:gaussian}. The local solution fails its required global continuation
with positive probability. The fixed pair also satisfies
$s_{\mathrm{cr}}=6-2/(p-1)>1$. Hence no finite threshold for this pair,
and therefore no universal assertion of the stated form, is possible.
\end{proof}


\appendix
\section{Exact finite arithmetic}\label{cert:arithmetic}

This appendix specifies the finite calculations used in
Proposition~\ref{free:main}.  They involve integers and rational numbers
only.  The accompanying file
\texttt{verification/}\allowbreak\texttt{certificates/}\allowbreak\texttt{exact\_verifier.py}
implements the recipes below using the Python standard library; its
deterministic output is
\texttt{verification/}\allowbreak\texttt{certificates/}\allowbreak\texttt{RESULTS.json}.  In particular,
neither special-function evaluation nor a floating-point root finder is
used to decide any inequality.  The mathematical identities that permit
coefficient truncation are proved in Lemma~\ref{free:boundary} and below.

Reproduction commands and the exact input and artifact bindings are in
\texttt{README.md} and
\texttt{verification/}\allowbreak\texttt{certificates/}\allowbreak\texttt{README.md}, supplied with the
manuscript. The verifier checks the finite comparisons below. Their
connection to the profile degree and spectral count is established by
the analytic arguments of Section~\ref{sec:free}.

\subsection{Coefficient enclosures}

For a complex polynomial $X$, write $|X|_1$ for the polynomial whose
$j$th coefficient is $|\operatorname{Re}X_j|+|\operatorname{Im}X_j|$.
An enclosure $[X;e]$, with nonnegative real coefficients $e_j$, means
that the actual $j$th coefficient differs from $X_j$ by a complex number
of $\ell^1$ norm at most $e_j$.  Addition adds the errors; negation,
conjugation, and reflection of the variable preserve their absolute
bounds.  Polynomial multiplication uses convolution and the rule
\begin{equation}\label{cert:product-rule}
 [X;e][Y;f]=[XY;\ e(|Y|_1+f)+|X|_1f].
\end{equation}
This follows by expanding $(X+\varepsilon)(Y+\zeta)-XY$ and using
$|zw|_1\le|z|_1|w|_1$ coefficientwise before convolution.
The program stores a coefficient as the triple
$(\operatorname{Re}X_j,\operatorname{Im}X_j,e_j)$.

For the boundary and exterior tests set
\begin{equation}\label{cert:integer-box}
 H_0=10^8,\qquad B=33477607,\qquad S=270506819,
 \qquad |H_0b-B|<2,\quad |H_0Z-S|<2.
\end{equation}
The disk \eqref{free:disk} lies strictly inside this box: substitution of
its rational center shows
$|H_0b_*-B|+H_0\delta<2$ and
$|H_0Z_*-S|+H_0\delta<2$.  Repeated occurrences of an uncertain
coefficient may be bounded independently in
\eqref{cert:product-rule}; this enlarges the enclosure.  The polynomial
identities used to discard coefficients concern the coherent family
with the same actual $b,Z$ in every occurrence, and do not rely on
independent error variables satisfying those identities.

\subsection{The left boundary forms}\label{cert:boundary}

Take $K=8$ and $0\le\ell\le3$.  Every quantity in the following
recursion is a polynomial in the real variable $v$.  Use the central
values
\[
 t_n=H_0(\ell/2+n-1/32)+i(H_0v-B),\quad
 M=H_0(\ell+5),\quad s=iS,\quad d_n=t_n-M.
\]
Give the constant coefficient of $t_n,d_n,s$ error $2$, and every
other input coefficient error zero.  Starting from row vectors
$x=(1,0)$, $y=(0,1)$, do, for $n=0,\ldots,7$,
\begin{equation}\label{cert:forward-recipe}
 x_{\rm new}=t_nx+sy,\qquad
 y_{\rm new}=d_ny-x_{\rm new},\qquad (x,y)\leftarrow(x_{\rm new},y_{\rm new}).
\end{equation}
Define the sesquilinear expression
\[
 \begin{gathered}
 h(a,b',c,d)=M\overline a c+s(\overline a d-\overline{b'}c),
 \\
 A=h(x_0,y_0,x_0,y_0),\quad
 D=h(x_1,y_1,x_1,y_1),\quad
 C=h(x_0,y_0,x_1,y_1).
 \end{gathered}
\]
For any polynomial let $f^\#(v)=\overline{f(-v)}$, entrywise on
coefficients.  Replace each of $A,D,C$ by its sum with its sharp,
and form $AD-C\overline C$.  Relative to
\eqref{free:boundary-form}, this simply scales each forward matrix by
$H_0$ and the cone matrix by $H_0$, so the determinant is multiplied
by the positive constant $H_0^{4K+2}$.

Lemma~\ref{free:boundary} proves that the actual determinant is real,
even, and of degree at most 30 throughout the parameter box.  Retain
its even real coefficient enclosures and write it as
$D_\ell(w)=\sum_{j=0}^{15}D_jw^j$, where $w=v^2$.
There is no inference from small computed high coefficients here:
their vanishing follows from the structural proof.  The program also
checks the central vanishings as a consistency check.

For each of
\begin{equation}\label{cert:substitutions}
 (U,V,J)=(0,1,1),\quad(1,10,1),\quad(10,1,0),
\end{equation}
compute the coefficient enclosures of
\begin{equation}\label{cert:transformed-polynomial}
 (1+Jw)^{15}D_\ell\left(\frac{U+Vw}{1+Jw}\right)
 =\sum_{j=0}^{15}D_j(U+Vw)^j(1+Jw)^{15-j}.
\end{equation}
For example, a multiplication-only implementation starts $Q=0$, $E=1$
and, for $j=15,\ldots,0$, replaces
$Q\leftarrow Q(U+Vw)+ED_j$, $E\leftarrow E(1+Jw)$.
If $[X;e]$ is the resulting enclosure, the minima over coefficients
$0\le j\le15$ of $\lfloor\operatorname{Re}X_j/(1+e_j)\rfloor$ are
\begin{equation}\label{cert:boundary-table}
 \begin{array}{c|rrr}
 \ell&(0,1,1)&(1,10,1)&(10,1,0)\\\hline
 0&63&60&94\\
 1&21&178&199\\
 2&4&13&323\\
 3&1115&173&298
 \end{array}.
\end{equation}
Every entry is positive, hence every actual coefficient in
\eqref{cert:transformed-polynomial} is strictly positive.  The three
maps $(U+Vw)/(1+Jw)$ for $w\ge0$ cover $[0,1)$, $[1,10)$, and
$[10,\infty)$, respectively.  The clearing factors are positive.
Thus $D_\ell(v^2)>0$ for every real $v$.  This is precisely the
determinant inequality used to make the boundary form positive
definite in Lemma~\ref{free:boundary}.

\subsection{Rational separators and winding}\label{cert:winding}

The count uses no interval errors once the zero-tail parameters have
been moved to $(B/H_0,S/H_0)$.  Start $x=0$, $y=1$ as polynomials in
$X$.  For $n=7,\ldots,0$ set
\[
 t=H_0(X-1/32+\ell/2+n)-iB,
 \qquad
 (x,y)\leftarrow
 \big((t-H_0(\ell+5)-iS)x-iSy,\ t(x+y)\big),
\]
using the old pair on the right.  If $p_j$ denotes
\begin{equation}\label{cert:p-coefficients}
 p_j=\operatorname{Im}\sum_{k=0}^j x_k\overline{y_{j-k}},
\end{equation}
then $P_\ell(X)=\sum p_jX^j$ is the polynomial of
Lemma~\ref{free:winding}, up to a positive common scaling.
Its degree is 15 and $p_{15}=-8SH_0^{15}$; $p_{16}=0$.
The two axis polynomials are therefore
\[
 \mathcal R_\ell(w)=\sum_{j=0}^7(-1)^jp_{2j}w^j,
 \qquad
 \mathcal I_\ell(w)=\sum_{j=0}^7(-1)^jp_{2j+1}w^j.
\]

Evaluate both at $w=(q/100)^2$, taking endpoints $q=0,7500$
and the following intermediate lists:
\[
\begin{array}{c|l}
 \ell&q\\\hline
 0&45,90,145,220,315,450,635,920,1780,3100\\
 1&30,68,116,173,250,350,485,680,1005,1880,3210\\
 2&51,78,94,121,165,218,278,369,490,720,1080,1930,3340\\
 3&57,95,122,157,205,265,333,431,559,784,1148,2020,3521
\end{array}
\]
In the order including endpoints, the signs follow the repeating
cycle $++,-+,--,+-$, starting at $-+$ in row zero, $++$ in row one,
and $--$ in rows two and three.  At the two negative arguments
$w=-(q/100)^2$, $q=50,200$, rows zero and one have signs $++,--$,
in that order.

Here is an integer-only recipe for every sign, including a relative
margin against cancellation in the sum.  For coordinate $c=0,1$ and
argument $w=\varepsilon(q/100)^2$, $\varepsilon\in\{1,-1\}$, set
\[
 d_j=p_{2j+c}100^{14-2j}(-\varepsilon q^2)^j,
 \qquad j=0,\ldots,7.
\]
The sign of $\sum d_j$ is the required sign.  For its specified sign
$e\in\{1,-1\}$, the minimum, over all tests in that row, of
\[
 \left\lfloor\frac{10^4e\sum d_j}{\sum|d_j|}\right\rfloor
\]
is respectively
\begin{equation}\label{cert:separator-margins}
 583,\qquad432,\qquad80,\qquad68.
\end{equation}
These strictly positive integer margins verify all stated signs.
The sign changes give $5,6,7,7$ positive root brackets for
$\mathcal R_\ell$ and $6,6,7,7$ for $\mathcal I_\ell$.
After the stated negative brackets are included, both have seven
distinct brackets.  The degree bounds then prove simplicity and
exhaustion.  Lemma~\ref{free:winding} gives the argument change and
its orientation, converting this finite sign table into the matching
counts $2,1,0,0$.

\subsection{The high angular multiplier}\label{cert:high-angular}

The path constants in Lemma~\ref{free:high-angular} are bounded using
only the rational endpoint values $\cos\vartheta_0=99/101$ and
$\sin\vartheta_0=20/101$, and monotonicity of sine and cosine on the
initial arc.  For example,
\[
 \sin(2\vartheta_0)<.4,\quad \cos(2\vartheta_0)>.92,
 \quad \cos\vartheta_0>.98,
 \quad N\le.04-2.70\cdot.92=-2.444.
\]
Thus $L'>.98-4(-2.444)>10.7$.  The lower bound on the arc potential
after dropping its nonnegative terms is
\[
 -3-1/32+.3(10.7)=143/800>0.
\]
At the join, substitute
$u=10/101$, $v=Z+1/101$ in the formulas for $L,N$.
The disk bounds give $1.13<L<1.19$, so the jump lower bound is
\[
 -1.5-.3(1.19)+2.444=587/1000>0.
\]
On the ray, $u$ and $L$ have slope $99/101\in(.98,.981)$
and $v$ has slope $20/101<.20$; these prove every affine bound used
in \eqref{free:high-polynomial}.  Also $M/2-Z\ge9/2-Z_*-\delta>0$.

For transparency, let $P(t)$ be exactly the rational polynomial in
\eqref{free:high-polynomial} and put
\[
 Q(t)=2+t(33-72.2t+44.8t^2)+t^4(17.9-4.61t+.54t^2).
\]
The seven coefficients of $P-Q$, in ascending order, are
\begin{equation}\label{cert:high-difference}
 \begin{split}
 &\frac{142518401}{500000000},\quad
 \frac{325932989}{1000000000},\quad
 \frac{15777215591}{200000000000},\quad
 \frac{8728703679}{200000000000},\\
 &\frac{34209914771}{800000000000},\quad
 \frac{16113799}{3125000000},\quad
 \frac{60018849}{40000000000}.
 \end{split}
\end{equation}
They are strictly positive.  The two quadratic terms in $Q$ have
positive leading coefficient and discriminants
\[
 (72.2)^2-4(33)(44.8)=-17519/25,\qquad
 (4.61)^2-4(17.9)(.54)=-174119/10000.
\]
Thus $P>Q>0$ for $t\ge0$.  The denominator $L_0d_0(1+1.5t)^2$
is positive there.  These comparisons supply exactly the potential
positivity needed in the multiplier identity.

\subsection{The profile product and degree margin}\label{cert:profile}

This calculation uses exact complex rational scalars and ordinary
$2\times2$ matrix multiplication.  For an exact complex matrix $F$,
$|F|_1$ means entrywise complex $\ell^1$ norm.  Initialize
$U=I$, $X=Y=E=0$.  For $n=33,\ldots,0$ set
\[
 \begin{gathered}
 t=n-ib_*,\quad
 F=\begin{pmatrix}t-5-iZ_*&-iZ_*\\t&t\end{pmatrix},\\
 G=-i\begin{pmatrix}1&0\\1&1\end{pmatrix},\quad
 H_1=-i\begin{pmatrix}1&1\\0&0\end{pmatrix},\quad N=|G|_1+|H_1|_1,
 \end{gathered}
\]
and replace the four matrices simultaneously by
\begin{equation}\label{cert:profile-recipe}
 \frac1{n+1}\left(
 FU,\quad GU+FX,\quad H_1U+FY,\quad
 (|F|_1+\delta N)E+N(|X|_1+|Y|_1)\right).
\end{equation}
Then $U=U_0$, $X=X_0$, $Y=Y_0$ in
Lemma~\ref{free:profile-degree}, while the nonnegative real matrix
$\delta^2E$ bounds the entrywise complex $\ell^1$ error of the
linear approximation.  Indeed multiplying a current error by the
perturbed factor costs $(|F|_1+\delta N)E$; the product of the two
linear perturbations costs $N(|X|_1+|Y|_1)$ after division by
$\delta^2$.  This proves the majorant recurrence, including all
higher order terms.

With $m_0=U_{01}$, the comparisons in
\eqref{free:profile-enclosures} are checked by ordinary rational
division, or by squaring positive modulus bounds.  For example
$9<|m_0|<10$ means $81<|m_0|^2<100$.  The determinant of a backward
factor is $t(t-5)$, independent of $Z$, so a rational upper bound for
$|\det L|^2$ throughout the disk is
\begin{equation}\label{cert:profile-determinant}
 \frac{\displaystyle\prod_{n=0}^{33}
 (n^2+(b_*+\delta)^2)((n-5)^2+(b_*+\delta)^2)}{(34!)^4}.
\end{equation}
Comparison with $(6.2\cdot10^{-12})^2|m_0|^4$ proves the last
inequality of \eqref{free:profile-enclosures}.

Here are all the additional comparisons and the formulas propagating
them to the analytic quotient.  Put
\[
 x=X_{11}/m_0,\qquad y=Y_{11}/m_0,\qquad u=U_{00}/m_0,
 \qquad \kappa=(.048)^2.
\]
The exact rational calculations give
\begin{gather}
 |x|^2>\kappa,\qquad
 (|x|^2-\kappa)(|y|^2-\kappa)
       >(\operatorname{Re}(\overline xy))^2,\notag\\
 |1-iZ_*u/5|>2.34,\qquad |u|<3.18,\qquad
 |x|<.159,\qquad |y|<.124.\label{cert:profile-extra}
\end{gather}
The first line is Sylvester's criterion for the Gram matrix of the
real-linear map $(\eta,\zeta)\mapsto\eta x+\zeta y$, minus
$\kappa I$, to be positive definite.  It proves minimum stretch
greater than $.048$.  The rectangles in
\eqref{free:profile-enclosures} also show
$\operatorname{Re}x\operatorname{Im}y-\operatorname{Im}x\operatorname{Re}y<0$.

To make the error propagation fully reproducible, define the positive
rationals
\begin{align*}
 \chi&=40\delta+50000\delta^2,\\
 v_1&=(1+(Z_*+\delta)/5)\chi+3.18\delta/5,\\
 v_2&=(Z_*+\delta)\chi/5+3.18\delta/5,\\
 \gamma&=2.34-3.18Z_*/5-v_1-v_2.
\end{align*}
The entrywise change of $L/m_0$ from $U/m_0$ is less than $\chi$.
The quantities $v_1$ and $v_2$ bound the change in the two terms of
the denominator gap, respectively.  Exact comparison gives
\[
 \gamma>.6,\qquad \chi<5\cdot10^{-7},
\]
which proves \eqref{free:denominator-gap}.  In particular all divisions
used in the quotient estimate have already been justified.
Set
\begin{align}\label{cert:ratio-errors}
 e_{\rm jet}&=
 \frac{10^{-12}+50000\delta^2+(.283\delta)\chi}{1-\chi},\notag\\
 e_{\rm tail}&=
 \frac{(2(Z_*+\delta)/5)(6.2\cdot10^{-12})}{.6(1-\chi)}.
\end{align}
The first bounds the distance of $L_{11}/L_{01}$ from the linear map:
its numerator accounts for $U_{11}$, the quadratic remainder, and
the change of denominator.  The second is the determinant estimate
for the tail-ratio correction in Lemma~\ref{free:profile-degree}.
The remaining comparisons are
\begin{equation}\label{cert:degree-comparisons}
 e_{\rm jet}+e_{\rm tail}<2.5\cdot10^{-11},\qquad
 .283\delta+2.5\cdot10^{-11}<3\cdot10^{-9},\qquad
 .048\delta>2.5\cdot10^{-11}.
\end{equation}
Together with \eqref{cert:integer-box}, these are the complete finite
comparisons behind the disk estimate and nonzero degree.  The actual
boundary margin is greater than
$4.8\cdot10^{-10}-2.5\cdot10^{-11}=4.55\cdot10^{-10}$.
The verifier emits every comparison separately and also emits the
rational values of $e_{\rm jet}+e_{\rm tail}$ and $\gamma$.

\subsection{Exterior positivity and its degree cancellations}\label{cert:exterior}

Here $M=5$, $t_n=n-ib$, $d_n=t_n-5$, $s=iz$ and $K=5$ in
the unscaled forward recurrence.  We first justify the degree bounds
used in the arithmetic.  Write $T^{(n)}$ for the product of its first
$n$ factors.  For $n\ge1$ its column zero has degree at most $n-1$
in $z$, and column one has degree at most $n$.  Their leading terms are
\begin{equation}\label{cert:exterior-leading}
 T^{(n)}_{\cdot0}=-b a_nz^{n-1}\binom{1}{-1}
     +O(z^{n-2}),\qquad
 T^{(n)}_{\cdot1}=a_nz^n\binom{1}{-1}+O(z^{n-1}),
 \quad a_n=i(-i)^{n-1}.
\end{equation}
For $n=1$ the first column is exactly $-ib(1,-1)^t$, so the
remainder in that first assertion is zero.  Induction follows by
retaining the $z$ part
$\left(\begin{smallmatrix}0&i\\0&-i\end{smallmatrix}\right)$
of the next factor; its action on $(1,-1)^t$ is multiplication by
$-i$.

Let $x,y$ be the rows after $n$ factors and $y^{\rm old}$ the second
row after $n-1$.  The recurrence gives
$y=d_{n-1}y^{\rm old}-x$.  Thus
\begin{equation}\label{cert:last-row-cancellation}
 \overline{x_i}y_j-\overline{y_i}x_j
 =d_{n-1}\overline{x_i}y_j^{\rm old}
  -\overline{d_{n-1}}\overline{y_i^{\rm old}}x_j.
\end{equation}
In the Hermitian form
$M\overline{x_i}x_j+iz(\overline{x_i}y_j-\overline{y_i}x_j)$,
this proves degree at most $2n$ in its last diagonal entry $d$,
and at most $2n-1$ in its entry $c$ with indices $(1,0)$.
For the additional cancellation take $n\ge2$, as is the case for
$K=5$.  By
\eqref{cert:exterior-leading}, the leading coefficient of
$\overline{x_1}x_0$ is real, namely $-b$.  The coefficient of
$z^{2n-2}$ on the right of
\eqref{cert:last-row-cancellation} for $(i,j)=(1,0)$ is
$2ib\operatorname{Re}d_{n-1}$, which is purely imaginary.
Multiplication by $iz$ makes its highest coefficient real as well.
Consequently $\operatorname{Im}c$ has degree at most $2n-2$.
At $n=5$ the bounds are therefore
\begin{equation}\label{cert:exterior-degrees}
 \deg d\le10,\qquad \deg c\le9,\qquad
 \deg\operatorname{Im}c\le8.
\end{equation}
These are identities for every real $b$, and translating $z$ preserves
them.  In particular, discarded error bounds above these degrees do
not represent actual uncertainty in a nonzero coefficient.

For the finite tests put $z=z_1+v$, $v\ge0$, with $z_1=2.704$ or
$3$, and repeat \eqref{cert:forward-recipe} five times with
\[
 M=5H_0,\qquad t_n=H_0n-iB,\qquad d_n=t_n-M,
 \qquad s=iH_0(z_1+v).
\]
Only the constant coefficients of $t_n,d_n$ have error $2$; $s$ is
exact.  Form $d=h(x_1,y_1,x_1,y_1)$ and
$c=h(x_1,y_1,x_0,y_0)$ by the same recipe as above.  These are
$H_0^{11}$ times the unscaled entries.  The three tests, including
their truncation degrees and interval-ratio minima, are
\begin{equation}\label{cert:exterior-table}
\begin{array}{c|l|c|r}
 z_1&\text{polynomial}&\text{degree}&
 \min\lfloor\operatorname{Re}X_j/(1+e_j)\rfloor\\\hline
 2.704&\operatorname{Re}d&10&288247\\
 2.704&43\operatorname{Re}d+
 1000(\operatorname{Re}c-1130H_0^{11}(z_1+v))&10&17080\\
 3&-\operatorname{Im}c-1130H_0^{11}(z_1+v)&8&31628
\end{array}
\end{equation}
Every actual coefficient is positive, so these polynomials are
strictly positive for $v\ge0$.  Dividing the second by
$1000H_0^{11}$ gives $.043d+\operatorname{Re}c-1130z>0$ in
unscaled notation; the other two give the remaining form inequalities.
Finally $0<b<.335$ throughout the disk, and the exact comparison
\begin{equation}\label{cert:exterior-determinant}
 \prod_{n=0}^4(n^2+.335^2)((n-5)^2+.335^2)<1130^2
\end{equation}
implies $|\det T|<1130$.  This completes the finite verification of
\eqref{free:exterior-forms} and all numerical assertions used in
Section~\ref{sec:free}.


\bibliographystyle{plain}
\bibliography{references}
\end{document}
